% Family and social life — Anglais 1re Tech — Cours signé OnBuch+
% Programme officiel MINESEC. Compiler avec : tectonic N01-family-and-social-life.tex
\newcommand{\DOCMATIERE}{Anglais}
\newcommand{\DOCNIVEAU}{1re Tech}
\newcommand{\DOCMODULE}{Anglais – classe de Première technique}
\newcommand{\DOCLECON}{N01}
\newcommand{\DOCTITRE}{Family and social life}
% OnBuch+ — préambule « cours signés » ENRICHI (compilé avec tectonic / XeLaTeX)
% Système visuel « Pop » d'OnBuch+ : fond crème (jamais blanc), bordures encre
% épaisses, ombres « dures » décalées, titres Archivo Black en capitales.
% Version MATHÉMATIQUES : graphes (pgfplots/tikz), boîtes pédagogiques
% (définition, propriété, méthode, attention, le savais-tu, exercice résolu,
% exercice, TP, bilan), page de garde et signature OnBuch+ ancrée en bas.
%
% Macros attendues dans main.tex AVANT \input{preamble} :
%   \DOCMATIERE  \DOCNIVEAU  \DOCMODULE  \DOCLECON (numéro)  \DOCTITRE
\documentclass[11pt]{article}

\usepackage{fontspec}
\usepackage[french]{babel}
\usepackage[a4paper,margin=2cm,top=3.4cm,bottom=2.8cm,headheight=1.4cm,footskip=1.3cm]{geometry}
\usepackage{amsmath,amssymb,amsfonts}
\usepackage{mathtools} % \xmapsto, \coloneqq... souvent utilisés par les modèles
\usepackage{cancel} % simplifications barrées (bilans de forces)
% \abs{z} (module, valeur absolue) et \norm{u} (norme), utilisés par les modèles
\providecommand{\abs}{}\let\abs\relax\DeclarePairedDelimiter{\abs}{\lvert}{\rvert}
\providecommand{\norm}{}\let\norm\relax\DeclarePairedDelimiter{\norm}{\lVert}{\rVert}
\usepackage[table]{xcolor} % [table] : fournit aussi \rowcolors (alternance de lignes)
\usepackage{pagecolor}
\usepackage{tikz}
\usetikzlibrary{arrows.meta,positioning,calc,shapes.geometric,shapes.misc,shapes.arrows,decorations.pathmorphing,decorations.pathreplacing,decorations.markings,patterns,shadows,mindmap,trees,fit,backgrounds,matrix,intersections,angles,quotes}
\usepackage{pgfplots}
\usepackage[version=4]{mhchem} % équations nucléaires \ce{^{235}_{92}U}
\usepackage[european,siunitx]{circuitikz} % schémas électriques
\pgfplotsset{compat=1.18}
\usepgfplotslibrary{fillbetween,groupplots} % aires entre courbes (calcul intégral)
\usepackage{enumitem}
\usepackage{fancyhdr}
% [most] charge toutes les bibliothèques tcolorbox (listings, minted, raster,
% documentation...) et gonfle inutilement la mémoire TeX sur un document long
% (~300 boîtes) : on ne charge que ce qui est réellement utilisé.
\usepackage[skins,breakable]{tcolorbox}
\usepackage{graphicx}
\usepackage{colortbl}
\usepackage{booktabs}
\usepackage{tabularx}
\usepackage{multirow}
\usepackage{diagbox} % case d'en-tête diagonale des tableaux à double entrée
% Colonnes extensibles centrée / gauche / droite (C, L, R), souvent utilisées par les modèles
\newcolumntype{C}{>{\centering\arraybackslash}X}
\newcolumntype{L}{>{\raggedright\arraybackslash}X}
\newcolumntype{R}{>{\raggedleft\arraybackslash}X}
\usepackage{array}
\usepackage{multicol}
\usepackage{siunitx}
% parse-numbers=false : accepte aussi \SI{2,5 \times 10^{-3}}{...}
\sisetup{locale=FR,output-decimal-marker={,},group-minimum-digits=4,parse-numbers=false}
\DeclareSIUnit\ohm{\ensuremath{\symup{\Omega}}} % U+2126 absent de la police
\DeclareSIUnit\division{div} % réglages d'oscilloscope (ms/div, V/div)
\DeclareSIUnit\tr{tr} % tours (vitesse de rotation, ex. \SI{1500}{\tr\per\minute})
\DeclareSIUnit\molar{M} % concentration molaire (ex. \SI{3}{\molar}, \SI{1}{\micro\molar})
\DeclareSIUnit\an{an} % année (le modèle confond parfois avec \year, un compteur TeX) — ex. \SI{5}{\kilo\meter\per\an}
\DeclareSIUnit\FCFA{FCFA} % franc CFA (ex. \SI{500}{\FCFA\per\kilo\gram})
\DeclareSIUnit\degreeBrix{\textdegree Brix} % teneur en sucre d'un sirop/jus (réfractométrie)
\DeclareSIUnit\month{mois} % le modèle utilise parfois \month, un compteur TeX, comme unité de durée
\DeclareSIUnit\microliter{\micro\liter} % le modèle écrit parfois \microliter en un seul token
\DeclareSIUnit\million{millions} % dénombrement (ex. \SI{15}{\million\per\mL} spermatozoïdes)
\usepackage{qrcode}
% \degree : commande fréquemment utilisée par les modèles pour un angle ou une
% température en dehors de \SI{}{} (non fournie par les paquets chargés).
\providecommand{\degree}{^{\circ}}
% \qed (fin de démonstration), utilisé par les modèles sans amsthm
\providecommand{\qed}{\hfill$\square$}
% \barre : double barre des valeurs interdites dans un tableau de variations (tkz-tab)
\providecommand{\barre}{\ensuremath{\|}}
% environnement proof (amsthm non chargé) utilisé par les modèles
\ifdefined\proof\else\newenvironment{proof}[1][Démonstration]{\par\noindent\textit{#1.}\ }{\qed\par}\fi
\usepackage{lastpage}
\usepackage{hyperref}
\hypersetup{hidelinks}

% ============================================================
% Palette « Pop » OnBuch+ — mêmes hex que la classe P (plus_ui.dart)
% ============================================================
\definecolor{poporange}{HTML}{F59321}
\definecolor{popdark}{HTML}{E07A0C}
\definecolor{popcream}{HTML}{FFF6E8}
\definecolor{popink}{HTML}{1C1714}
\definecolor{poppurple}{HTML}{7A5AE0}
\definecolor{popgreen}{HTML}{2E9E6A}
\definecolor{poppink}{HTML}{E0457B}
\definecolor{popblue}{HTML}{2F7DE1}
\definecolor{popgold}{HTML}{C98A12}
\definecolor{popmuted}{HTML}{8A7F76}
\definecolor{popline}{HTML}{EADFD2}
\definecolor{popgray}{HTML}{8A7F76}
\colorlet{popgrey}{popgray}
% Alias tolérants (le modèle nomme parfois les couleurs différemment)
\colorlet{popwhite}{white}
\colorlet{popblack}{popink}
\colorlet{popred}{poppink}
\colorlet{popyellow}{popgold}
% Teintes claires (fonds de tableaux/encadrés) — le modèle nomme parfois
% les couleurs avec un suffixe L (ex. popblueL) sans les avoir définies.
\colorlet{poporangeL}{poporange!20}
\colorlet{popdarkL}{popdark!20}
\colorlet{poppurpleL}{poppurple!20}
\colorlet{popgreenL}{popgreen!20}
\colorlet{poppinkL}{poppink!20}
\colorlet{popblueL}{popblue!20}
\colorlet{popgoldL}{popgold!20}
\colorlet{popredL}{popred!20}
\colorlet{popgrayL}{popgray!20}
% Teintes claires pour fonds de tableaux / zones
\colorlet{popblueL}{popblue!10!white}
\colorlet{poporangeL}{poporange!14!white}
\colorlet{popgreenL}{popgreen!12!white}
\colorlet{poppurpleL}{poppurple!12!white}
\colorlet{poppinkL}{poppink!12!white}
\colorlet{popgoldL}{popgold!14!white}

\pagecolor{popcream}

% ============================================================
% Typographie
% ============================================================
\defaultfontfeatures{Ligatures=TeX}
\setmainfont{PlusJakartaSans-Medium.ttf}[Path=../../../fonts/,BoldFont=PlusJakartaSans-Bold.ttf]
% Maths en sans-serif (Fira Math) avec chiffres et lettres droites de Plus
% Jakarta, pour que les formules se fondent dans le texte.
\usepackage[math-style=ISO,bold-style=ISO]{unicode-math}
\setmathfont{FiraMath-Regular.otf}[Path=../../../fonts/]
\setmathfont{PlusJakartaSans-Medium.ttf}[Path=../../../fonts/,range={up/num,up/{Latin,latin},\mathup}]
\usepackage{icomma} % virgule décimale sans espace en mode math
% Caractères mathématiques Unicode absents des polices de texte (Plus Jakarta,
% Archivo Black) : rendus par leur équivalent mathématique, en texte comme en
% formule — sinon ils s'affichent en carré vide (ex. « Arithmétique dans ℤ »).
\usepackage{newunicodechar}
\newunicodechar{ℝ}{\ensuremath{\mathbb{R}}}
\newunicodechar{ℤ}{\ensuremath{\mathbb{Z}}}
\newunicodechar{ℂ}{\ensuremath{\mathbb{C}}}
\newunicodechar{ℕ}{\ensuremath{\mathbb{N}}}
\newunicodechar{ℚ}{\ensuremath{\mathbb{Q}}}
\newunicodechar{𝔻}{\ensuremath{\mathbb{D}}}
\newunicodechar{α}{\ensuremath{\alpha}}
\newunicodechar{β}{\ensuremath{\beta}}
\newunicodechar{γ}{\ensuremath{\gamma}}
\newunicodechar{δ}{\ensuremath{\delta}}
\newunicodechar{ε}{\ensuremath{\varepsilon}}
\newunicodechar{θ}{\ensuremath{\theta}}
\newunicodechar{λ}{\ensuremath{\lambda}}
\newunicodechar{μ}{\ensuremath{\mu}}
\newunicodechar{σ}{\ensuremath{\sigma}}
\newunicodechar{τ}{\ensuremath{\tau}}
\newunicodechar{φ}{\ensuremath{\varphi}}
\newunicodechar{ψ}{\ensuremath{\psi}}
\newunicodechar{ω}{\ensuremath{\omega}}
\newunicodechar{Σ}{\ensuremath{\Sigma}}
% majuscules produites par \MakeUppercase dans les titres (α-aminés -> Α)
\newunicodechar{Α}{\ensuremath{\alpha}}
\newunicodechar{Β}{\ensuremath{\beta}}
\newunicodechar{Γ}{\ensuremath{\Gamma}}
\newunicodechar{Δ}{\ensuremath{\Delta}}
\newunicodechar{Ε}{\ensuremath{\varepsilon}}
\newunicodechar{Θ}{\ensuremath{\Theta}}
\newunicodechar{Λ}{\ensuremath{\Lambda}}
\newunicodechar{Μ}{\ensuremath{\mu}}
\newunicodechar{Τ}{\ensuremath{\tau}}
\newunicodechar{Φ}{\ensuremath{\Phi}}
\newunicodechar{Ψ}{\ensuremath{\Psi}}
\newunicodechar{Ω}{\ensuremath{\Omega}}
\newunicodechar{↦}{\ensuremath{\mapsto}}
\newunicodechar{∈}{\ensuremath{\in}}
\newunicodechar{∉}{\ensuremath{\notin}}
\newunicodechar{∧}{\ensuremath{\wedge}}
\newunicodechar{⇔}{\ensuremath{\Leftrightarrow}}
\newunicodechar{⟺}{\ensuremath{\Longleftrightarrow}}
\newunicodechar{⇒}{\ensuremath{\Rightarrow}}
\newunicodechar{⟹}{\ensuremath{\Longrightarrow}}
\newunicodechar{∘}{\ensuremath{\circ}}
\newunicodechar{∪}{\ensuremath{\cup}}
\newunicodechar{∩}{\ensuremath{\cap}}
\newunicodechar{≡}{\ensuremath{\equiv}}
\newunicodechar{⊂}{\ensuremath{\subset}}
\newunicodechar{⊥}{\ensuremath{\perp}}
\newunicodechar{∃}{\ensuremath{\exists}}
\newunicodechar{∀}{\ensuremath{\forall}}
\newunicodechar{∛}{\ensuremath{\sqrt[3]{\,}}}
\newunicodechar{∜}{\ensuremath{\sqrt[4]{\,}}}
\newunicodechar{⃗}{\ensuremath{{}^{\to}}}
\newunicodechar{ⁿ}{\ifmmode^{n}\else\textsuperscript{\ensuremath{n}}\fi}
\newunicodechar{ˣ}{\ifmmode^{x}\else\textsuperscript{\ensuremath{x}}\fi}
\newunicodechar{ᵇ}{\ifmmode^{b}\else\textsuperscript{\ensuremath{b}}\fi}
\newunicodechar{ʳ}{\ifmmode^{r}\else\textsuperscript{\ensuremath{r}}\fi}
\newunicodechar{ᵘ}{\ifmmode^{u}\else\textsuperscript{\ensuremath{u}}\fi}
\newunicodechar{ʸ}{\ifmmode^{y}\else\textsuperscript{\ensuremath{y}}\fi}
\newunicodechar{ᵃ}{\ifmmode^{a}\else\textsuperscript{\ensuremath{a}}\fi}
\newunicodechar{ᵉ}{\ifmmode^{e}\else\textsuperscript{\ensuremath{e}}\fi}
\newunicodechar{ᵏ}{\ifmmode^{k}\else\textsuperscript{\ensuremath{k}}\fi}
\newunicodechar{ᵖ}{\ifmmode^{p}\else\textsuperscript{\ensuremath{p}}\fi}
\newunicodechar{ᴺ}{\ifmmode^{N}\else\textsuperscript{\ensuremath{N}}\fi}
\newunicodechar{ᶜ}{\ifmmode^{c}\else\textsuperscript{\ensuremath{c}}\fi}
\newunicodechar{ʰ}{\ifmmode^{h}\else\textsuperscript{\ensuremath{h}}\fi}
\newunicodechar{ᵠ}{\ifmmode^{q}\else\textsuperscript{\ensuremath{q}}\fi}
\newunicodechar{ₙ}{\ifmmode_{n}\else\textsubscript{\ensuremath{n}}\fi}
\newunicodechar{ₐ}{\ifmmode_{a}\else\textsubscript{\ensuremath{a}}\fi}
\newunicodechar{ᵢ}{\ifmmode_{i}\else\textsubscript{\ensuremath{i}}\fi}
\newunicodechar{ᵣ}{\ifmmode_{r}\else\textsubscript{\ensuremath{r}}\fi}
\newunicodechar{ₖ}{\ifmmode_{k}\else\textsubscript{\ensuremath{k}}\fi}
\newunicodechar{ᵦ}{\ifmmode_{\beta}\else\textsubscript{\ensuremath{\beta}}\fi}
\newunicodechar{ₓ}{\ifmmode_{x}\else\textsubscript{\ensuremath{x}}\fi}
\newfontfamily\popdisplay{ArchivoBlack-Regular.ttf}[Path=../../../fonts/]
\newfontfamily\poplogo{SpaceGrotesk-Bold.ttf}[Path=../../../fonts/]
\newfontfamily\popsemi{PlusJakartaSans-SemiBold.ttf}[Path=../../../fonts/]
\newfontfamily\popxbold{PlusJakartaSans-ExtraBold.ttf}[Path=../../../fonts/]
\newfontfamily\popxboldls{PlusJakartaSans-ExtraBold.ttf}[Path=../../../fonts/,LetterSpace=3.0]

\setlist[enumerate]{leftmargin=1.7em,itemsep=5pt,topsep=4pt}
\setlist[itemize]{leftmargin=1.7em,itemsep=4pt,topsep=4pt,label={\color{poporange}\textbullet}}
\setlength{\parindent}{0pt}
\setlength{\parskip}{5pt}
\renewcommand{\arraystretch}{1.25}

% pgfplots : style Pop par défaut
\pgfplotsset{
  every axis/.append style={axis line style={popink,line width=1pt},tick style={popink},
    grid style={popline,line width=0.5pt},label style={font=\small},tick label style={font=\footnotesize}},
  popcurve/.style={poporange,line width=1.8pt,no markers,smooth},
}
% Même style disponible dans un \draw TikZ simple (hors pgfplots)
\tikzset{popcurve/.style={poporange,line width=1.8pt}}
\tikzset{popgrid/.style={popline,very thin}}
\colorlet{popcurve}{poporange} % le nom du style est parfois utilisé comme couleur

% ============================================================
% Pill / squiggle
% ============================================================
\newcommand{\poppill}[3]{%
  \tikz[baseline=-0.55ex]\node[draw=popink,line width=1.2pt,rounded corners=2.6pt,%
    fill=#2,inner xsep=6.5pt,inner ysep=3pt]{%
    \popxboldls\fontsize{7}{7}\selectfont\color{#3}\MakeUppercase{#1}};%
}
\newcommand{\popsquiggle}[1]{%
  \tikz[baseline=-0.4ex]{
    \draw[#1,line width=2.6pt,line cap=round]
      plot [smooth] coordinates {(0,0.09) (0.34,0.01) (0.68,0.21) (1.02,0.01) (1.36,0.10)};
  }%
}
% Mot-clé surligné (marqueur orange sous le texte)
\newcommand{\cle}[1]{{\popxbold\color{popdark}#1}}

% ============================================================
% En-tête / pied
% ============================================================
\pagestyle{fancy}
\fancyhf{}
\renewcommand{\headrulewidth}{0pt}
\renewcommand{\footrulewidth}{0pt}
\lhead{\raisebox{-0.15ex}{\poplogo\fontsize{13}{13}\selectfont\color{popink}OnBuch\color{poporange}+}}
\rhead{\poppill{\DOCMATIERE\ \textbullet\ \DOCNIVEAU\ \textbullet\ Leçon \DOCLECON}{white}{popink}}
\lfoot{\popxbold\fontsize{7.6}{7.6}\selectfont\color{popmuted}\MakeUppercase{Cours signé OnBuch+}\ \textbullet\ {\popsemi\DOCTITRE}}
\rfoot{\tikz[baseline=-0.6ex]\node[rounded corners=8.5pt,fill=popink,minimum height=17pt,minimum width=17pt,inner xsep=5pt,inner sep=0]{\popxbold\fontsize{8}{8}\selectfont\color{popcream}\thepage\,/\,\pageref*{LastPage}};}

% ============================================================
% Titres
% ============================================================
\newcommand{\chaptitre}[2]{%
  \begin{tcolorbox}[enhanced,colback=poporange,colframe=popink,boxrule=2.6pt,arc=11pt,
      drop shadow=popink,left=15pt,right=15pt,top=12pt,bottom=11pt,before skip=3pt,after skip=18pt]
    \poppill{#2}{popink}{popcream}\par\vspace{9pt}
    {\raggedright\hyphenpenalty=10000\exhyphenpenalty=10000\popdisplay\fontsize{22}{24}\selectfont\color{popcream}\MakeUppercase{#1}\par}
  \end{tcolorbox}
}
% Section numérotée : pastille orange avec le numéro + titre Archivo
\newcounter{popsec}
\newcommand{\coursec}[1]{%
  \stepcounter{popsec}\setcounter{popsub}{0}%
  \par\vspace{16pt}\noindent
  \begin{minipage}{\linewidth}
  \tikz[baseline=-0.7ex]\node[draw=popink,line width=1.6pt,fill=poporange,rounded corners=4pt,minimum size=22pt,inner sep=0]{\popdisplay\fontsize{12}{12}\selectfont\color{popcream}\thepopsec};%
  \hspace{8pt}{\popdisplay\fontsize{15}{17}\selectfont\color{popink}\MakeUppercase{#1}}\par
  \vspace{4pt}\noindent{\color{poporange}\rule{36pt}{3.6pt}}
  \end{minipage}\par\nopagebreak\vspace{8pt}
}
\newcounter{popsub}
\newcommand{\courssub}[1]{%
  \stepcounter{popsub}%
  \par\vspace{9pt}\noindent{\popxbold\fontsize{11.5}{13}\selectfont\color{popink}{\color{poporange}\thepopsec.\thepopsub}\hspace{6pt}#1}\par\nopagebreak\vspace{4pt}
}

% ============================================================
% Boîtes pédagogiques — même recette : fond blanc, bordure épaisse colorée,
% ombre dure encre, badge plein. Titre optionnel [#1].
% ============================================================
\tcbset{popbase/.style={breakable,enhanced,colback=white,boxrule=2.3pt,arc=9pt,
  shadow={3pt}{-3pt}{0pt}{popink},left=14pt,right=14pt,top=11pt,bottom=11pt,before skip=11pt,after skip=15pt,
  fontupper=\popsemi\fontsize{10.6}{14.5}\selectfont\color{popink},
  fontlower=\popsemi\fontsize{10.6}{14.5}\selectfont\color{popink}}}
% Coche dessinée (le glyphe ✓ n'existe pas dans les polices du document)
\newcommand{\popcheck}[1]{\tikz[baseline=-0.5ex]\draw[#1,line width=1.6pt,line cap=round,line join=round](-0.13,0.0)--(-0.04,-0.09)--(0.13,0.1);}
\providecommand{\coche}{\popcheck{popgreen}}
\providecommand{\resultat}[1]{\par\smallskip\noindent\fcolorbox{popgreen}{popgreenL}{\parbox{\dimexpr\linewidth-2\fboxsep-2\fboxrule\relax}{\textbf{Résultat.} #1}}\par\smallskip}
\newcommand{\popboxhead}[3]{\poppill{#1}{#2}{white}\ifx\relax#3\relax\else\hspace{6pt}{\popxbold\fontsize{10.6}{12}\selectfont\color{popink}#3}\fi\par\vspace{7pt}}

\newtcolorbox{aretenir}[1][]{popbase,colframe=popgreen,before upper={\popboxhead{À retenir}{popgreen}{#1}}}
\newtcolorbox{exemplebox}[1][]{popbase,colframe=poporange,before upper={\popboxhead{Exemple}{poporange}{#1}}}
\newtcolorbox{definition}[1][]{popbase,colframe=popblue,before upper={\popboxhead{Définition}{popblue}{#1}}}
\newtcolorbox{propriete}[1][]{popbase,colframe=poppurple,before upper={\popboxhead{Propriété}{poppurple}{#1}}}
\newtcolorbox{methode}[1][]{popbase,colframe=popgold,before upper={\popboxhead{Méthode}{popgold}{#1}}}
\newtcolorbox{attention}[1][]{popbase,colframe=poppink,before upper={\popboxhead{Attention}{poppink}{#1}}}
\newtcolorbox{experience}[1][]{popbase,colframe=popblue,colback=popblueL,before upper={\popboxhead{Expérience}{popblue}{#1}}}
% « Le savais-tu ? » : carte violette pleine (contexte, culture, Cameroun)
\newtcolorbox{savaistu}[1][]{popbase,colback=poppurple,colframe=popink,
  fontupper=\popsemi\fontsize{10.4}{14.2}\selectfont\color{white},
  before upper={\poppill{Le savais-tu\,?}{white}{poppurple}\ifx\relax#1\relax\else\hspace{6pt}{\popxbold\color{white}#1}\fi\par\vspace{7pt}}}
% Exercice résolu : énoncé en haut, \tcblower puis la solution sur fond crème
\newcounter{popexo}\newcounter{popexr}
\newtcolorbox{exoresolu}[1][]{popbase,colframe=poporange,colbacklower=popcream,
  segmentation style={popink,line width=1.2pt,dashed},
  before upper={\stepcounter{popexr}\popboxhead{Exercice résolu \thepopexr}{poporange}{#1}},
  before lower={\poppill{Solution}{popink}{popcream}\par\vspace{6pt}}}
% Exercice (non résolu en place) — la correction est dans la partie Corrigés
\newtcolorbox{exercice}[1][]{popbase,colframe=popink,
  before upper={\stepcounter{popexo}\popboxhead{Exercice \thepopexo}{popink}{#1}}}
% Corrigé
\newtcolorbox{corrige}[1][]{popbase,colframe=popgreen,colback=popgreenL,
  before upper={\popboxhead{Corrigé}{popgreen}{#1}}}
% Objectifs / prérequis / bilan
\newtcolorbox{objectifs}{popbase,colframe=popink,colback=poporangeL,
  before upper={\popboxhead{Objectifs de la leçon}{popink}{}}}
\newtcolorbox{prerequis}{popbase,colframe=popink,colback=popblueL,
  before upper={\popboxhead{Prérequis}{popblue}{}}}
\newtcolorbox{bilan}{popbase,colframe=popink,colback=white,boxrule=2.8pt,
  before upper={\popboxhead{Fiche bilan}{poporange}{L'essentiel à connaître pour l'examen}}}
% Figure encadrée avec légende
\newtcolorbox{popfigure}[1][]{enhanced,breakable=false,colback=white,colframe=popline,boxrule=1.4pt,arc=8pt,
  left=10pt,right=10pt,top=10pt,bottom=8pt,before skip=10pt,after skip=12pt,halign=center,
  lower separated=false,halign lower=center,
  fontlower=\popsemi\fontsize{9}{11.5}\selectfont\color{popmuted},
  #1}
\newcommand{\legende}[1]{\par\vspace{6pt}{\popsemi\fontsize{9}{11.5}\selectfont\color{popmuted}#1}}

% ============================================================
% Signature OnBuch+
% ============================================================
\newcommand{\poplogoinline}[1]{{\poplogo\fontsize{#1}{#1}\selectfont\color{popink}OnBuch\color{poporange}+}}
\newcommand{\signatureonbuch}{%
  \begin{tcolorbox}[enhanced,colback=popink,colframe=popink,boxrule=2.2pt,arc=10pt,
      drop shadow=poporange,left=14pt,right=14pt,top=10pt,bottom=10pt,before skip=0pt,after skip=0pt]
    \tikz[baseline=-0.7ex]\node[circle,fill=popgreen,draw=popcream,line width=1.3pt,minimum size=22pt,inner sep=0]{\popcheck{white}};%
    \hspace{9pt}\begin{minipage}[c]{0.62\linewidth}
      {\popxbold\fontsize{12}{13}\selectfont\color{popcream}Signé\ \textbullet\ {\poplogo OnBuch\color{poporange}+}}\par\vspace{2pt}
      {\raggedright\popsemi\fontsize{8}{10.5}\selectfont\color{popline}\MakeUppercase{Équipe pédagogique OnBuch+}\\\MakeUppercase{Conforme au programme officiel MINESEC}\par}
    \end{minipage}\hfill
    {\poplogo\fontsize{22}{22}\selectfont\color{popcream}OnBuch\color{poporange}+}
  \end{tcolorbox}%
}

% ============================================================
% Page de garde
% #1 titre · #2 matière · #3 niveau · #4 module · #5 n° leçon · #6 durée ·
% #7 description / objectifs (texte ou itemize)
% ============================================================
\newcommand{\pagegarde}[7]{%
  \thispagestyle{empty}
  \newgeometry{a4paper,margin=1.7cm,top=1.6cm,bottom=1.4cm}%
  \begin{tikzpicture}[remember picture,overlay]
    \fill[poporange!18] ([xshift=-3.2cm,yshift=-3.6cm]current page.north east) circle (4.6cm);
    \fill[poppurple!14] ([xshift=2.4cm,yshift=7.6cm]current page.south west) circle (3.8cm);
  \end{tikzpicture}%
  {\poplogo\fontsize{30}{30}\selectfont\color{popink}OnBuch\color{poporange}+}\hfill
  \raisebox{0.6ex}{\poppill{Cours signé \textbullet\ Premium}{popink}{popcream}}\par
  \vspace{1pt}\popsquiggle{poppurple}\par
  \vspace{8pt}
  \begin{tcolorbox}[enhanced,colback=poporange,colframe=popink,boxrule=2.8pt,arc=13pt,
      drop shadow=popink,left=18pt,right=18pt,top=12pt,bottom=13pt,after skip=10pt]
    \poppill{#2 \textbullet\ #3}{popink}{popcream}\hspace{5pt}\poppill{#4}{white}{popink}\par\vspace{8pt}
    {\popdisplay\fontsize{14}{14}\selectfont\color{popink}LEÇON #5}\par\vspace{4pt}
    {\raggedright\hyphenpenalty=10000\exhyphenpenalty=10000\popdisplay\fontsize{25}{27}\selectfont\color{popcream}\MakeUppercase{#1}\par}
    \vspace{8pt}
    {\popxbold\fontsize{9.5}{9.5}\selectfont\color{popink}\MakeUppercase{Durée conseillée : #6}}
  \end{tcolorbox}
  \begin{tcolorbox}[enhanced,colback=white,colframe=popink,boxrule=2.2pt,arc=10pt,
      drop shadow=popink,left=16pt,right=16pt,top=10pt,bottom=10pt,after skip=8pt]
    \poppill{Dans cette leçon}{poppurple}{white}\par\vspace{5pt}
    \popsemi\fontsize{9.3}{12.6}\selectfont\color{popink}#7
  \end{tcolorbox}
  \noindent\begin{minipage}[t]{0.487\linewidth}
    \begin{tcolorbox}[enhanced,colback=popgreen,colframe=popink,boxrule=2.2pt,arc=10pt,
        drop shadow=popink,left=11pt,right=11pt,top=10pt,bottom=10pt,height=3.1cm,valign=center]
      \tcbox[colback=white,colframe=popink,boxrule=1.3pt,arc=5pt,boxsep=0pt,nobeforeafter,
        left=3pt,right=3pt,top=3pt,bottom=3pt,valign=center]{\qrcode[height=1.9cm]{https://play.google.com/store/apps/details?id=com.luvvix.onbuch&hl=fr}}%
      \hspace{8pt}\begin{minipage}[c]{0.5\linewidth}
        {\popdisplay\fontsize{11}{12.5}\selectfont\color{white}\MakeUppercase{Télécharge OnBuch}}\par\vspace{4pt}
        {\raggedright\popsemi\fontsize{8.4}{11}\selectfont\color{white}Gratuit \textbullet\ Google Play\\Tuteur IA, annales, concours, résultats\par}
      \end{minipage}
    \end{tcolorbox}
  \end{minipage}\hfill
  \begin{minipage}[t]{0.487\linewidth}
    \begin{tcolorbox}[enhanced,colback=poppurple,colframe=popink,boxrule=2.2pt,arc=10pt,
        drop shadow=popink,left=11pt,right=11pt,top=10pt,bottom=10pt,height=3.1cm,valign=center]
      \tikz[baseline=-0.7ex]\node[circle,fill=white,draw=popink,line width=1.3pt,minimum size=1.6cm,inner sep=0]{\popdisplay\fontsize{24}{24}\selectfont\color{poppurple}+};%
      \hspace{8pt}\begin{minipage}[c]{0.6\linewidth}
        {\popdisplay\fontsize{11}{12.5}\selectfont\color{white}\MakeUppercase{Passe à OnBuch+}}\par\vspace{4pt}
        {\raggedright\popsemi\fontsize{8.4}{11}\selectfont\color{white}Tous les cours signés, Tuteur IA illimité, fiches PDF — onglet OnBuch+ de l'app\par}
      \end{minipage}
    \end{tcolorbox}
  \end{minipage}
  \vspace{8pt}
  \popcontenu
  \vfill
  \signatureonbuch
  \restoregeometry
  \newpage
}

% Rangée « Ce cours contient » (page de garde)
\newcommand{\poptile}[3]{% #1 couleur #2 gros texte #3 légende
  \begin{tcolorbox}[enhanced,colback=#1,colframe=popink,boxrule=2pt,arc=9pt,shadow={2.5pt}{-2.5pt}{0pt}{popink},
      left=6pt,right=6pt,top=9pt,bottom=9pt,halign=center,valign=center,height=1.75cm,width=\linewidth,nobeforeafter]
    {\popdisplay\fontsize{12.5}{13}\selectfont\color{white}\MakeUppercase{#2}}\par\vspace{3pt}
    {\centering\popsemi\fontsize{7.8}{9.5}\selectfont\color{white}#3\par}
  \end{tcolorbox}}
\newcommand{\popcontenu}{%
  \par\noindent{\popxboldls\fontsize{7.5}{7.5}\selectfont\color{popmuted}CE COURS SIGNÉ CONTIENT}\par\vspace{7pt}
  \noindent\begin{minipage}[t]{0.235\linewidth}\poptile{popblue}{Cours}{complet, illustré, pas à pas}\end{minipage}\hfill
  \begin{minipage}[t]{0.235\linewidth}\poptile{poporange}{Méthodes}{et exercices résolus}\end{minipage}\hfill
  \begin{minipage}[t]{0.235\linewidth}\poptile{poppink}{Exercices}{type examen + corrigés}\end{minipage}\hfill
  \begin{minipage}[t]{0.235\linewidth}\poptile{popgreen}{Fiche bilan}{l'essentiel à retenir}\end{minipage}\par}

% Sommaire
\newtcolorbox{sommaire}{enhanced,colback=white,colframe=popink,boxrule=2.2pt,arc=10pt,
  drop shadow=popink,left=18pt,right=18pt,top=13pt,bottom=13pt,before skip=6pt,after skip=20pt,
  before upper={\poppill{Sommaire}{poppurple}{white}\par\vspace{9pt}\popsemi\fontsize{10.6}{14.5}\selectfont\color{popink}}}

% Signature de fin de cours — poussée tout en bas de la dernière page
\newcommand{\findecours}{%
  \par\vfill\nopagebreak
  \begin{center}\popsquiggle{poporange}\end{center}\vspace{4pt}
  \signatureonbuch
}
\AtBeginDocument{\def\llbracket{\mathopen{[\mkern-3.5mu[}}\def\rrbracket{\mathclose{]\mkern-3.5mu]}}} % intervalles d'entiers (glyphe ⟦ absent de la police)
% Glyphes absents de Fira Math : redéfinis à partir de glyphes disponibles
\AtBeginDocument{%
  \def\setminus{\mathbin{\backslash}}\let\smallsetminus\setminus
  \def\vdots{\mathord{\vbox{\baselineskip3.2pt\lineskiplimit0pt\kern4pt\hbox{.}\hbox{.}\hbox{.}}}}%
  \def\ddots{\mathinner{\mkern1mu\raise6.5pt\hbox{.}\mkern2mu\raise3.6pt\hbox{.}\mkern2mu\raise0.7pt\hbox{.}\mkern1mu}}%
  \def\triangle{\mathord{\scalebox{0.9}{$\symup{\Delta}$}}}%
  \def\nabla{\mathord{\raisebox{\depth}{\scalebox{1}[-1]{$\symup{\Delta}$}}}}%
  \def\checkmark{\popcheck{popdark}}%
  \def\star{\mathbin{\ast}}\def\bigstar{\mathord{\ast}}%
  \def\diamond{\mathbin{\scalebox{0.6}{$\Diamond$}}}%
  \def\bigcup{\mathop{\mathchoice{\raisebox{-0.3em}{\scalebox{1.7}{$\cup$}}}{\scalebox{1.25}{$\cup$}}{\cup}{\cup}}\displaylimits}%
  \def\bigcap{\mathop{\mathchoice{\raisebox{-0.3em}{\scalebox{1.7}{$\cap$}}}{\scalebox{1.25}{$\cap$}}{\cap}{\cap}}\displaylimits}%
  \def\lesseqqgtr{\mathrel{\lessgtr}}\def\bigoplus{\mathop{\oplus}\displaylimits}%
}
\providecommand{\ssi}{si et seulement si} % abréviation fréquente
\AtBeginDocument{\providecommand{\wideparen}{\overparen}} % arc de cercle

%% FIGFIX-BEGIN v4 — correctifs globaux des figures (docs/onbuch-generation/tools/figfix)
\usepackage{adjustbox}\usepackage{etoolbox}
% toute figure TikZ/pgfplots plus large que la ligne est réduite à la largeur disponible
\BeforeBeginEnvironment{tikzpicture}{\begin{adjustbox}{max width=\linewidth}}
\AfterEndEnvironment{tikzpicture}{\end{adjustbox}}
% légendes pgfplots lisibles par-dessus les courbes
\pgfplotsset{every axis/.append style={legend style={fill=white,fill opacity=0.92,text opacity=1,draw=black!15,inner sep=3pt}}}
% étiquettes d'arêtes (syntaxe edge["..."]) avec fond blanc
\tikzset{every edge quotes/.append style={fill=white,inner sep=1.2pt}}
%% FIGFIX-END
\begin{document}

\pagegarde{Family and social life}{Anglais}{1re Tech}{Anglais – classe de Première technique}{N01}{14 séances (fiche de progression harmonisée)}{Cette leçon développe les compétences communicationnelles en anglais pour interagir efficacement lors des événements communautaires camerounais : fêtes, travaux collectifs, règlements de conflits et activités récréatives. Elle articule grammaire (modaux, temps composés, expression du but/résultat), vocabulaire contextuel et production orale/écrite ancrés dans le vécu local.
\vspace{4pt}
\begin{multicols}{2}\small
\begin{itemize}
\item À la fin de cette leçon, je sais échanger en anglais avec des camarades, voisins, adultes et autorités scolaires sur l'organisation de fêtes communautaires camerounaises.
\item Je sais utiliser correctement les modaux (can/may, can/could, ought to, need, have to, must) pour demander la permission, exprimer la possibilité, l'obligation ou la déduction.
\item Je sais comprendre des textes oraux sur le travail communautaire (appel au « bush clearing », « njangi ») et le règlement de différends par les notables.
\item Je sais employer les structures d'expression du but (to, in order to, so that), du résultat (so... that, such... that) et de la préférence (would rather, prefer).
\item Je sais lire et analyser des textes sur les activités récréatives traditionnelles (lutte traditionnelle, danse, jeux de société) et modernes.
\item Je sais conjuguer et utiliser le present perfect et le past perfect pour relier passé et présent dans des récits d'événements communautaires.
\end{itemize}\end{multicols}}

\begin{sommaire}
\begin{multicols}{2}
\begin{enumerate}
\item 1. Fêtes communautaires : interactions sociales \& vocabulaire
\item 2. Modaux : permission, possibilité, obligation, déduction
\item 3. Travail communautaire \& règlement de conflits : écoute \& vocabulaire
\item 4. Exprimer but, résultat, préférence \& obligation
\item 5. Activités récréatives : lecture, vocabulaire \& présent perfect / past perfect
\item 6. Let's / Let's not \& écriture argumentée
\item 7. Intégration : organisation d'un festival interculturel scolaire
\item Activité d'intégration
\item Méthodes et exercices résolus
\item Exercices
\item Corrigés des exercices
\item Fiche bilan
\end{enumerate}
\end{multicols}
\end{sommaire}

\begin{objectifs}
\begin{itemize}
\item À la fin de cette leçon, je sais échanger en anglais avec des camarades, voisins, adultes et autorités scolaires sur l'organisation de fêtes communautaires camerounaises.
\item Je sais utiliser correctement les modaux (can/may, can/could, ought to, need, have to, must) pour demander la permission, exprimer la possibilité, l'obligation ou la déduction.
\item Je sais comprendre des textes oraux sur le travail communautaire (appel au « bush clearing », « njangi ») et le règlement de différends par les notables.
\item Je sais employer les structures d'expression du but (to, in order to, so that), du résultat (so... that, such... that) et de la préférence (would rather, prefer).
\item Je sais lire et analyser des textes sur les activités récréatives traditionnelles (lutte traditionnelle, danse, jeux de société) et modernes.
\item Je sais conjuguer et utiliser le present perfect et le past perfect pour relier passé et présent dans des récits d'événements communautaires.
\item Je sais construire des phrases avec « Let's / Let's not » pour proposer des actions collectives.
\item Je sais rédiger une composition argumentée sur le rôle pacificateur des activités récréatives dans ma communauté.
\item Je sais mobiliser l'ensemble de ces acquis dans une tâche d'intégration simulant l'organisation d'un festival interculturel au Cameroun.
\end{itemize}
\end{objectifs}

\begin{prerequis}
\begin{itemize}
\item Modaux de base (can, must, should) étudiés en 4ème et 3ème.
\item Present simple, past simple, future simple (will) maîtrisés.
\item Vocabulaire de la famille, de l'école et du quartier (noms de lieux, métiers, relations sociales).
\item Connecteurs logiques simples (and, but, because, so).
\item Notions de base sur la culture camerounaise : chefferies, associations de développement, fêtes traditionnelles.
\end{itemize}
\end{prerequis}

\newpage

\coursec{Fêtes communautaires : interactions sociales \& vocabulaire}

\courssub{Introduction : a festival in Ngousso}

In the Ngousso neighbourhood of Yaoundé, the young people of the association \emph{``Jeunesse Unie''} (United Youth) are organising the maize harvest festival. It is a big project! They must ask the \cle{quarter head} (chef de quartier) for permission to use the community hall. They must calm down two \cle{neighbours} (voisins) who are angry about the noise of the drums. They must plan \cle{recreational activities} (activités récréatives) — football, dancing, singing — to bring the community together. To succeed, they must master the art of communicating with everybody: \cle{mates} (camarades), \cle{adults} (adultes) and \cle{authorities} (autorités).

\textbf{Problematic:} How can we express permission, obligation, purpose and result in English so that this festival strengthens social cohesion instead of creating tensions?

In this first part of the lesson, we will identify the actors of a community festival, learn the vocabulary of celebrations, and practise the polite formulas adapted to each person we speak to.

\courssub{The actors of a community festival}

A festival is never organised by one person. It is a \cle{collective} work. Let us identify the actors with their English names and their roles.

\begin{center}
\begin{tabularx}{\linewidth}{|l|X|X|}\hline
\rowcolor{popblueL} \textbf{Actor (English)} & \textbf{French gloss} & \textbf{Role in the festival} \\ \hline
Mates / peers & camarades, jeunes du même âge & They give ideas, dance, play music, do the cleanup. \\ \hline
Neighbours & voisins & They host visitors, lend chairs, but may complain about noise. \\ \hline
Adults / parents & adultes, parents & They cook, give advice, contribute money. \\ \hline
School authorities & chef d'établissement, surveillants & The Principal authorises the use of the school hall or field. \\ \hline
Traditional authorities & chef de quartier, notable, lamido & They bless the event, settle disputes, give official permission. \\ \hline
Administrative authorities & sous-préfet, maire & They give the official authorisation for public events. \\ \hline
\end{tabularx}
\end{center}

\begin{definition}[Community festival]
A \cle{community festival} is a recurrent public event that celebrates the cultural, agricultural or historical identity of a community. It implies collective preparation (planning, fundraising, rehearsals) and intergenerational participation: young people, adults and elders all take part.
\end{definition}

\begin{attention}[Mates : un mot à manier avec prudence]
A frequent mistake: in informal British and Australian English, \emph{mates} means ``friends'' or ``peers''. But in other contexts, \emph{cell mates} means ``compagnons de cellule en prison''! In formal writing in Cameroon, prefer \cle{peers} or \cle{fellow youth}. Keep \emph{mates} for friendly spoken exchanges: \emph{``Hey mates, let's meet at 4 p.m.!''}
\end{attention}

\courssub{Vocabulary of festivals in Cameroon}

Cameroon is rich in community festivals. Each region has its own celebration, and each one is an opportunity to interact with mates, neighbours and authorities.

\begin{center}
\begin{tabularx}{\linewidth}{|l|X|X|}\hline
\rowcolor{popgreenL} \textbf{Festival} & \textbf{Place / People} & \textbf{What is celebrated} \\ \hline
Ngondo & Douala (Sawa people) & Water cult, canoe races, messages of the ancestors (December). \\ \hline
Ngoun & Foumban (Bamoun people) & Great cultural and royal festival of the Bamoun kingdom. \\ \hline
Nyem-Nyem & Ngaoundéré (Adamaoua) & Festival of the Nyem-Nyem people: dances, horses, traditions. \\ \hline
Medumba & Bangangté (West) & Cultural festival of the Medumba people. \\ \hline
Maize festival & Bafia (Centre) & Harvest festival: the new maize is presented and shared. \\ \hline
Cultural week & Schools everywhere & Dances, sketches, traditional meals, Miss/Mister culture. \\ \hline
\end{tabularx}
\end{center}

\begin{popfigure}
\begin{center}
\begin{tikzpicture}[scale=1.0]
% Contour très simplifié du Cameroun
\draw[thick,popdark,fill=popcream] (1.2,0) -- (3.2,0) -- (3.6,0.8) -- (3.4,1.8) -- (4.2,2.6) -- (4.6,3.6) -- (4.2,4.6) -- (3.2,5.4) -- (2.4,5.2) -- (2.2,6.4) -- (1.6,7.4) -- (1.2,6.6) -- (0.8,5.6) -- (0.4,4.6) -- (0.2,3.4) -- (0.0,2.4) -- (0.6,1.4) -- cycle;
% Flèche du nord
\draw[->,thick,popdark] (4.9,6.6) -- (4.9,7.3) node[above]{\textbf{N}};
% Villes et festivals
\fill[popblue] (0.9,1.1) circle (2.2pt) node[right=1pt,popdark]{\footnotesize \textbf{Douala -- Ngondo} (Dec.)};
\fill[poporange] (2.9,3.1) circle (2.2pt) node[right=1pt,popdark]{\footnotesize \textbf{Foumban -- Ngoun}};
\fill[poppurple] (3.3,4.9) circle (2.2pt) node[right=1pt,popdark]{\footnotesize \textbf{Ngaoundéré -- Nyem-Nyem}};
\fill[popgreen] (1.7,2.7) circle (2.2pt) node[left=1pt,popdark]{\footnotesize \textbf{Bangangté -- Medumba}};
\fill[poppink] (1.9,1.9) circle (2.2pt) node[right=1pt,popdark]{\footnotesize \textbf{Bafia -- Fête du Maïs}};
\fill[popdark] (1.5,1.5) circle (1.6pt) node[below left=0pt,popdark]{\footnotesize Yaoundé};
\end{tikzpicture}
\end{center}
\legende{Schematic map of Cameroon locating five major community festivals: Ngondo (Douala), Ngoun (Foumban), Nyem-Nyem (Ngaoundéré), Medumba (Bangangté) and the Maize festival (Bafia). Simplified outline, indicative scale.}
\end{popfigure}

\begin{popfigure}
\begin{center}
\begin{tikzpicture}
\begin{axis}[
    ybar, bar width=14pt,
    width=0.95\linewidth, height=6.2cm,
    ymin=0, ymax=550,
    ylabel={Estimated participants (thousands)},
    symbolic x coords={Ngondo, Ngoun, Nyem-Nyem, Medumba, Maize fest.},
    xtick=data,
    x tick label style={font=\footnotesize, rotate=12, anchor=east},
    y tick label style={font=\footnotesize},
    nodes near coords, nodes near coords style={font=\footnotesize, popdark},
    every node near coord/.append style={/pgf/number format/fixed},
    axis lines=left,
    ymajorgrids=true, grid style={popline!40},
]
\addplot[fill=popblue!70, draw=popblue] coordinates {(Ngondo,500) (Ngoun,200) (Nyem-Nyem,80) (Medumba,120) (Maize fest.,60)};
\end{axis}
\end{tikzpicture}
\end{center}
\legende{Estimated number of participants (in thousands) for five Cameroonian festivals, approximate figures from press reports. The Ngondo is by far the biggest gathering.}
\end{popfigure}

\begin{exemplebox}[The Ngondo in figures]
The Ngondo 2023 mobilised about \num{500000} spectators over 3 weeks. The estimated budget was 150 million FCFA. The organisation required 12 sectoral committees and about 800 volunteers (bénévoles). This shows that a festival is a real project, with planning, money management and teamwork.
\end{exemplebox}

\begin{savaistu}[The spirit of ``njangi'']
The word \cle{njangi} (tontine, collective work) comes from the Bassa word \emph{njanji}, which means ``to help one another'' (s'entraider). This spirit of mutual aid still structures the organisation of festivals in the Littoral and the South-West: each family contributes money, food or work, exactly like the young people of ``Jeunesse Unie'' in Ngousso.
\end{savaistu}

\courssub{The stages of a festival : from planning to cleanup}

Organising a festival follows clear stages. Learn this vocabulary chain — you will need it in all your compositions.

\begin{enumerate}
\item \cle{Planning} (planification) : the committee meets, fixes the date, the place and the programme. \emph{``We should plan the festival two months ahead.''}
\item \cle{Fundraising} (collecte de fonds) : members contribute, a njangi is launched, sponsors are contacted. \emph{``We need to raise 200 000 FCFA for the drums and the food.''}
\item \cle{Invitation} : letters are sent to the authorities, posters are designed. \emph{``We would like to invite the Divisional Officer to our festival.''}
\item \cle{Rehearsal} (répétition) : dancers and musicians practise. \emph{``The dance group must rehearse every Saturday.''}
\item \cle{D-day} (le jour J) : the festival takes place. \emph{``On D-day, everybody must be at the hall at 7 a.m.''}
\item \cle{Cleanup} (nettoyage) : the place is cleaned after the event. \emph{``Let's not leave the field dirty; let's clean up together.''}
\end{enumerate}

\courssub{Polite formulas : adapting your language to your interlocutor}

This is the heart of social interaction. In English, as in our communities, you do not speak to the Principal as you speak to your best friend. The level of politeness changes.

\begin{center}
\begin{tabularx}{\linewidth}{|l|X|X|}\hline
\rowcolor{poppurpleL} \textbf{Interlocutor} & \textbf{Formula} & \textbf{Example} \\ \hline
A mate / peer & Hey, can you...? / Let's... & \emph{``Hey Bih, can you bring the speaker tomorrow?''} \\ \hline
An adult / neighbour & Excuse me, Sir/Madam, may I...? / Could you please...? & \emph{``Excuse me, Madam, could you please lend us some chairs?''} \\ \hline
An authority & Good morning + title. We would like to request... / May we...? & \emph{``Good morning, Principal. We would like to request the use of the school hall.''} \\ \hline
\end{tabularx}
\end{center}

\begin{propriete}[Modals of permission and possibility -- preview]
\begin{itemize}
\item \textbf{can / may} : permission. \emph{May} is more formal than \emph{can}. \emph{``May we use the hall?''} (to an authority) vs \emph{``Can I borrow your pen?''} (to a mate).
\item \textbf{can / could} : possibility. \emph{Could} is more polite/tentative. \emph{``Could you help us?''}
\item \textbf{ought to / should} : advice, moral obligation. \emph{``We ought to respect the neighbours.''}
\item \textbf{must} : strong obligation (internal) or deduction. \emph{``We must clean up.''} / \emph{``He must be the chief.''}
\item \textbf{have to / don't have to} : external obligation / absence of obligation.
\item \textbf{need / needn't} : necessity / absence of necessity.
\end{itemize}
Detailed practice in Section 2.
\end{propriete}

\begin{methode}[How to address an authority]
To make a request to an authority (Principal, quarter head, Divisional Officer), follow this five-step structure :
\begin{enumerate}
\item \textbf{Greeting} : \emph{Good morning,} + \textbf{title} : \emph{Principal / Sir / Chief.}
\item \textbf{Polite request} with \emph{may / could / would} : \emph{We would like to request... / May we use...?}
\item \textbf{Justification} : \emph{...because we are organising the harvest festival for the youth of the quarter.}
\item \textbf{Practical details} : date, time, number of people.
\item \textbf{Anticipated thanks} : \emph{Thank you very much for your kind consideration.}
\end{enumerate}
\end{methode}

\begin{exemplebox}[Three typical scenarios]
\textbf{1. Asking for the community hall:} \emph{``Good morning, Chief. We are the members of `Jeunesse Unie'. May we use the community hall on Saturday 14th for our maize festival? We shall clean it after the event. Thank you, Sir.''}

\textbf{2. Negotiating the sound volume with neighbours:} \emph{``Excuse me, Mr Essomba. We are sorry about the noise. Could we agree on this: we will reduce the volume after 9 p.m.?''}

\textbf{3. Coordinating rehearsals with mates:} \emph{``Hey mates, let's rehearse on Wednesday at 4 p.m. behind the school. Don't be late!''}
\end{exemplebox}

\courssub{Document analysis : a real poster}

Here is the text of a poster for the \emph{Fête du Maïs} in Bafia. Read it and extract the key information.

\begin{exemplebox}[Poster : Bafia Maize Festival]
\textbf{BAFIA MAIZE FESTIVAL -- 12th EDITION}\\
Date: Saturday 21 and Sunday 22 December, from 9 a.m.\\
Place: Bafia Municipal Stadium\\
Organisers: Bafia Youth Association, with the support of the City Council\\
Activities: traditional dances, maize dishes competition, football match, canoe race on the Mbam river, Miss Maize election.\\
Contacts: President 6XX XX XX XX -- free entrance for children under 10.
\end{exemplebox}

\begin{exoresolu}[Extracting information from the poster]
Énoncé : Answer in full sentences. (a) When does the festival take place? (b) Where does it take place? (c) Who organises it? (d) Name two activities. (e) Who can enter free of charge?
\tcblower
Solution détaillée :
\begin{enumerate}
\item[(a)] The festival takes place on Saturday 21 and Sunday 22 December, from 9 a.m.
\item[(b)] It takes place at the Bafia Municipal Stadium.
\item[(c)] It is organised by the Bafia Youth Association, with the support of the City Council.
\item[(d)] Two activities are the traditional dances and the maize dishes competition (also: football match, canoe race, Miss Maize election).
\item[(e)] Children under 10 can enter free of charge.
\end{enumerate}
Method: to extract information, first scan the poster for \cle{wh-words} answers: date (when), place (where), organisers (who), activities (what).
\end{exoresolu}

\courssub{Guided role-play : the organising committee}

A festival is managed by a \cle{committee} (comité). Each member has a function and must speak accordingly.

\begin{popfigure}
\begin{center}
\begin{tikzpicture}[
  box/.style={draw=popdark, fill=popblueL, rounded corners=2pt, align=center, font=\footnotesize, minimum width=3.1cm, minimum height=0.85cm},
  sub/.style={draw=popdark, fill=popgreenL, rounded corners=2pt, align=center, font=\footnotesize, minimum width=2.5cm, minimum height=0.85cm},
  link/.style={thick, popdark}]
\node[box, fill=poporangeL, align=center] (pres) at (0,0){\textbf{President}\\ (porte-parole)};
\node[box, align=center] (sec) at (-4.2,-1.8){\textbf{Secretary}\\ writes letters,\\ keeps minutes};
\node[box, align=center] (tres) at (4.2,-1.8){\textbf{Treasurer}\\ manages funds,\\ njangi};
\node[sub, align=center] (log) at (-4.5,-3.6){Logistics\\ commissioner};
\node[sub, align=center] (com) at (-1.5,-3.6){Communication\\ commissioner};
\node[sub, align=center] (secu) at (1.5,-3.6){Security\\ commissioner};
\node[sub, align=center] (anim) at (4.5,-3.6){Entertainment\\ commissioner};
\draw[link] (pres) -- (sec);
\draw[link] (pres) -- (tres);
\draw[link] (pres) -- (0,-1.2) -- (0,-1.5);
\draw[link] (0,-1.5) -- (-4.5,-1.5) -- (log.north);
\draw[link] (0,-1.5) -- (-1.5,-1.5) -- (com.north);
\draw[link] (0,-1.5) -- (1.5,-1.5) -- (secu.north);
\draw[link] (0,-1.5) -- (4.5,-1.5) -- (anim.north);
\end{tikzpicture}
\end{center}
\legende{Organisation chart of a typical festival committee: the President (spokesperson) coordinates the Secretary, the Treasurer and four commissioners (Logistics, Communication, Security, Entertainment).}
\end{popfigure}

\begin{experience}[Simulating a committee meeting]
\textbf{Material:} role cards (President, Secretary, Treasurer, Logistics commissioner, spokesperson), the Bafia poster, a meeting agenda.

\textbf{Protocol:}
\begin{enumerate}
\item Form groups of five and distribute the roles.
\item The President opens the meeting: \emph{``Good afternoon, mates. Let's plan our harvest festival.''}
\item The Treasurer reports: \emph{``We have raised 85 000 FCFA, but we need 120 000. We should organise a njangi.''}
\item The Logistics commissioner asks a mate: \emph{``Can you book the drums?''} then an adult: \emph{``Excuse me, Sir, could you lend us your generator?''}
\item The spokesperson rehearses the request to the quarter head, using the five-step method.
\item The Secretary writes the minutes (compte rendu) in English.
\end{enumerate}

\textbf{Observation:} the level of language changes with the interlocutor: \emph{can} with mates, \emph{could/may} with adults, \emph{would like to request} with authorities.

\textbf{Interpretation:} choosing the right register is a social skill: it shows respect, obtains permissions, and prevents conflicts — exactly what the ``Jeunesse Unie'' needs in Ngousso.
\end{experience}

\begin{exercice}[Practising interactions]
\begin{enumerate}
\item Rewrite these sentences with the correct level of politeness. (a) To a mate: ask him to bring chairs. (b) To a neighbour: ask her to reduce her generator noise. (c) To the Principal: ask for the school field.
\item Complete the dialogue with: \emph{may, can, would like, let's}.\\
-- A: Good morning, Chief. We \ldots\ request your blessing for our festival.\\
-- B: You are welcome, my children. \ldots\ you organise it peacefully?\\
-- C: Yes, Sir. \ldots\ meet next week to prepare everything!
\item Match the stage to the definition: planning, fundraising, rehearsal, D-day, cleanup / (i) the event itself, (ii) collecting money, (iii) preparing the programme, (iv) practising dances, (v) washing and tidying after the event.
\end{enumerate}
\end{exercice}

\begin{corrige}[Exercice : Practising interactions]
\begin{enumerate}
\item (a) \emph{Hey, can you bring some chairs, please?} (b) \emph{Excuse me, Madam, could you please reduce the noise of your generator?} (c) \emph{Good morning, Principal. We would like to request the use of the school field for our festival.}
\item \emph{may} (or \emph{would like to request}: ``We would like to request''), \emph{Can}, \emph{Let's}.
\item planning -- (iii) ; fundraising -- (ii) ; rehearsal -- (iv) ; D-day -- (i) ; cleanup -- (v).
\end{enumerate}
\end{corrige}

\begin{aretenir}[Key points of this section]
\begin{itemize}
\item A \cle{community festival} celebrates the identity of a community and requires collective preparation: planning, fundraising, invitation, rehearsal, D-day, cleanup.
\item The actors are \cle{mates/peers}, \cle{neighbours}, \cle{adults}, \cle{school authorities} and \cle{traditional authorities}.
\item Adapt your politeness: \emph{can} with mates, \emph{could/may} with adults, \emph{we would like to request} with authorities.
\item To address an authority: Greeting + title + polite request + justification + anticipated thanks.
\item Cameroonian festivals (Ngondo, Ngoun, Nyem-Nyem, Medumba, Maize festival) are real projects, organised by committees in the spirit of \cle{njangi} (mutual aid).
\item Grammar preview: modals \emph{can/may} (permission), \emph{can/could} (possibility), \emph{ought to/should} (advice), \emph{must} (obligation/deduction), \emph{have to/don't have to}, \emph{need/needn't}. Detailed study in Section 2.
\item \emph{Let's / Let's not} + bare infinitive for suggestions (e.g., \emph{Let's rehearse}, \emph{Let's not leave the field dirty}).
\end{itemize}
\end{aretenir}

\coursec{Modaux : permission, possibilité, obligation, déduction}

Dans la section précédente, nous avons acquis le vocabulaire essentiel pour décrire les fêtes communautaires, les rôles des participants (le chef traditionnel, les notables, la jeunesse, l'administration) et les préparatifs matériels. Désormais, pour \emph{interagir} efficacement dans ces contextes — demander l'autorisation d'utiliser la salle des fêtes, négocier le prêt d'un groupe électrogène, conseiller un camarade sur son comportement, ou déduire ce qui s'est passé en notre absence — nous avons besoin d'outils grammaticaux précis : les \textbf{verbes modaux}. En anglais technique et professionnel, le choix du modal détermine le niveau de politesse, la clarté de l'instruction et la crédibilité de l'interlocuteur. Cette section les passe en revue un par un, avec des contextes camerounais authentiques.

\courssub{Permission : \texorpdfstring{\cle{can} (informel) vs \cle{may} (formel)}{can vs may}}

Lorsque vous vous adressez à un pair (camarade, voisin, frère), \cle{can} est la norme. Face à une autorité (proviseur, chef de quartier, chef traditionnel, parent), \cle{may} marque le respect dû au rang.

\begin{exemplebox}[Demander l'accès à la salle des fêtes]
\textbf{Contexte :} Le délégué de classe s'adresse au proviseur pour le festival interculturel.
\begin{itemize}
    \item \textbf{Élève à élève :} \emph{``Can I use the microphone for the rehearsal?''} (Tu peux / Est-ce que je peux utiliser le micro pour la répétition ?) — \textit{Registre familier / standard.}
    \item \textbf{Élève à proviseur :} \emph{``May we have the key to the school hall, Principal?''} (Pourrions-nous avoir la clé de la salle des fêtes, Monsieur le Proviseur ?) — \textit{Registre formel, marque de déférence.}
\end{itemize}
\emph{Vocabulaire clé :} \cle{rehearsal} (répétition), \cle{deference} (déférence, respect), \cle{school hall} (salle des fêtes / amphithéâtre).
\end{exemplebox}

\begin{propriete}[\cle{May} : restriction temporelle]
\cle{May} s'utilise \textbf{uniquement au présent} pour exprimer la permission. Pour le passé, on ne dit pas *``I mayed go'' mais \textbf{``I was allowed to go''} (j'ai eu la permission d'aller) ou \textbf{``I had permission to go''}.
\begin{center}
    \begin{tabular}{|l|l|l|}\hline
    \textbf{Temps} & \textbf{Permission accordée} & \textbf{Permission refusée} \\ \hline
    Présent & You \textbf{may} enter. / You \textbf{can} enter. & You \textbf{may not} / \textbf{cannot} enter. \\ \hline
    Passé & You \textbf{were allowed to} enter. & You \textbf{were not allowed to} enter. \\ \hline
    \end{tabular}
\end{center}
\end{propriete}


\coursec{Travail communautaire \& règlement de conflits : écoute \& vocabulaire}

In the previous section, we learnt how to use \cle{modals} to ask for permission (\emph{May I join the work?}), to express obligation (\emph{Everyone must attend}) and to make deductions (\emph{He must be the elder}). These tools are exactly what we need now, because in this section we enter the heart of village life in Cameroon: \cle{communal work} and the \cle{settlement of community disputes}. When a village calls its people to clear the bush, or when two neighbours quarrel about the boundary of a farm, English is one of the languages used to announce, to discuss, to judge and to reconcile. Let us learn how.

\courssub{Key vocabulary: communal work and dispute settlement}

\begin{definition}[Communal work and dispute settlement]
\cle{Communal work} (or \emph{community work}) is unpaid collective work done by the members of a community for the common good: clearing the village square, repairing a bridge, digging a water source. A \cle{dispute} is a serious disagreement or conflict between people (a \emph{grievance} is a complaint, something you feel is unfair). To \cle{settle a dispute} means to find a solution accepted by both sides, often with the help of a \cle{mediator} (in traditional society, an \emph{elder} or a \emph{notable}).
\end{definition}

\begin{center}
\begin{tabularx}{\linewidth}{|l|X|X|}
\hline
\rowcolor{popblueL} \textbf{English term} & \textbf{Meaning} & \textbf{French gloss} \\
\hline
bush clearing & cutting grass and cleaning public places & débroussaillage \\
call to community work & official announcement summoning people to work & appel au travail communautaire \\
fine / penalty & money (or goods) paid as a punishment & amende / sanction \\
mediator / elder & person who helps two sides to agree & médiateur / ancien \\
dispute / grievance & conflict / complaint & conflit / grief \\
consensus & general agreement of everybody & consensus \\
apology & words said to show regret (« I am sorry ») & excuse \\
reconciliation & restoration of friendship after a conflict & réconciliation \\
peace plant & plant (dracaena) used to seal peace & plante de paix \\
\hline
\end{tabularx}
\end{center}

\begin{definition}[The peace plant]
The \cle{peace plant} (\emph{Dracaena arborea}) is a symbolic plant used in the Grassfields regions and in the South of Cameroon to seal a reconciliation. During the ceremony, each party holds one end of a leaf or stem; the breaking of the plant seals the agreement: the conflict is broken like the stem, and peace is planted like the dracaena, which grows again even from a small cutting.
\end{definition}

\begin{savaistu}[Did you know? The njangi]
In many Cameroonian villages, the \cle{njangi} (collective work or rotating savings group) is regulated by internal rules. If a member is absent without excuse, the fine can be paid \emph{in kind} --- a bottle of palm oil, a bag of maize, a fowl --- or \emph{in cash}, according to the village's internal regulations. The fine is not a revenge: it is a way to protect the common good.
\end{savaistu}

\courssub{Structure of a call to community work}

A \cle{call to community work} is a short official announcement, read by the town crier or the secretary of the traditional council. It always contains three blocks:

\begin{enumerate}
\item \textbf{Announcement} (\emph{l'annonce}): date, time, place, tools required.
\item \textbf{Motivation} (\emph{la motivation}): why the work matters for the common good.
\item \textbf{Sanction} (\emph{la sanction}): the fine for anyone absent without a valid excuse.
\end{enumerate}

\begin{exemplebox}[A model call to community work]
« \emph{Attention please! The Traditional Council of Bapa informs all the inhabitants that there will be a general bush clearing on \textbf{Saturday 14th June at 7 a.m.} at the \textbf{market square}. Every family \textbf{must} send one person with a \textbf{cutlass, a hoe or a rake}. This work will make our market clean and healthy before the rainy season. Anyone absent without a valid excuse \textbf{will pay a fine of 1 000 FCFA}. Long live our village!} »
\end{exemplebox}

\begin{attention}[Common mistake with « apologise »]
Many students write: « \emph{He apologised me} ». This is \textbf{incorrect}. The verb \emph{to apologise} is followed by the preposition \textbf{to}: « He apologised \textbf{to} me » or « He said sorry \textbf{to} me ». Similarly: « She apologised \textbf{for} her lateness » (for + the reason).
\end{attention}

\courssub{Listening: a boundary dispute settled by the elder}

Here is the script of a typical listening text. Read it aloud in groups of four (one student is Ngong, one is Fopa, one is Mama Rose, one is Elder Tabi), then answer the comprehension questions.

\begin{exemplebox}[Listening script: « The farm boundary »]
\textbf{Elder Tabi:} My children, we are here to settle the dispute between Ngong and Fopa about the boundary of their farms. \emph{Let's hear both sides.} Ngong, speak first.\par
\textbf{Ngong:} Papa, Fopa planted cassava on my land. He crossed the old boundary by three metres!\par
\textbf{Fopa:} It is not true! My father planted there before I was born. That strip of land belongs to my family.\par
\textbf{Elder Tabi:} Do you have a witness?\par
\textbf{Ngong:} Yes. Mama Rose saw where the old boundary stones were placed.\par
\textbf{Mama Rose:} I confirm: the stones were near the big mango tree, not near the stream.\par
\textbf{Elder Tabi:} The council of elders has deliberated. The boundary is the line of the old stones, near the mango tree. Fopa, you \textbf{ought to} remove your cassava after the harvest, and you \textbf{must not} plant there again. \emph{What do you propose} to repair the offence?\par
\textbf{Fopa:} I apologise to Ngong. I offer him two baskets of cassava from this harvest.\par
\textbf{Ngong:} I accept his apology. \emph{We agree that} the mango tree line is the boundary.\par
\textbf{Elder Tabi:} \emph{The matter is settled.} Both of you will now hold the peace plant and swear to live as brothers. After that, we shall share the kola nut.
\end{exemplebox}

\begin{methode}[How to take notes on an oral mediation text]
When you listen to a mediation dialogue, do not write full sentences. Note five points only:
\begin{enumerate}
\item \textbf{Parties in conflict}: who is quarrelling? (Ngong vs Fopa)
\item \textbf{Object of the dispute}: what is it about? (a farm boundary)
\item \textbf{Key arguments}: what does each side claim? (witness, old stones)
\item \textbf{Decision}: what is the verdict? (boundary = mango tree line)
\item \textbf{Sanction / commitment}: fine, apology, oath? (cassava offered, oath on the peace plant)
\end{enumerate}
This five-point grid is exactly what the examiner expects in a listening comprehension at the Probatoire technique.
\end{methode}

\begin{aretenir}[Mediation expressions]
\begin{itemize}
\item \emph{Let's hear both sides.} = Écoutons les deux parties.
\item \emph{What do you propose?} = Que proposes-tu ?
\item \emph{We agree that\ldots} = Nous convenons que\ldots
\item \emph{The matter is settled.} = L'affaire est réglée.
\item \emph{I accept your apology.} = J'accepte tes excuses.
\item \emph{Let us swear on the peace plant.} = Jurons sur la plante de paix.
\end{itemize}
\end{aretenir}

\courssub{The stages of a traditional mediation session}

A traditional mediation follows a fixed order, like a scientific protocol. Observe the timeline below and memorise the seven stages.

\begin{popfigure}
\begin{center}
\begin{tikzpicture}[x=1.55cm, y=1cm]
\draw[-{Stealth}, thick, popdark] (0,0) -- (7.6,0) node[below left, font=\small] {time};
\foreach \x/\num/\txt/\col in {
 0.4/1/{Call}/popblue,
 1.5/2/{Facts}/popgreen,
 2.6/3/{Testimonies}/poporange,
 3.7/4/{Deliberation}/poppurple,
 4.8/5/{Verdict}/popink,
 5.9/6/{Oath \& peace plant}/popgold,
 7.0/7/{Closing feast}/poppink}{
  \fill[\col] (\x,0) circle (3pt);
  \node[above, font=\footnotesize\bfseries, text=\col] at (\x,0.12) {\num};
  \node[below, font=\footnotesize, text width=1.35cm, align=center] at (\x,-0.12) {\txt};
  \draw[\col!60] (\x,0) -- (\x,0.5);
}
\node[font=\small\itshape, text=popmuted] at (3.8,0.9) {A traditional mediation session, step by step};
\end{tikzpicture}
\end{center}
\legende{Timeline of a traditional mediation: 1. Call of the parties; 2. Statement of the facts; 3. Testimonies of witnesses; 4. Deliberation of the elders; 5. Verdict; 6. Oath and reconciliation on the peace plant; 7. Closing feast (kola nut, food and drink shared by all).}
\end{popfigure}

Notice how the modal verbs of Section 2 appear at each stage: the parties \emph{must} attend (obligation), witnesses \emph{may} speak (permission), the elders \emph{ought to} be fair (advice), and « the thief \emph{must} be ashamed » (deduction).

\courssub{Analysing a real reconciliation report (PV)}

After the oath, the secretary writes an official report, called in Cameroon a \emph{procès-verbal} (PV) of reconciliation. Here is an anonymised authentic model. Study its structure.

\begin{exemplebox}[PV of reconciliation (anonymised)]
\textbf{RECONCILIATION REPORT --- Village of B.}\par
Date: 22/10/2022. Place: Palace of the Traditional Council.\par
\textbf{Parties:} Mr A. N. (complainant) vs Mr B. T. (defendant).\par
\textbf{Object:} Destruction of a fence and grazing of goats in a maize farm.\par
\textbf{Facts heard:} Both parties and two witnesses were heard.\par
\textbf{Decision of the Council of Elders:} The defendant shall rebuild the fence within two weeks and pay compensation of 15 000 FCFA. The complainant withdraws his grievance.\par
\textbf{Commitment:} The two parties held the peace plant, exchanged the kola nut and swore to live in peace.\par
Signed: the parties (or their fingerprints), two witnesses, the secretary, the president of the council.
\end{exemplebox}

\begin{propriete}[Structure of a reconciliation report]
A good PV of reconciliation always contains: (1) the \cle{heading} (place, date), (2) the \cle{parties}, (3) the \cle{object} of the dispute, (4) the \cle{facts heard}, (5) the \cle{decision} with sanctions or compensation, (6) the \cle{commitment} (oath, peace plant, kola), (7) the \cle{signatures}. In Cameroon it may be written in English, in French, or with formulas in the local language, and the kola nut exchanged is the living symbol of the agreement.
\end{propriete}

\begin{exemplebox}[A convincing figure]
In 2022, the Bapa community (West Region) settled \textbf{47 land disputes} through traditional mediation; \textbf{0 case} was taken to court. The average cost of a mediation was about \textbf{5 000 FCFA} (kola, drinks, small gifts to the elders), against about \textbf{200 000 FCFA} for a court case (lawyer, transport, files). Mediation is therefore about \textbf{40 times cheaper} ($200\,000 \div 5\,000 = 40$) --- and much faster.
\end{exemplebox}

To confirm these figures, researchers analysed a corpus of 50 reconciliation reports and counted how often the key words appear. The word cloud below shows the result: the bigger the word, the more frequent it is.

\begin{popfigure}
\begin{center}
\begin{tikzpicture}
\node[font=\Huge\bfseries, text=popblue] at (0,0.2) {peace};
\node[font=\huge\bfseries, text=popgreen] at (-2.6,-0.5) {elder};
\node[font=\LARGE\bfseries, text=poporange] at (2.7,-0.4) {land};
\node[font=\Large\bfseries, text=poppurple] at (-1.7,1.1) {kola};
\node[font=\large\bfseries, text=popink] at (1.9,1.2) {fine};
\node[font=\large\bfseries, text=popgold] at (-3.1,0.7) {forgive};
\node[font=\Large\bfseries, text=poppink] at (0.3,-1.2) {unity};
\node[font=\normalsize\itshape, text=popmuted] at (3.3,0.6) {witness};
\node[font=\small\itshape, text=popmuted] at (-0.6,-0.8) {boundary};
\node[font=\small\itshape, text=popmuted] at (-3.4,-1.2) {oath};
\node[font=\footnotesize\itshape, text=popmuted] at (2.2,-1.3) {apology};
\end{tikzpicture}
\end{center}
\legende{Word cloud built from a corpus of 50 reconciliation reports (West and North-West Regions, 2019--2023). Word size is proportional to frequency: \emph{peace} (48 occurrences), \emph{elder} (41), \emph{land} (33), \emph{kola} (26), \emph{fine} (19), \emph{forgive} (17), \emph{unity} (15).}
\end{popfigure}

\courssub{Lexical trap: « settle a dispute » vs « settle in »}

\begin{attention}[One verb, two meanings]
The verb \emph{to settle} changes meaning with the particle:
\begin{itemize}
\item \emph{to \cle{settle a dispute}} = to resolve a conflict. « The elders settled the dispute in one morning. » (résoudre)
\item \emph{to \cle{settle in}} = to install yourself comfortably in a new place. « My uncle settled in Douala last year. » (s'installer)
\end{itemize}
Never translate « s'installer » by \emph{settle a dispute}! Context decides the meaning.
\end{attention}

\begin{exoresolu}[Guided exercise: complete the mediation dialogue]
Complete the dialogue with: \emph{both sides, propose, apology, settled, fine, consensus}.\par
Elder: « Let's hear \ldots\ldots\ldots\ Mr Ali, what do you \ldots\ldots\ldots ? »\par
Mr Ali: « I offer my \ldots\ldots\ldots to Mr Bello and I will pay the \ldots\ldots\ldots of 2 000 FCFA. »\par
Elder: « Is there a \ldots\ldots\ldots ? --- All: Yes! --- Then the matter is \ldots\ldots\ldots ! »
\tcblower
\textbf{Solution:} Let's hear \textbf{both sides}. Mr Ali, what do you \textbf{propose}? --- I offer my \textbf{apology} to Mr Bello and I will pay the \textbf{fine} of 2 000 FCFA. --- Is there a \textbf{consensus}? --- Then the matter is \textbf{settled}!\par
\emph{Explanation:} « both sides » follows « hear »; « propose » is the verb of suggestion used by mediators; « apology » is offered (and accepted); the « fine » is paid; the « consensus » is the agreement of all; « settled » is the past participle in the fixed formula « The matter is settled ».
\end{exoresolu}

\begin{exercice}[Practice]
\textbf{Exercise 1 (Vocabulary).} Match the words to their definitions: bush clearing, mediator, grievance, reconciliation, peace plant.\par
a) a person who helps two enemies to agree; b) a complaint; c) collective cleaning of public places; d) the plant broken to seal peace; e) the return of friendship after a conflict.\par
\textbf{Exercise 2 (Listening).} Listen to the script « The farm boundary » read by your teacher and fill the five-point note grid: parties, object, key arguments, decision, sanction/commitment.\par
\textbf{Exercise 3 (Writing).} Write a call to community work (8--10 lines) for the cleaning of your school campus before the end-of-year festival: announcement (date, time, place, tools), motivation, sanction.
\end{exercice}

\begin{corrige}[Exercise 1]
a) mediator; b) grievance; c) bush clearing; d) peace plant; e) reconciliation.
\end{corrige}

\begin{corrige}[Exercises 2 and 3]
\textbf{Exercise 2:} Parties: Ngong vs Fopa. Object: farm boundary (three metres of land). Key arguments: Ngong claims Fopa crossed the boundary; Fopa claims family ownership; Mama Rose confirms the old stones near the mango tree. Decision: boundary = line of the old stones; Fopa removes the cassava after harvest. Commitment: apology, two baskets of cassava, oath on the peace plant, kola shared.\par
\textbf{Exercise 3 (model answer):} « Attention please! The Principal and the Students' Council inform all students that there will be a general cleaning of the school campus on Friday 20th June at 7 a.m. Every class must bring brooms, rakes and cutlasses. A clean campus will make our end-of-year festival beautiful and healthy. Any class absent without a valid excuse will clean the toilets for one week. Long live our school! »
\end{corrige}

\begin{aretenir}[To remember]
Communal work is announced by a \cle{call to community work} (announcement + motivation + sanction). Disputes are settled by \cle{mediation}: both sides are heard, witnesses testify, the elders deliberate, the verdict is given, and the parties swear on the \cle{peace plant} and share the \cle{kola nut}. The report of the session is the \cle{PV of reconciliation}. Remember: \emph{apologise to somebody}, \emph{settle a dispute} (résoudre) $\neq$ \emph{settle in} (s'installer).
\end{aretenir}

In the next section, we will learn how to express \cle{purpose}, \cle{result}, \cle{preference} and \cle{obligation} --- the grammar you need to explain \emph{why} people take part in communal work: « They clear the bush \textbf{in order to} keep the village clean. »

\coursec{Exprimer but, résultat, préférence \& obligation}

In the previous section, we listened to texts about \cle{communal work} and \cle{settling community disputes}. We heard sentences like: \emph{``The villagers came together so that the dispute could be settled peacefully''} and \emph{``The meeting was so long that some people left early''}. These sentences do more than describe facts: they explain \textbf{why} people act (purpose), \textbf{what happens} as a consequence (result), \textbf{what people prefer} (preference) and \textbf{what people must do} (obligation). In this section, we learn to build these four types of links between ideas. They are essential tools for the Probatoire and Baccalauréat technique, both in the grammar exercise and in essay writing.

\courssub{Expressing purpose: to / in order to / so as to}

\begin{definition}[Purpose — le but]
A \cle{purpose} expression answers the question \emph{why?} or \emph{what for?}. When the subject of the main clause and the subject of the purpose action are \textbf{the same person or thing}, we use:
\begin{itemize}
\item \textbf{to} + infinitive (neutral, very common);
\item \textbf{in order to} + infinitive (more formal, used in speeches and official letters);
\item \textbf{so as to} + infinitive (formal, slightly less frequent).
\end{itemize}
\end{definition}

\begin{exemplebox}[Same subject — purpose with infinitive]
\begin{itemize}
\item We organise this festival \textbf{to promote} unity in our quarter.
\item The students cleaned the market square \textbf{in order to prepare} the celebration.
\item The chief arrived early \textbf{so as to welcome} the guests.
\end{itemize}
French gloss: \emph{in order to} = afin de ; \emph{so as to} = de manière à.
\end{exemplebox}

\begin{attention}[Negative purpose]
To express a negative purpose, do \textbf{not} say *\emph{``to not offend''}. Say: \textbf{so as not to} / \textbf{in order not to} + infinitive.
\begin{itemize}
\item[\checkmark] We spoke quietly \textbf{so as not to disturb} the elders.
\item[$\times$] *We spoke quietly to not disturb the elders.
\end{itemize}
\end{attention}

\courssub{Purpose with a different subject: so that / in order that}

Now observe this sentence from a community meeting in Bertoua: \emph{``We start early so that the chief can attend.''} Who starts early? \emph{We}. Who attends? \emph{The chief}. The two subjects are \textbf{different}, so an infinitive is impossible. We need a full clause introduced by \textbf{so that} or \textbf{in order that}.

\begin{propriete}[Rule: same subject or different subject?]
\begin{itemize}
\item \textbf{Same subject} in both actions $\rightarrow$ \textbf{to / in order to / so as to} + infinitive.\\ \emph{We start early to welcome the chief.} (we / we)
\item \textbf{Different subjects} $\rightarrow$ \textbf{so that / in order that} + subject + \textbf{can/could} (present/future meaning) or \textbf{will/would} (after a past tense).\\ \emph{We start early so that the chief \textbf{can} attend.}\\ \emph{We started early so that the chief \textbf{could} attend.}
\end{itemize}
\end{propriete}

\begin{methode}[The ``so that'' test — but ou résultat ?]
To check whether \emph{so that} expresses \textbf{purpose}:
\begin{enumerate}
\item Replace \emph{so that} by \emph{in order that}.
\item If the sentence is still correct and still answers the question \emph{why?}, it is a purpose clause.
\item If the sentence answers \emph{what happened next?}, it is a result clause (see below).
\end{enumerate}
Example: \emph{``We spoke slowly so that everyone could understand''} $\rightarrow$ \emph{``in order that everyone could understand''} is correct $\rightarrow$ \textbf{purpose}.
\end{methode}

\courssub{Expressing result: so \ldots\ that and such \ldots\ that}

\begin{definition}[Result — la conséquence]
A \cle{result} clause expresses the consequence of a quality or a situation. Two structures:
\begin{itemize}
\item \textbf{so} + adjective / adverb + \textbf{that} + clause;
\item \textbf{such} + (a/an) + (adjective) + noun + \textbf{that} + clause.
\end{itemize}
\end{definition}

\begin{exemplebox}[Result in context]
\begin{itemize}
\item The music was \textbf{so loud that} the neighbours complained. (so + adjective)
\item The dancers performed \textbf{so well that} the crowd applauded for ten minutes. (so + adverb)
\item It was \textbf{such a successful festival that} we will repeat it next year. (such + a + adjective + noun)
\item They made \textbf{such progress that} the dispute was settled in one day. (such + uncountable noun, no article)
\end{itemize}
\end{exemplebox}

\begin{methode}[The ``so \ldots\ that'' test]
If you can replace the whole structure by \emph{``and therefore''} or \emph{``and as a result''}, it is a \textbf{result} clause, not a purpose clause.\\ \emph{The music was so loud that the neighbours complained} $\rightarrow$ \emph{The music was very loud, and therefore the neighbours complained.} Correct $\rightarrow$ \textbf{result}.
\end{methode}

\begin{attention}[so or such?]
Choose according to the word that follows:
\begin{itemize}
\item \textbf{so} + adjective/adverb alone: \emph{so tired, so quickly}.
\item \textbf{such} + noun group: \emph{such a good idea, such noise, such beautiful dances}.
\end{itemize}
Common error: *\emph{``so a good idea''} $\rightarrow$ \textbf{such a good idea}.
\end{attention}

\begin{popfigure}
\begin{center}
\begin{tikzpicture}[font=\small, >=Stealth, node distance=0.5cm]
% PURPOSE
\node[draw=popblue, fill=popblueL, rounded corners, text width=3.1cm, align=center] (p1) {\textbf{PURPOSE}\\ same subject};
\node[draw=popblue, fill=white, rounded corners, right=0.6cm of p1, text width=4.6cm, align=center] (p2) {\textbf{to / in order to / so as to}\\ + infinitive};
\draw[->, popblue, thick] (p1) -- (p2);
\node[draw=popblue, fill=popblueL, rounded corners, below=0.5cm of p1, text width=3.1cm, align=center] (p3) {\textbf{PURPOSE}\\ different subjects};
\node[draw=popblue, fill=white, rounded corners, right=0.6cm of p3, text width=4.6cm, align=center] (p4) {\textbf{so that / in order that}\\ + can / could / will / would};
\draw[->, popblue, thick] (p3) -- (p4);
% RESULT
\node[draw=poporange, fill=poporangeL, rounded corners, below=1.1cm of p3, text width=3.1cm, align=center] (r1) {\textbf{RESULT}\\ cause $\rightarrow$ consequence};
\node[draw=poporange, fill=white, rounded corners, right=0.6cm of r1, text width=4.6cm, align=center] (r2) {\textbf{so} + adj/adv + \textbf{that}\\ \textbf{such} + (a) noun + \textbf{that}};
\draw[->, poporange, thick] (r1) -- (r2);
% PREFERENCE
\node[draw=popgreen, fill=popgreenL, rounded corners, below=1.1cm of r1, text width=3.1cm, align=center] (pr1) {\textbf{PREFERENCE}\\ choice between two actions};
\node[draw=popgreen, fill=white, rounded corners, right=0.6cm of pr1, text width=4.6cm, align=center] (pr2) {\textbf{would rather} + bare infinitive\\ \textbf{prefer} + V-ing/noun + \textbf{to}};
\draw[->, popgreen, thick] (pr1) -- (pr2);
\end{tikzpicture}
\end{center}
\legende{Syntax map: how to choose the right structure. 1. Purpose with the same subject $\rightarrow$ infinitive. 2. Purpose with different subjects $\rightarrow$ \emph{so that} + clause. 3. Result $\rightarrow$ \emph{so/such \ldots\ that}. 4. Preference $\rightarrow$ \emph{would rather} or \emph{prefer \ldots\ to}.}
\end{popfigure}

\courssub{Expressing preference: would rather and prefer}

\begin{definition}[Preference — la préférence]
To say what we like or choose between two options:
\begin{itemize}
\item \textbf{would rather} + bare infinitive (infinitive \textbf{without} \emph{to}); comparison with \textbf{than} + bare infinitive.
\item \textbf{prefer} + noun/V-ing + \textbf{to} + noun/V-ing.
\end{itemize}
\end{definition}

\begin{exemplebox}[Preferences at the festival]
\begin{itemize}
\item I'\textbf{d rather dance than watch}. (would rather + V \ldots\ than + V)
\item We'\textbf{d rather not play} music after midnight. (negative: would rather \textbf{not} + V)
\item I \textbf{prefer traditional wrestling to football}. (prefer + noun + to + noun)
\item She \textbf{prefers singing to dancing}. (prefer + V-ing + to + V-ing)
\end{itemize}
French gloss: \emph{I'd rather} = je préférerais ; \emph{I prefer X to Y} = je préfère X à Y.
\end{exemplebox}

\begin{attention}[The most frequent error]
\emph{Would rather} is followed by the infinitive \textbf{without} \emph{to}:
\begin{itemize}
\item[$\times$] *I'd rather \textbf{to go} home.
\item[\checkmark] I'd rather \textbf{go} home.
\end{itemize}
Also note: after \emph{prefer \ldots\ to}, the word \emph{to} is a preposition, so use a noun or \textbf{V-ing}, never the bare infinitive: *\emph{``I prefer wrestling to play football''} $\rightarrow$ \emph{I prefer wrestling to football} / \emph{I prefer wrestling to playing football}.
\end{attention}

\courssub{External obligation: have to / have got to vs must}

\begin{definition}[Obligation — l'obligation]
\begin{itemize}
\item \textbf{must} + bare infinitive: the obligation comes from the \textbf{speaker} (internal obligation, personal conviction).
\item \textbf{have to} / \textbf{have got to} + bare infinitive: the obligation comes from an \textbf{outside rule or authority} (external obligation: law, school rules, council decisions).
\end{itemize}
\end{definition}

\begin{exemplebox}[Internal or external obligation?]
\begin{itemize}
\item We \textbf{have to get} a permit from the council before the festival. (external: the law requires it)
\item Students \textbf{have to wear} their uniform at the school celebration. (external: school rule)
\item I \textbf{must apologise} to my neighbour; I was wrong. (internal: my own decision)
\end{itemize}
French gloss: \emph{a permit} = une autorisation ; \emph{the council} = la mairie / la communauté urbaine.
\end{exemplebox}

\begin{aretenir}[Summary table]
\begin{center}
\begin{tabularx}{\linewidth}{|l|X|X|}
\hline
\rowcolor{popblueL} \textbf{Function} & \textbf{Structure} & \textbf{Example} \\
\hline
Purpose (same subject) & to / in order to / so as to + V & We meet to plan the festival. \\
\hline
Purpose (different subjects) & so that + can/could/will/would & We meet early so that the chief can come. \\
\hline
Result & so + adj/adv + that & It was so noisy that babies cried. \\
\hline
Result & such + (a) + noun + that & It was such a success that we will repeat it. \\
\hline
Preference & would rather + V (than + V) & I'd rather dance than watch. \\
\hline
Preference & prefer + noun/V-ing + to & I prefer wrestling to football. \\
\hline
Obligation (external) & have to + V & We have to inform the authorities. \\
\hline
Obligation (internal) & must + V & I must respect the elders. \\
\hline
\end{tabularx}
\end{center}
\end{aretenir}

\begin{exemplebox}[A pedagogical illustration — Ngondo festival speeches]
In the opening speeches of the \cle{Ngondo} festival (Douala, Sawa people), purpose clauses (\emph{so that}, \emph{in order that}) dominate because official speeches explain \emph{why} the community acts. Result clauses (\emph{so \ldots\ that}, \emph{such \ldots\ that}) describe the impact of past editions. Preference structures (\emph{would rather}, \emph{prefer}) appear when speakers compare traditional and modern forms of celebration. This mirrors exactly the skills you need in your essays: explain purpose, show results, express preferences.
\end{exemplebox}

\begin{popfigure}
\begin{center}
\begin{tikzpicture}
% Manual pie chart using TikZ (no pgf-pie package)
\def\angleA{0}
\def\radius{2.6}
% Wrestling 35% = 126 degrees
\fill[popblue] (0,0) -- (\angleA:\radius) arc (\angleA:126:\radius) -- cycle;
\node[popblue] at (63:1.8) {35\%};
% Dance 25% = 90 degrees
\def\angleB{126}
\fill[poporange] (0,0) -- (\angleB:\radius) arc (\angleB:216:\radius) -- cycle;
\node[poporange] at (171:1.8) {25\%};
% Football 20% = 72 degrees
\def\angleC{216}
\fill[popgreen] (0,0) -- (\angleC:\radius) arc (\angleC:288:\radius) -- cycle;
\node[popgreen] at (252:1.8) {20\%};
% Board games 15% = 54 degrees
\def\angleD{288}
\fill[poppurple] (0,0) -- (\angleD:\radius) arc (\angleD:342:\radius) -- cycle;
\node[poppurple] at (315:1.8) {15\%};
% Other 5% = 18 degrees
\def\angleE{342}
\fill[popgold] (0,0) -- (\angleE:\radius) arc (\angleE:360:\radius) -- cycle;
\node[popgold] at (351:1.8) {5\%};
% Legend
\begin{scope}[yshift=-3.5cm, font=\small]
\node[anchor=west] at (0,0) {\textcolor{popblue}{\rule{0.5cm}{0.3cm}} Traditional wrestling (35\%)};
\node[anchor=west] at (0,-0.4) {\textcolor{poporange}{\rule{0.5cm}{0.3cm}} Dance (25\%)};
\node[anchor=west] at (0,-0.8) {\textcolor{popgreen}{\rule{0.5cm}{0.3cm}} Football (20\%)};
\node[anchor=west] at (0,-1.2) {\textcolor{poppurple}{\rule{0.5cm}{0.3cm}} Board games (15\%)};
\node[anchor=west] at (0,-1.6) {\textcolor{popgold}{\rule{0.5cm}{0.3cm}} Other (5\%)};
\end{scope}
\end{tikzpicture}
\end{center}
\legende{Fictional survey: ``Favourite recreational activity'' among 200 secondary school students in Bertoua. Wrestling 35\,\% (70 students), dance 25\,\% (50), football 20\,\% (40), board games 15\,\% (30), other 5\,\% (10). Model sentence: \emph{Most students prefer traditional wrestling to football.}}
\end{popfigure}

\begin{savaistu}[Had better — le conseil pressant (bonus)]
\emph{Had better} + bare infinitive expresses strong advice with an implicit negative consequence: \emph{``You'd better ask the chief before using the square''} means \emph{if you don't, there may be trouble}. It is stronger than \emph{should}. Negative form: \emph{You'd better not arrive late.} Note that despite the past form \emph{had}, it refers to the present or future. \emph{Note: this structure is not in the official Première Tech syllabus but appears frequently in authentic texts.}
\end{savaistu}

\begin{exoresolu}[Combining sentences — Probatoire style]
\textbf{Task:} Join each pair of sentences using the word(s) in brackets.
\begin{enumerate}
\item The youths cleaned the village square. They wanted to welcome the delegation. (in order to)
\item The hall was very small. Not everybody could enter. (so \ldots\ that)
\item I like dancing. I like watching football less. (would rather)
\item The meeting started at 6 a.m. The workers could attend before going to the fields. (so that)
\end{enumerate}
\tcblower
\textbf{Detailed solution:}
\begin{enumerate}
\item Same subject (\emph{the youths}) $\rightarrow$ infinitive of purpose: \emph{The youths cleaned the village square \textbf{in order to welcome} the delegation.}
\item Result of a quality (\emph{small} = adjective) $\rightarrow$ \emph{so + adjective + that}: \emph{The hall was \textbf{so small that} not everybody could enter.}
\item Preference between two actions $\rightarrow$ \emph{would rather + bare infinitive + than}: \emph{I \textbf{would rather dance than watch} football.}
\item Different subjects (\emph{the meeting} / \emph{the workers}) $\rightarrow$ \emph{so that + could}: \emph{The meeting started at 6 a.m. \textbf{so that} the workers \textbf{could attend} before going to the fields.}
\end{enumerate}
\end{exoresolu}

\begin{exercice}[Practice — combine and transform]
Join each pair of sentences with the structure in brackets.
\begin{enumerate}
\item The festival was a great success. The mayor promised to finance it next year. (such \ldots\ that)
\item We must ask for a permit. The council requires it. (have to)
\item She prefers storytelling. She prefers it more than singing. (prefer \ldots\ to)
\item The elders spoke slowly. The young people could take notes. (in order that)
\item The drummers played very loudly. The whole village heard them. (so \ldots\ that)
\end{enumerate}
\end{exercice}

\begin{corrige}[Exercice — combine and transform]
\begin{enumerate}
\item It was \textbf{such a successful festival that} the mayor promised to finance it next year. (\emph{such + a + adjective + noun + that})
\item We \textbf{have to ask} for a permit because the council requires it. (external obligation)
\item She \textbf{prefers storytelling to singing}. (\emph{prefer + noun + to + noun})
\item The elders spoke slowly \textbf{in order that} the young people \textbf{could take} notes. (different subjects $\rightarrow$ purpose clause)
\item The drummers played \textbf{so loudly that} the whole village heard them. (\emph{so + adverb + that})
\end{enumerate}
\end{corrige}

\begin{aretenir}[Key points to remember]
\begin{itemize}
\item Purpose, same subject: \textbf{to / in order to / so as to} + infinitive; negative: \textbf{so as not to}.
\item Purpose, different subjects: \textbf{so that / in order that} + can/could/will/would.
\item Result: \textbf{so} + adjective/adverb + \textbf{that}; \textbf{such} + (a) + noun + \textbf{that}.
\item Preference: \textbf{would rather} + infinitive \emph{without} \emph{to}; \textbf{prefer} + noun/V-ing + \textbf{to}.
\item Obligation: \textbf{must} = internal; \textbf{have to} = external (rules, law, authorities).
\end{itemize}
\end{aretenir}

\coursec{Activités récréatives : lecture, vocabulaire \& present perfect / past perfect}

After studying how to express purpose, result, preference and obligation, we now reach the cultural heart of this unit: traditional and modern \cle{recreational activities} (\emph{activités récréatives}). Understanding how Cameroonian communities celebrate, play and dance together will give us authentic content to master two essential tenses for narration and testimony: the \textbf{present perfect} (link between past and present) and the \textbf{past perfect} (an action completed before another past action).

\begin{aretenir}[Objectifs de la séance]
À la fin de cette séquence, tu seras capable de :
\begin{itemize}
\item comprendre des textes sur les loisirs traditionnels et modernes au Cameroun ;
\item utiliser correctement le vocabulaire des activités récréatives ;
\item employer le \cle{present perfect} et le \cle{past perfect} dans un récit ou un témoignage ;
\item rédiger un court paragraphe sur une activité récréative de ta communauté.
\end{itemize}
\end{aretenir}

\courssub{Lecture : trois regards sur les loisirs traditionnels}

Read the three texts below for general understanding. Then find the target vocabulary words (in bold with \cle{...}) and the verb forms in \emph{present perfect} or \emph{past perfect} (in bold).

\begin{exemplebox}[Text 1: Newspaper article --- ``Traditional wrestling, a social cement in the North'' (adapted)]
Every dry season, the villages of the Bénoué area vibrate to the rhythm of \cle{traditional wrestling} (\emph{lutte traditionnelle}). Last Sunday, the \cle{wrestling arena} (\emph{arène de lutte}) of Garoua welcomed more than 5,000 spectators. The \cle{referee} (\emph{arbitre}) whistled the start of the first \cle{round} (\emph{manche}) at exactly 4 p.m. The wrestlers, covered with oil and kaolin, fought with impressive technique. The winner received a \cle{trophy} (\emph{trophée}) offered by the Lamido. \textbf{The festival has become} a symbol of unity. \textbf{Young people have never missed} an edition since 2010. \textbf{By the time the chief arrived}, the wrestlers \textbf{had already performed} the ceremonial dance. \textbf{The crowd had cheered} loudly before the final match started.
\end{exemplebox}

\begin{exemplebox}[Text 2: Youth blog --- ``My first Ngondo'' (fictional testimony)]
My name is Mireille, I am 17 and I live in Douala. \textbf{I have just come back} from the \cle{Ngondo} festival on the Wouri river banks. It \textbf{was} magical! The \cle{drumbeat} (\emph{battement de tam-tam}) of the drums guided the \cle{choreography} (\emph{chorégraphie}) of the dancers. \textbf{I have never seen} such colourful pirogues racing on the water. \textbf{My grandmother had told} me many stories about the \emph{Jengu} water spirits before I \textbf{went} there. \textbf{When the sacred vase emerged}, the initiates \textbf{had already disappeared} into the water. \textbf{Since that day}, I \textbf{have felt} a stronger connection to my Sawa heritage.
\end{exemplebox}

\begin{exemplebox}[Text 3: Cultural report --- ``Ancestral board games: songo and nku'']
In the Centre and South regions, evening gatherings are never silent. Around a wooden \cle{board game} (\emph{jeu de plateau}), two players compete at \cle{songo} (also called \emph{awalé}). Every \cle{seed} (\emph{graine / pion}) counts. The \cle{strategy} (\emph{stratégie}) consists in capturing the opponent's seeds. \textbf{Researchers have studied} this game for decades. \textbf{It has spread} in many variants across Cameroon. \textbf{Long before modern schools existed}, elders \textbf{had used} games like the \emph{nku} to teach counting and quick thinking. \textbf{By the time the funeral wake started}, the players \textbf{had set up} the board to entertain the guests, as tradition recommends.
\end{exemplebox}

\begin{exercice}[Reading comprehension]
Answer these questions in full sentences.
\begin{enumerate}
\item Where does the wrestling festival in Text 1 take place, and when?
\item Who offered the trophy to the winner?
\item In Text 2, what had Mireille's grandmother done before Mireille went to the festival?
\item What is the Ngondo festival connected to?
\item In Text 3, what is the aim of the strategy in the songo game?
\item Find one sentence in each text with the present perfect and one with the past perfect.
\end{enumerate}
\end{exercice}

\begin{corrige}[Exercice : Reading comprehension]
\begin{enumerate}
\item It takes place in Garoua (Bénoué area, North of Cameroon), during the dry season.
\item The Lamido offered the trophy.
\item She had told Mireille many stories about the Jengu water spirits.
\item It is connected to the Sawa people and the Wouri river (water spirits / Jengu cult).
\item The aim is to capture the opponent's seeds.
\item Examples: Text 1 --- present perfect: ``The festival has become a symbol of unity''; past perfect: ``the wrestlers had already performed the ceremonial dance''. Text 2 --- present perfect: ``I have never seen such colourful pirogues''; past perfect: ``My grandmother had told me many stories''. Text 3 --- present perfect: ``It has spread in many variants''; past perfect: ``the players had set up the board''.
\end{enumerate}
\end{corrige}

\courssub{Vocabulaire des activités récréatives}

\begin{popfigure}
\centering
\begin{tikzpicture}[
  mindmap,
  concept color=popblue!80,
  text=white,
  font=\small\bfseries,
  level 1 concept/.append style={level distance=4.2cm, sibling angle=90, concept color=poporange!80},
  level 2 concept/.append style={level distance=3.1cm, sibling angle=45, concept color=popgreen!70, font=\scriptsize\bfseries},
  every node/.style={scale=0.85}
]
\node[concept, align=center]{Recreational\\ Activities}
  child[clockwise from=90] {
    node[concept, align=center]{Traditional\\ Sports}
    child { node[concept, align=center]{wrestling arena\\ (arène de lutte)} }
    child { node[concept, align=center]{referee\\ (arbitre)} }
    child { node[concept, align=center]{round\\ (manche)} }
    child { node[concept, align=center]{trophy\\ (trophée)} }
  }
  child[clockwise from=180] {
    node[concept] {Dance \& Music}
    child { node[concept, align=center]{cheer\\ (acclamer)} }
    child { node[concept, align=center]{drumbeat\\ (battement de tam-tam)} }
    child { node[concept, align=center]{choreography\\ (chorégraphie)} }
    child { node[concept, align=center]{festival\\ (festival)} }
  }
  child[clockwise from=270] {
    node[concept] {Board Games}
    child { node[concept, align=center]{board game\\ (jeu de plateau)} }
    child { node[concept, align=center]{seed\\ (graine / pion)} }
    child { node[concept, align=center]{strategy\\ (stratégie)} }
    child { node[concept] {songo / awalé} }
  }
  child[clockwise from=0] {
    node[concept] {Modern Leisure}
    child { node[concept, align=center]{concert\\ (concert)} }
    child { node[concept, align=center]{cinema\\ (cinéma)} }
    child { node[concept, align=center]{picnic\\ (pique-nique)} }
    child { node[concept, align=center]{talent show\\ (concours de talents)} }
  };
\end{tikzpicture}
\legende{Carte mentale du vocabulaire des activités récréatives (anglais / français / contexte local).}
\end{popfigure}

\begin{center}
\begin{tabularx}{\linewidth}{|l|X|X|}\hline
\rowcolor{popblueL}
\textbf{English Term} & \textbf{Français} & \textbf{Example in context} \\ \hline
wrestling arena & arène de lutte & The \textbf{wrestling arena} was full of spectators. \\ \hline
referee & arbitre & The \textbf{referee} blew the whistle. \\ \hline
round & manche / tour & He won the first \textbf{round}. \\ \hline
trophy & trophée / coupe & She lifted the \textbf{trophy} proudly. \\ \hline
cheer & acclamer & The crowd began to \textbf{cheer}. \\ \hline
drumbeat & battement de tam-tam & The \textbf{drumbeat} accelerated. \\ \hline
choreography & chorégraphie & The \textbf{choreography} was complex. \\ \hline
board game & jeu de plateau & \textbf{Songo} is a popular \textbf{board game}. \\ \hline
seed & graine / pion & Capture the opponent's \textbf{seeds}. \\ \hline
strategy & stratégie & You need a good \textbf{strategy} to win. \\ \hline
heritage & patrimoine & The Ngondo is part of the Sawa \textbf{heritage}. \\ \hline
spectator & spectateur & Five thousand \textbf{spectators} attended. \\ \hline
\end{tabularx}
\end{center}

\begin{exercice}[Vocabulary: gap filling]
Complete the sentences with words from the table above (use the correct form).
\begin{enumerate}
\item The \_\_\_\_\_\_\_\_\_\_\_\_ blew his whistle to stop the match.
\item Our team won the \_\_\_\_\_\_\_\_\_\_\_\_ at the end of the tournament.
\item The \_\_\_\_\_\_\_\_\_\_\_\_ of the tam-tams could be heard from far away.
\item To win at songo, you must capture your opponent's \_\_\_\_\_\_\_\_\_\_\_\_.
\item The dancers repeated their \_\_\_\_\_\_\_\_\_\_\_\_ all afternoon.
\item More than 5,000 \_\_\_\_\_\_\_\_\_\_\_\_ came to the wrestling arena.
\item The Ngondo festival is part of the Sawa cultural \_\_\_\_\_\_\_\_\_\_\_\_.
\item The wrestler lost the first \_\_\_\_\_\_\_\_\_\_\_\_ but won the fight.
\end{enumerate}
\end{exercice}

\begin{corrige}[Exercice : Vocabulary]
\begin{enumerate}
\item referee
\item trophy
\item drumbeat
\item seeds
\item choreography
\item spectators
\item heritage
\item round
\end{enumerate}
\end{corrige}

\courssub{Grammaire 1 : le present perfect}

\begin{definition}[Present perfect --- forme]
\textbf{Affirmative :} sujet + \textbf{have/has} + participe passé (V-ed ou 3\up{e} colonne).\\
\textbf{Negative :} sujet + have/has + \textbf{not} + participe passé.\\
\textbf{Interrogative :} Have/Has + sujet + participe passé ?
\begin{itemize}
\item \emph{I \textbf{have visited} the Ngondo festival twice.}
\item \emph{She \textbf{has never missed} an edition.}
\item \emph{\textbf{Have} you \textbf{ever played} songo?}
\end{itemize}
\end{definition}

\begin{propriete}[Present perfect --- emplois]
On utilise le present perfect pour :
\begin{enumerate}
\item une action passée \textbf{liée au présent} (résultat visible maintenant) : \emph{I have lost my trophy.} (je ne l'ai plus) ;
\item une \textbf{expérience de vie} (avec \emph{ever}, \emph{never}) : \emph{Have you ever seen traditional wrestling?} ;
\item une action qui \textbf{continue jusqu'à maintenant} (avec \emph{since} + point de départ, \emph{for} + durée) : \emph{The festival has existed \textbf{since} 1950. / \textbf{for} more than 70 years.} ;
\item une action \textbf{très récente} (avec \emph{just}, \emph{already}, \emph{yet}) : \emph{I have \textbf{just} come back from the concert.}
\end{enumerate}
\end{propriete}

\begin{attention}[Pièges fréquents]
\begin{itemize}
\item Jamais de present perfect avec un \textbf{moment passé précis} (\emph{yesterday, last year, in 2010}) : on dit \emph{I \textbf{saw} the match yesterday} (past simple), pas \emph{I have seen the match yesterday}.
\item \emph{Since} + point de départ (\emph{since 2010}) ; \emph{for} + durée (\emph{for ten years}).
\item N'oublie pas le participe passé : \emph{gone} (et non \emph{goed}), \emph{seen}, \emph{taken}, \emph{won}.
\end{itemize}
\end{attention}

\courssub{Grammaire 2 : le past perfect}

\begin{definition}[Past perfect --- forme et emploi]
\textbf{Forme :} sujet + \textbf{had} + participe passé.\\
\textbf{Emploi :} le past perfect exprime une action \textbf{accomplie avant} une autre action passée (antériorité). L'autre action est au past simple.
\begin{itemize}
\item \emph{By the time the chief \textbf{arrived}, the wrestlers \textbf{had performed} the dance.} (d'abord la danse, ensuite l'arrivée du chef)
\item \emph{My grandmother \textbf{had told} me stories before I \textbf{went} to the festival.}
\end{itemize}
Mots-repères : \emph{before, after, by the time, when, already, just}.
\end{definition}

\begin{exemplebox}[Chronologie : qui fait quoi en premier ?]
\emph{When the referee arrived, the match had already started.}
\begin{enumerate}
\item First action: the match started. $\rightarrow$ \textbf{past perfect} (\emph{had started})
\item Second action: the referee arrived. $\rightarrow$ \textbf{past simple} (\emph{arrived})
\end{enumerate}
Astuce : l'action au past perfect est toujours la \textbf{première} dans la chronologie réelle.
\end{exemplebox}

\begin{methode}[Choisir entre past simple, present perfect et past perfect]
\begin{enumerate}
\item L'action a-t-elle un lien avec le présent (résultat, expérience, \emph{since/for}) ? $\rightarrow$ \textbf{present perfect}.
\item L'action est-elle située à un moment passé précis et terminé ? $\rightarrow$ \textbf{past simple}.
\item L'action est-elle antérieure à une autre action passée ? $\rightarrow$ \textbf{past perfect}.
\end{enumerate}
\end{methode}

\begin{exoresolu}[Conjugaison : au bon temps]
Énoncé : Mets les verbes entre parenthèses au temps correct (past simple, present perfect ou past perfect).\\
1. The dancers (rehearse) \_\_\_\_\_\_\_\_\_\_\_\_ for two hours before the show (begin) \_\_\_\_\_\_\_\_\_\_\_\_.\\
2. I (never / see) \_\_\_\_\_\_\_\_\_\_\_\_ a pirogue race in my life.\\
3. Last year, our school (organise) \_\_\_\_\_\_\_\_\_\_\_\_ a cultural day.
\tcblower
Solution détaillée :
\begin{enumerate}
\item \textbf{had rehearsed} --- \textbf{began}. La répétition est antérieure au début du spectacle (repère \emph{before}) : past perfect pour la première action, past simple pour la seconde.
\item \textbf{have never seen}. Expérience de vie avec \emph{never}, sans date précise : present perfect.
\item \textbf{organised}. Moment passé précis (\emph{last year}) : past simple obligatoire.
\end{enumerate}
\end{exoresolu}

\begin{exercice}[Grammar: choose the correct tense]
Put the verbs in brackets in the correct tense (past simple, present perfect or past perfect).
\begin{enumerate}
\item The Ngondo festival (exist) \_\_\_\_\_\_\_\_\_\_\_\_ for many decades.
\item When we (arrive) \_\_\_\_\_\_\_\_\_\_\_\_ at the arena, the first round (already / finish) \_\_\_\_\_\_\_\_\_\_\_\_.
\item (you / ever / play) \_\_\_\_\_\_\_\_\_\_\_\_ songo with your grandfather?
\item The crowd (cheer) \_\_\_\_\_\_\_\_\_\_\_\_ loudly when the winner lifted the trophy last Sunday.
\item She (just / come back) \_\_\_\_\_\_\_\_\_\_\_\_ from the village festival.
\item Before the referee (blow) \_\_\_\_\_\_\_\_\_\_\_\_ the whistle, the wrestlers (enter) \_\_\_\_\_\_\_\_\_\_\_\_ the arena.
\end{enumerate}
\end{exercice}

\begin{corrige}[Exercice : Grammar]
\begin{enumerate}
\item \textbf{has existed} (durée jusqu'à maintenant, \emph{for} + durée : present perfect).
\item \textbf{arrived} --- \textbf{had already finished} (la fin de la manche est antérieure à l'arrivée).
\item \textbf{Have you ever played} (expérience de vie : present perfect).
\item \textbf{cheered} (moment passé précis : \emph{last Sunday}).
\item \textbf{has just come back} (action très récente avec \emph{just}).
\item \textbf{blew} --- \textbf{had entered} (les lutteurs entrent avant le coup de sifflet).
\end{enumerate}
\end{corrige}

\courssub{Expression écrite : un paragraphe-témoignage}

\begin{methode}[Rédiger un court témoignage sur une activité récréative]
\begin{enumerate}
\item Présente l'activité et le lieu (\emph{Last month, I attended...}).
\item Raconte les événements au past simple.
\item Utilise au moins un past perfect pour l'antériorité (\emph{By the time we arrived, ... had already started}).
\item Termine par un present perfect pour le bilan ou l'effet présent (\emph{Since that day, I have felt...}).
\item Utilise au moins quatre mots du vocabulaire de la séance.
\end{enumerate}
\end{methode}

\begin{exercice}[Writing task]
Write a short paragraph (8--10 lines) about a recreational activity you have taken part in (a festival, a match, a dance, a board game competition). Use the past simple, at least one past perfect, one present perfect, and four words from the vocabulary list.
\end{exercice}

\begin{corrige}[Exercice : Writing --- exemple de production]
Last December, I attended the wrestling festival in my village. When my friends and I arrived at the \textbf{wrestling arena}, the first \textbf{round} had already started. The \textbf{referee} was watching the two wrestlers carefully. The \textbf{drumbeat} of the tam-tams made everybody excited, and the crowd began to \textbf{cheer}. At the end, the strongest wrestler received a beautiful \textbf{trophy}. I have never seen such an exciting event in my life. Since that day, I have decided to train every week, and I have already won two matches in my neighbourhood.
\end{corrige}

\begin{aretenir}[À retenir]
\begin{itemize}
\item \textbf{Present perfect} = have/has + participe passé : lien avec le présent, expérience (\emph{ever/never}), durée (\emph{since/for}), action récente (\emph{just/already/yet}).
\item \textbf{Past perfect} = had + participe passé : action accomplie \textbf{avant} une autre action passée (\emph{before, by the time, when}).
\item Jamais de present perfect avec une date passée précise (\emph{yesterday, last year}) $\rightarrow$ past simple.
\item Vocabulaire clé : \emph{wrestling arena, referee, round, trophy, cheer, drumbeat, choreography, board game, seed, strategy, heritage, spectator}.
\end{itemize}
\end{aretenir}

\coursec{Let's / Let's not \& écriture argumentée}

\courssub{Transition depuis la section 5}
In Section 5, we explored \cle{recreational activities} (activités récréatives) through reading texts and the \cle{present perfect} / \cle{past perfect} tenses. We saw how events like traditional wrestling (\emph{lutte traditionnelle}), \emph{njangi} meetings, or the \emph{Ngondo} festival create social bonds. Now, we move from \emph{describing} these activities to \emph{proposing} them and \emph{arguing} for their value. To do this, we need two tools: the structure \textbf{Let's / Let's not} for making inclusive suggestions, and the \textbf{argumentative essay} structure to convince an audience (school authorities, community leaders) that these activities bring peace.

\courssub{La structure \emph{Let's / Let's not} : proposer et décider ensemble}

\begin{definition}[Forme affirmative : \emph{Let's + infinitif sans \textbf{to}}]
\emph{Let's} est la forme contractée de \emph{Let us}. Elle exprime une \cle{suggestion inclusive} (suggestion qui inclut le locuteur) : « Et si nous... ? », « Allons-y... ».
\begin{center}
\textbf{Structure :} Let's + verbe (base verbale) + complément.
\end{center}
\end{definition}

\begin{exemplebox}[Exemples contextuels (Fêtes communautaires)]
\begin{itemize}
    \item \textbf{Let's organise} a cleaning day before the festival. (Organisons une journée de nettoyage...)
    \item \textbf{Let's invite} the neighbourhood chiefs to the opening ceremony. (Invitons les chefs de quartier...)
    \item \textbf{Let's prepare} traditional dishes together. (Préparons des plats traditionnels ensemble.)
\end{itemize}
\emph{Traduction française :} « Et si nous organisions... », « Organisons... », « Allons-y pour... ».
\end{exemplebox}

\begin{definition}[Forme négative : \emph{Let's not + infinitif sans \textbf{to}}]
Pour exprimer une suggestion de \emph{ne pas faire} quelque chose, on place \textbf{not} après \emph{Let's}, sans \textbf{to}.
\begin{center}
\textbf{Structure :} Let's not + verbe (base verbale) + complément.
\end{center}
\end{definition}

\begin{exemplebox}[Exemples contextuels (Règlement de conflits)]
\begin{itemize}
    \item \textbf{Let's not invite} politics into this celebration. (N'invitons pas la politique dans cette fête.)
    \item \textbf{Let's not argue} about the referee's decision during the match. (Ne contestons pas la décision de l'arbitre...)
    \item \textbf{Let's not forget} the elders' advice. (N'oublions pas les conseils des aînés.)
\end{itemize}
\end{exemplebox}

\begin{definition}[Question-tag : \emph{shall we?}]
Comme \emph{Let's} implique une invitation à l'action commune, le \cle{tag question} (question d'appui) associé est toujours \textbf{shall we?}.
\begin{center}
Let's start the meeting, \textbf{shall we?} / Let's not be late, \textbf{shall we?}
\end{center}
\end{definition}

\begin{propriete}[Différence cruciale : \emph{Let's} vs \emph{Let us}]
Bien que \emph{Let's} = \emph{Let us} (contracté), l'usage et le sens diffèrent radicalement.
\begin{center}
\begin{tabularx}{\linewidth}{|l|X|X|}\hline
\textbf{Critère} & \textbf{\emph{Let's} (Suggestion collective)} & \textbf{\emph{Let us} (Demande d'autorisation)} \\ \hline
\textbf{Sens} & « Faisons-le ensemble » (Nous \textbf{inclus}) & « Permettez-nous de le faire » (Nous \textbf{exclus}) \\ \hline
\textbf{Sujet implicite} & \emph{We} (locuteur + auditeur) & \emph{We} (locuteur + groupe, s'adressant à un tiers autoritaire) \\ \hline
\textbf{Exemple} & \emph{Let's dance!} (Dansons !) & \emph{Let us go home, Sir.} (Laissez-nous rentrer, Monsieur.) \\ \hline
\textbf{Négation} & \emph{Let's not dance.} & \emph{Do not let us down.} (Ne nous décevez pas.) \\ \hline
\end{tabularx}
\end{center}
\emph{Note d'examen :} Au Probatoire/Bac Tech, on teste souvent la transformation \emph{Let us} $\to$ \emph{Let's} (suggestion) ou le choix du tag \emph{shall we?}.
\end{propriete}

\begin{attention}[Piège fréquent : \emph{*Let's not to go}]
\emph{Erreur :} « Let's not \textbf{to} go. » \\
\emph{Correct :} \textbf{Let's not go.} \\
\emph{Règle :} Après \emph{Let's not}, le verbe est à la \textbf{base verbale} (infinitif sans \textbf{to}). Jamais de \textbf{to}. Autre erreur fréquente : écrire \emph{Lets} sans apostrophe — il faut toujours \textbf{Let's} (contraction de \emph{Let us}).
\end{attention}

\begin{exoresolu}[Transformation et tags]
\textbf{Énoncé :} Transformez les phrases suivantes en suggestions avec \emph{Let's} (affirmatif ou négatif) et ajoutez le tag approprié.
\begin{enumerate}
    \item We should organise a football tournament. (Affirmatif)
    \item We must not bring alcohol to the village square. (Négatif)
    \item Let us ask the principal for permission. (Suggestion collective)
\end{enumerate}
\tcblower
\textbf{Solution détaillée :}
\begin{enumerate}
    \item \textbf{Let's organise a football tournament, shall we?} (Base verbale \emph{organise}, tag \emph{shall we?})
    \item \textbf{Let's not bring alcohol to the village square, shall we?} (\emph{not} + base verbale \emph{bring})
    \item \textbf{Let's ask the principal for permission, shall we?} (Contraction de \emph{Let us} en \emph{Let's} car c'est une suggestion \emph{nous inclus} ; tag \emph{shall we?})
\end{enumerate}
\end{exoresolu}

\begin{exercice}[Entraînement : Let's / Let's not]
\textbf{Énoncé :} Complétez avec \emph{Let's} ou \emph{Let's not} + le verbe entre parenthèses (forme correcte). Ajoutez le tag.
\begin{enumerate}
    \item \_\_\_\_\_\_\_\_\_\_\_\_\_\_\_\_\_\_\_\_ (clean) the community hall before the guests arrive, \_\_\_\_\_\_\_\_\_\_?
    \item \_\_\_\_\_\_\_\_\_\_\_\_\_\_\_\_\_\_\_\_ (make) too much noise during the traditional dance, \_\_\_\_\_\_\_\_\_\_?
    \item \_\_\_\_\_\_\_\_\_\_\_\_\_\_\_\_\_\_\_\_ (invite) the neighbouring village to our cultural week, \_\_\_\_\_\_\_\_\_\_?
    \item \_\_\_\_\_\_\_\_\_\_\_\_\_\_\_\_\_\_\_\_ (forget) to thank the sponsors in the closing speech, \_\_\_\_\_\_\_\_\_\_?
\end{enumerate}
\end{exercice}

\begin{corrige}[Exercice n°1]
\begin{enumerate}
    \item \textbf{Let's clean} the community hall before the guests arrive, \textbf{shall we?}
    \item \textbf{Let's not make} too much noise during the traditional dance, \textbf{shall we?}
    \item \textbf{Let's invite} the neighbouring village to our cultural week, \textbf{shall we?}
    \item \textbf{Let's not forget} to thank the sponsors in the closing speech, \textbf{shall we?}
\end{enumerate}
\end{corrige}

\begin{popfigure}
\begin{tikzpicture}[
    box/.style={draw=popblue, thick, rounded corners=4pt, fill=popblue!5, minimum height=1.8cm, text width=3.6cm, align=center, font=\small},
    arrow/.style={-Stealth, thick, popdark},
    labelbox/.style={draw=poporange, thick, rounded corners=4pt, fill=poporange!10, minimum height=1.2cm, text width=3.2cm, align=center, font=\small\bfseries}
]
% Nœuds principaux
\node[box, align=center] (intro) at (0,0){\textbf{INTRODUCTION}\\\cle{Hook} (Proverbe)\\\cle{Définition} du sujet\\\cle{Thèse} (Réponse)\\\cle{Annonce du plan}};
\node[box, align=center] (arg1) at (0,-2.6){\textbf{ARGUMENT 1}\\\emph{Espace neutre de rencontre}\\\cle{Exemple :} Lutte à Garoua\\\cle{Connecteur :} \emph{For instance,}};
\node[box, align=center] (arg2) at (0,-5.2){\textbf{ARGUMENT 2}\\\emph{Règles partagées = Respect}\\\cle{Exemple :} Njangi à Bafoussam\\\cle{Connecteur :} \emph{Furthermore,}};
\node[box, align=center] (arg3) at (0,-7.8){\textbf{ARGUMENT 3}\\\emph{Identité commune renforcée}\\\cle{Exemple :} Festival Ngondo\\\cle{Connecteur :} \emph{Moreover,}};
\node[box, align=center] (concl) at (0,-10.4){\textbf{CONCLUSION}\\\cle{Bilan} des 3 arguments\\\cle{Ouverture} (Perspective)\\\cle{Connecteur :} \emph{Ultimately,}};

% Flèches
\draw[arrow] (intro.south) -- (arg1.north);
\draw[arrow] (arg1.south) -- (arg2.north);
\draw[arrow] (arg2.south) -- (arg3.north);
\draw[arrow] (arg3.south) -- (concl.north);

% Étiquettes latérales (Plan type)
\node[labelbox, anchor=east] at (-4.6,0) {\textbf{1. INTRODUCTION}};
\node[labelbox, anchor=east, align=center] at (-4.6,-5.2){\textbf{2. DÉVELOPPEMENT}\\(3 arguments)};
\node[labelbox, anchor=east] at (-4.6,-10.4) {\textbf{3. CONCLUSION}};

% Connecteurs transversaux
\node[font=\small\itshape, text=poppurple, anchor=west] at (3.2,-1.3) {\emph{First,...}};
\node[font=\small\itshape, text=poppurple, anchor=west] at (3.2,-3.9) {\emph{Second,...}};
\node[font=\small\itshape, text=poppurple, anchor=west] at (3.2,-6.5) {\emph{Finally,...}};
\node[font=\small\itshape, text=poppurple, anchor=west] at (3.2,-9.1) {\emph{In conclusion,...}};
\end{tikzpicture}
\legende{Squelette de la dissertation argumentée : structure en 5 paragraphes avec connecteurs obligatoires et exemples camerounais intégrés.}
\end{popfigure}

\courssub{Écriture argumentée : « Les activités récréatives apaisent les conflits communautaires »}

\begin{methode}[Rédiger l'introduction (4 étapes obligatoires)]
Pour réussir l'introduction au Bac Tech, suivez ce protocole :
\begin{enumerate}
    \item \textbf{Accroche (Hook) :} Un proverbe, une statistique, une citation ou une observation générale. \emph{Ex: "In Cameroon, a popular saying goes: 'One hand washes the other'."}
    \item \textbf{Définition du sujet :} Définir \cle{recreational activities} (sport, danse, jeux, travaux communautaires) et \cle{community disputes} (conflits fonciers, générationnels, politiques).
    \item \textbf{Thèse (Position claire) :} Répondre OUI/NUANCÉ à la question. \emph{Ex: "I firmly believe that recreational activities are powerful tools to settle disputes and restore peace."}
    \item \textbf{Annonce du plan :} Utiliser \textbf{First, Second, Finally} (ou \emph{Firstly, Furthermore, Ultimately}). \emph{Ex: "First, they create neutral ground. Second, they teach shared rules. Finally, they build a common identity."}
\end{enumerate}
\end{methode}

\begin{aretenir}[Connecteurs argumentatifs (Logique du discours)]
\begin{center}
\begin{tabularx}{\linewidth}{|l|X|X|}\hline
\textbf{Fonction} & \textbf{Connecteurs (Mots-liens)} & \textbf{Exemple d'usage} \\ \hline
\cle{Ajout / Renforcement} & \textbf{Furthermore, Moreover, In addition, Besides} & \emph{Furthermore,} the njangi builds trust. \\ \hline
\cle{Illustration / Exemple} & \textbf{For instance, For example, Such as} & \emph{For instance,} wrestling in Garoua. \\ \hline
\cle{Explication / Cause} & \textbf{In fact, Indeed, Because} & \emph{In fact,} rules require respect. \\ \hline
\cle{Opposition / Nuance} & \textbf{However, Nevertheless, On the other hand} & \emph{However,} violence can occur if rules fail. \\ \hline
\cle{Conséquence / Résultat} & \textbf{Consequently, Therefore, As a result} & \emph{Consequently,} peace is restored. \\ \hline
\cle{Conclusion / Synthèse} & \textbf{Ultimately, In conclusion, To sum up} & \emph{Ultimately,} unity prevails. \\ \hline
\end{tabularx}
\end{center}
\emph{Consigne :} Utilisez au moins \textbf{5 connecteurs différents} dans votre développement pour viser la note maximale en \cle{organisation}.
\end{aretenir}

\begin{exemplebox}[Modèle de développement (Argument 1 \& 2) — Extrait]
\textbf{First,} recreational activities provide a \cle{neutral meeting ground} (espace neutre de rencontre). \textbf{For instance,} in Garoua, traditional wrestling (\emph{la lutte}) brings together young men from rival quarters. On the wrestling arena, social status and family grudges disappear; only technique and fair play matter. \textbf{In fact,} the referee (often an elder) is respected by all. \textbf{Furthermore,} the \emph{njangi} (tontine/association rotative) in Bafoussam acts as a conflict prevention mechanism. Members meet weekly, share a meal, and discuss community issues before they escalate. \textbf{Consequently,} disputes over land or money are often settled amicably during these gatherings rather than in courts.
\end{exemplebox}

\begin{savaistu}[Sagesse locale : \emph{One hand washes the other}]
Le proverbe ouest-camerounais \textbf{« One hand washes the other »} (litt. « Une main lave l'autre ») illustre parfaitement l'entraide récréative. Il signifie : « On a besoin les uns des autres pour réussir ». Il est souvent cité dans les discours d'ouverture des tournois inter-quartiers (football, lutte, danse) pour rappeler que la compétition sportive est une métaphore de la solidarité communautaire. Intégrer ce proverbe en \cle{accroche} ou en \cle{ouverture} de conclusion valorise votre copie (ancrage local apprécié au Bac Tech).
\end{savaistu}

\begin{exemplebox}[Exemples camerounais à mobiliser dans vos copies]
\begin{center}
\begin{tabularx}{\linewidth}{|l|X|}\hline
\textbf{Exemple} & \textbf{Argument qu'il illustre} \\ \hline
Lutte traditionnelle (Nord) & Espace de confrontation ritualisée, respect de l'adversaire, rôle de l'arbitre-aîné. \\ \hline
Njangi / Tontines (Ouest) & Réunions régulières, médiation par les aînés, entraide financière. \\ \hline
Festival Ngondo (Sawa, Douala) & Cohésion d'une communauté autour de rites partagés (course de pirogues, danses). \\ \hline
Danses traditionnelles & Cohésion intergénérationnelle, transmission des valeurs, fête partagée. \\ \hline
Jeux modernes (Football, Ludo) & Mixité sociale, règles universelles, esprit d'équipe. \\ \hline
\end{tabularx}
\end{center}
\emph{Conseil :} Variez les exemples (ne citez pas uniquement la lutte) et reliez chaque exemple à un connecteur : les copies qui diversifient les illustrations obtiennent les meilleures notes en \cle{contenu} et \cle{originalité}.
\end{exemplebox}

\begin{exoresolu}[Rédaction complète : Introduction + Plan détaillé]
\textbf{Sujet :} \emph{« Recreational activities can settle community disputes and bring peace. » Discuss. (250--300 words)}

\textbf{Introduction (Modèle) :}
\begin{quote}
\emph{« One hand washes the other »,} says a Cameroonian proverb. In Cameroon, community life is often shaken by disputes over land, inheritance or political differences. \textbf{Recreational activities} — such as traditional wrestling, \emph{njangi} meetings, cultural festivals or football tournaments — are organised leisure events where people meet freely. \textbf{I strongly believe that these activities are effective instruments to settle conflicts and restore lasting peace.} \textbf{First,} they create a neutral space for dialogue. \textbf{Second,} they impose shared rules that teach mutual respect. \textbf{Finally,} they strengthen a common identity beyond differences.
\end{quote}

\textbf{Plan du développement (Squelette pour l'élève) :}
\begin{itemize}
    \item \textbf{§2 (Arg 1 - Neutral ground) :} Lutte à Garoua / Tournois inter-quartiers. Connecteurs : \emph{First, For instance, In fact.}
    \item \textbf{§3 (Arg 2 - Shared rules) :} Njangi à Bafoussam / Règles du jeu = lois sociales. Connecteurs : \emph{Furthermore, Moreover, Consequently.}
    \item \textbf{§4 (Arg 3 - Common identity) :} Festival Ngondo (Sawa) / Danses / Chants. Connecteurs : \emph{Additionally, Ultimately.}
    \item \textbf{§5 (Contre-argument / Nuance - Optionnel pour excellents) :} \emph{However,} bad organisation can spark violence. \emph{Nevertheless,} with good leadership...
    \item \textbf{Conclusion :} Bilan + \emph{Let's encourage...} + Ouverture.
\end{itemize}
\end{exoresolu}

\courssub{Intégration : Auto-évaluation \& Production finale}

\begin{methode}[Grille d'auto-évaluation (10 critères)]
Avant de rendre votre copie, vérifiez chaque point de la grille ci-dessous. Objectif : \textbf{10/10} pour viser une excellente note. Cochez chaque critère validé sur votre brouillon.
\end{methode}

\begin{popfigure}
\begin{tikzpicture}[
    check/.style={draw=popgreen, thick, minimum size=0.55cm},
    item/.style={font=\small, anchor=west, text width=12.5cm},
    title/.style={font=\bfseries\large, text=popdark},
    cat/.style={font=\bfseries\small, text=popblue}
]
% Titre
\node[title] at (0, 6.5) {CHECKLIST AUTO-ÉVALUATION — DISSERTATION ARGUMENTÉE};

% Catégorie CONTENU
\node[cat, anchor=west] at (-7.5, 5.5) {\textbf{CONTENU}};
\node[check] (c1) at (-7.5, 4.8) {}; \node[item] at (-7.2, 4.8) {Thèse claire, nuancée et répondant exactement au sujet.};
\node[check] (c2) at (-7.5, 4.1) {}; \node[item] at (-7.2, 4.1) {3 arguments distincts + 3 exemples camerounais précis (Garoua, Bafoussam, Ngondo, etc.).};
\node[check] (c3) at (-7.5, 3.4) {}; \node[item] at (-7.2, 3.4) {Vocabulaire spécifique utilisé : \emph{neutral ground, shared rules, common identity, mediate, escalate, fair play}.};
\node[check] (c4) at (-7.5, 2.7) {}; \node[item] at (-7.2, 2.7) {Contre-argument (nuance) intégré si niveau « Excellent » (\emph{However, bad organisation...}).};

% Catégorie ORGANISATION
\node[cat, anchor=west] at (-7.5, 1.7) {\textbf{ORGANISATION}};
\node[check] (o1) at (-7.5, 1.0) {}; \node[item] at (-7.2, 1.0) {Introduction complète : Accroche + Définition + Thèse + Annonce du plan (\emph{First, Second, Finally}).};
\node[check] (o2) at (-7.5, 0.3) {}; \node[item] at (-7.2, 0.3) {Paragraphes distincts (1 idée = 1 paragraphe), connecteurs variés (\emph{Furthermore, For instance, Consequently, Ultimately}).};
\node[check] (o3) at (-7.5, -0.4) {}; \node[item] at (-7.2, -0.4) {Conclusion : Bilan synthétique + Ouverture avec \emph{Let's} / \emph{Let's not}.};

% Catégorie LANGUE
\node[cat, anchor=west] at (-7.5, -1.4) {\textbf{LANGUE}};
\node[check] (l1) at (-7.5, -2.1) {}; \node[item] at (-7.2, -2.1) {Grammaire : \emph{Present Perfect} (\emph{have settled}), \emph{Past Perfect} (\emph{had disputed}), \emph{Let's/Let's not} correct (sans \emph{to}).};
\node[check] (l2) at (-7.5, -2.8) {}; \node[item] at (-7.2, -2.8) {Orthographe / Ponctuation soignées (majuscules, points, virgules avant connecteurs).};
\node[check] (l3) at (-7.5, -3.5) {}; \node[item] at (-7.2, -3.5) {Nombre de mots : 250--300 mots (tolérance $\pm$ 10\%).};

% Légende cases
\node[font=\small, anchor=west] at (-7.5, -4.3) {Case cochée = Validé \quad Case vide = À revoir \quad Objectif : 10 cases cochées avant remise.};
\end{tikzpicture}
\legende{Grille d'auto-évaluation visuelle (10 critères). Imprimez-la ou recopiez-la sur votre brouillon.}
\end{popfigure}

\begin{exercice}[Production écrite guidée — Sujet d'examen type]
\textbf{Sujet :} \emph{« As the president of the Students' Cultural Club, write an argumentative essay (250--300 words) to convince the school administration to organise an Intercultural Festival. Use the thesis: "An intercultural festival strengthens unity and prevents violence in school." Include: Let's / Let's not, Present Perfect, Past Perfect, 5+ connectors, 3 local examples. »}

\textbf{Consignes :}
\begin{enumerate}
    \item Rédigez l'introduction complète (4 étapes) sur votre brouillon.
    \item Construisez le développement (3 arguments + exemples) en utilisant la checklist.
    \item Écrivez la conclusion avec une ouverture contenant \emph{Let's} ou \emph{Let's not}.
    \item Relisez-vous avec la grille d'auto-évaluation (Figure ci-dessus).
    \item Écrivez la version finale au propre.
\end{enumerate}
\end{exercice}

\begin{corrige}[Exercice n°2 — Proposition de corrigé type (environ 280 mots)]
\textbf{An Intercultural Festival: A Bridge to Peace in Our School}

\emph{« One hand washes the other »,} says a Cameroonian proverb. In our technical high school, tensions sometimes rise between students from different regions or specialities. An \textbf{intercultural festival} is a school event where students showcase their cultures through dance, food, wrestling and fashion. \textbf{I firmly believe that such a festival strengthens unity and prevents violence.} \textbf{First,} it creates a neutral meeting ground. \textbf{For instance,} during the traditional wrestling demonstrations, students from the North and the West compete respectfully under the same rules. \textbf{In fact,} they \textbf{have shared} the same arena for years without incident. \textbf{Furthermore,} the festival teaches shared rules. \textbf{For example,} the \emph{njangi}-style cooking competition requires teams to follow strict hygiene and timing rules. \textbf{Consequently,} students learn that discipline benefits the whole group. \textbf{Moreover,} it builds a common identity. \textbf{Last year,} the Ngondo-inspired canoe race \textbf{had brought} together Francophones and Anglophones before the closing ceremony even started. \textbf{Ultimately,} we realise we are one student body.

\textbf{However,} some fear that cultural pride might turn into rivalry. \textbf{Nevertheless,} with teacher supervision, this risk is low. \textbf{In conclusion,} the festival transforms differences into richness. \textbf{Let's organise} it annually, \textbf{shall we?} \textbf{Let's not allow} prejudice to divide us. Together, we build a peaceful campus.
\end{corrige}

\begin{aretenir}[Résumé des points clés pour le Bac Tech — Unité 1]
\begin{itemize}
    \item \textbf{Let's + BV} = Suggestion inclusive (\emph{Let's go!}) | Tag = \textbf{shall we?}
    \item \textbf{Let's not + BV} = Suggestion négative (\emph{Let's not fight!}) | \textbf{Pas de \emph{to}}.
    \item \textbf{Let us} (non contracté) = Demande d'autorisation à un supérieur (\emph{Let us leave, Sir.}).
    \item \textbf{Essai argumenté :} Intro (Hook, Définition, Thèse, Plan) $\to$ 3 § Développement (Arguments + Exemples camerounais + Connecteurs) $\to$ Conclusion (Bilan + Ouverture + \emph{Let's}).
    \item \textbf{Connecteurs stars :} \emph{Furthermore, For instance, Consequently, However, Ultimately}.
    \item \textbf{Grammaire intégrée :} \emph{Present Perfect} (expérience/résultat), \emph{Past Perfect} (antériorité), \emph{Let's/Let's not}.
\end{itemize}
\end{aretenir}

\coursec{Intégration : organisation d'un festival interculturel scolaire}

\courssub{Transition depuis la section 6}
After practising \cle{Let's / Let's not} for suggestions and writing argumentative paragraphs about how recreational activities foster peace, you are now ready to mobilise \emph{all} the language tools of this unit — vocabulary on festivals and community life, modals for permission/obligation/deduction, purpose/result connectors, present perfect and past perfect — in a single, realistic project: the \textbf{Intercultural School Festival}. This integration task mirrors the way technical professionals combine diverse competences to deliver a concrete outcome.

\courssub{Présentation de la tâche complexe}
\begin{definition}[Tâche d'intégration — Festival Interculturel du Lycée]
En groupes de 6 élèves, vous organisez un festival de 3 jours célébrant la diversité culturelle de l'établissement. La tâche exige la production de quatre livrables évalués :
\begin{enumerate}
\item Une \textbf{lettre formelle au Proviseur} (demande d'autorisation + budget prévisionnel).
\item Une \textbf{affiche bilingue} (français / anglais) avec date, lieu, activités, modalités d'inscription.
\item Un \textbf{scénario de médiation} d'un conflit fictif (ex. : bruit des répétitions vs besoin de repos des internes).
\item Un \textbf{discours d'ouverture} (2 min) intégrant \emph{present perfect}, modaux (\cle{must}, \cle{should}, \cle{ought to}), \cle{Let's}, connecteurs de but/résultat.
\end{enumerate}
\end{definition}

\courssub{Répartition des rôles au sein du groupe}
\begin{center}
\begin{tabularx}{\linewidth}{|l|X|}\hline
\textbf{Rôle} & \textbf{Responsabilités principales} \\ \hline
Coordonnateur & Pilotage global, respect du calendrier, liaison avec l'enseignant \\ \hline
Secrétaire & Rédaction de la lettre au Proviseur, compte-rendus de réunions \\ \hline
Trésorier & Budget prévisionnel (tableur), recherche de sponsoring, suivi dépenses \\ \hline
Communicant & Conception affiche, publications réseaux sociaux, annonces micro \\ \hline
Médiateur & Écriture du scénario de conflit + résolution, animation atelier médiation \\ \hline
Orateur & Rédaction et répétition du discours d'ouverture \\ \hline
\end{tabularx}
\end{center}

\begin{methode}[Travail de groupe efficace : 5 étapes clés]
\begin{enumerate}
\item \textbf{Définir l'objectif commun} : « Festival réussi, inclusif, sans incident, budget équilibré. »
\item \textbf{Répartir selon les forces} : à l'aise à l'oral $\to$ Orateur/Médiateur ; rigoureux chiffres $\to$ Trésorier ; créatif visuel $\to$ Communicant.
\item \textbf{Fixer des échéances intermédiaires} (jalons) : J+2 choix communautés, J+5 budget v1, J+8 lettre signée, J+10 affiche imprimée, J+12 scénario validé, J+13 discours prêt, J+14 répétition générale.
\item \textbf{Points de synchronisation quotidiens (5 min)} : tour de table « Qu'ai-je fait ? Qu'est-ce qui bloque ? Besoin d'aide ? ».
\item \textbf{Rédiger collectivement sur espace partagé} (Google Docs, Framapad, ou dossier partagé sur clé USB / disque local) : historique des versions, commentaires, traçabilité.
\end{enumerate}
\end{methode}

\courssub{Planification globale : diagramme de Gantt simplifié (14 séances)}
\begin{popfigure}
\begin{tikzpicture}[x=0.6cm, y=0.8cm, font=\small]
% Grille
\draw[popline, dashed] (0,0) grid (14,7);
% Numéros de séance
\foreach \x in {1,...,14} \node[below] at (\x-0.5, -0.3) {\x};
% Étiquettes des thèmes (axe Y)
\node[left] at (-0.5, 6.5) {Vocab/Interaction};
\node[left] at (-0.5, 5.5) {Modaux};
\node[left] at (-0.5, 4.5) {Écoute/Conflits};
\node[left] at (-0.5, 3.5) {But/Résultat};
\node[left] at (-0.5, 2.5) {Lecture/Perfect};
\node[left] at (-0.5, 1.5) {Let's/Écriture};
\node[left] at (-0.5, 0.5) {Intégration/Éval.};
% Barres horizontales (début inclusif, fin exclusive)
\fill[popblue!70] (0,6) rectangle (3,7);       % Séances 1-3
\fill[poporange!70] (3,5) rectangle (5,6);     % Séances 4-5
\fill[popgreen!70] (5,4) rectangle (7,5);      % Séances 6-7
\fill[poppurple!70] (7,3) rectangle (8,4);     % Séance 8
\fill[poppink!70] (8,2) rectangle (10,3);      % Séances 9-10
\fill[popgold!70] (10,1) rectangle (12,2);     % Séances 11-12
\fill[popdark!70] (12,0) rectangle (14,1);     % Séances 13-14
% Légende
\begin{scope}[yshift=-1.8cm]
\foreach \c/\l in {popblue!70/Vocab/Interaction, poporange!70/Modaux, popgreen!70/Écoute/Conflits, poppurple!70/But/Résultat, poppink!70/Lecture/Perfect, popgold!70/Let's/Écriture, popdark!70/Intégration/Éval.} {
    \fill[\c] (0,0) rectangle (0.5,0.3);
    \node[right] at (0.6,0.15) {\l};
    \pgftransformyshift{-0.4cm}
}
\end{scope}
\end{tikzpicture}
\legende{Diagramme de Gantt simplifié : répartition des 14 séances par bloc thématique. Les barres horizontales indiquent les séances consacrées à chaque module ; la phase d'intégration (S13-14) mobilise l'ensemble des acquis.}
\end{popfigure}

\courssub{Produit 1 — Lettre formelle au Proviseur}
\begin{definition}[Structure d'une lettre formelle administrative (camérounais)]
Une lettre au chef d'établissement suit un canevas codifié :
\begin{enumerate}
\item \textbf{En-tête expéditeur} : Nom du lycée, BP, Ville, Téléphone, Email.
\item \textbf{Date} (ex. : Yaoundé, le 15 octobre 2026).
\item \textbf{Destinataire} : Monsieur le Proviseur, Lycée Technique de [Ville].
\item \textbf{Objet} : clair, concis (ex. : \emph{Demande d'autorisation pour l'organisation du Festival Interculturel du Lycée}).
\item \textbf{Formule d'appel} : \emph{Monsieur le Proviseur,}
\item \textbf{Corps} : (a) Contexte (classe, unité pédagogique), (b) Description du projet (3 jours, 3 communautés, activités), (c) Budget prévisionnel synthétique, (d) Engagements (sécurité, propreté, encadrement).
\item \textbf{Formule de politesse} : \emph{Je vous prie d'agréer, Monsieur le Proviseur, l'expression de ma considération distinguée.}
\item \textbf{Signature} : Nom, Prénom, Classe, Fonction dans le groupe (ex. : Secrétaire du comité d'organisation).
\end{enumerate}
\end{definition}

\begin{attention}[Piège de registre — Ne pas mélanger familier et formel]
\textbf{Erreur fréquente} : « Hey Principal, can we do a festival? It'll be cool! » \\
\textbf{Correction} : Maintenez le registre formel tout au long. Utilisez \cle{may} pour la permission (\emph{We kindly request that we may organise...}), \cle{would like to} pour la courtoisie, \cle{undertake to} pour l'engagement. Pas de contractions (\emph{we're}, \emph{it's}) ; préférez \emph{we are}, \emph{it is}.
\end{attention}

\begin{popfigure}
\begin{tikzpicture}[node distance=4mm, font=\small]
% En-tête
\node[anchor=west, align=center] (header) at (0,10){\textbf{LYCÉE TECHNIQUE DE YAOUNDÉ}\\\textsc{bp 12345 — YAOUNDÉ — TÉL : (237) 2 22 33 44 55 — EMAIL : contact@ltyaounde.cm}};
\draw[popline, thick] (0,9.5) -- (\linewidth,9.5);
% Date / Destinataire
\node[anchor=west] (date) at (0,8.8) {Yaoundé, le 15 octobre 2026};
\node[anchor=west, align=center] (dest) at (0,8.1){\textbf{À l'attention de Monsieur le Proviseur}\\\textsc{Lycée Technique de Yaoundé}};
% Objet
\node[anchor=west] (objet) at (0,7.0) {\textbf{Objet : Demande d'autorisation pour l'organisation du Festival Interculturel du Lycée (3 jours)}};
% Appel
\node[anchor=west] (appel) at (0,6.3) {Monsieur le Proviseur,};
% Corps
\node[anchor=west, text width=0.92\linewidth, align=justify] (corps1) at (0,5.8) {Dans le cadre de l'unité pédagogique \emph{« Family and Social Life »} (Première Technique, Anglais), notre classe souhaite organiser un \textbf{Festival Interculturel} sur trois journées (du 12 au 14 novembre 2026) afin de valoriser la diversité culturelle de notre établissement et de mettre en pratique les compétences langagières acquises (interactions sociales, médiation, expression de but et d'obligation).};
\node[anchor=west, text width=0.92\linewidth, align=justify] (corps2) at (0,4.6) {Trois communautés seront représentées : \textbf{Bamiléké} (danse \emph{Njeng}, mets \emph{Nkui}), \textbf{Bassa} (contes \emph{Mvet}, plat \emph{Kpem}), \textbf{Peulh} (chants pastoraux, \emph{Lait caillé}). Les activités prévues : tournois sportifs (football, handball), ateliers culinaires, spectacles de danse/contes, stands d'artisanat.};
\node[anchor=west, text width=0.92\linewidth, align=justify] (corps3) at (0,3.4) {Le budget prévisionnel s'élève à \textbf{150 000 FCFA} (détail ci-joint). Les recettes attendues : cotisations élèves (100 000 FCFA) + sponsoring local (50 000 FCFA). Nous \textbf{nous engageons à} : (1) respecter le règlement intérieur, (2) assurer l'encadrement par 4 enseignants volontaires, (3) laisser les lieux propres, (4) remettre un rapport d'activités sous 8 jours.};
% Politesse
\node[anchor=west, text width=0.92\linewidth, align=justify] (politesse) at (0,2.0) {En espérant une suite favorable, je vous prie d'agréer, Monsieur le Proviseur, l'expression de ma considération distinguée.};
% Signature
\node[anchor=west, align=center] (sig) at (0,1.0){Fait à Yaoundé, le 15 octobre 2026\\[2mm] \textbf{Ngo Bitjek Marie}\\\textit{Élève de 1re TMA — Secrétaire du comité d'organisation}};
\end{tikzpicture}
\legende{Modèle de lettre formelle au Proviseur : en-tête institutionnel, date, destinataire, objet, formule d'appel, corps structuré (contexte, description, budget, engagements), formule de politesse, signature avec qualité.}
\end{popfigure}

\begin{exoresolu}[Rédaction de la lettre — Mise en application]
\textbf{Consigne} : À partir du modèle ci-dessus, rédigez la lettre pour votre propre établissement en adaptant les communautés représentées, les dates et le budget. Vérifiez : (1) absence de contractions, (2) modaux de permission (\emph{may}, \emph{could}), (3) connecteurs de but (\emph{in order to}, \emph{so that}), (4) \emph{present perfect} pour les actions déjà menées (\emph{We have contacted three local sponsors...}).

\textbf{Extrait corrigé (corps de lettre)} :
\begin{quote}
\emph{We \textbf{have selected} three cultural communities represented in our school: the Bamiléké, the Bassa and the Fulani. Each community \textbf{will} present traditional dances, dishes and tales. We \textbf{would like to} use the school courtyard and the multipurpose hall. We \textbf{have already contacted} two local sponsors who \textbf{have agreed to} support us. We \textbf{undertake to} ensure strict supervision and cleanliness. We \textbf{kindly request that you may} grant us authorisation.}
\end{quote}
\textbf{Analyse grammaticale} : \emph{have selected / have contacted / have agreed} = \cle{present perfect} (actions passées récentes, lien avec le présent) ; \emph{will present} = futur simple (programme) ; \emph{would like to} = courtoisie ; \emph{may grant} = permission formelle ; \emph{undertake to} = engagement formel.
\end{exoresolu}

\courssub{Produit 2 — Affiche bilingue (français / anglais)}
\begin{exemplebox}[Affiche type — Contenu minimal]
\begin{center}
\begin{tabularx}{\linewidth}{|X|X|}\hline
\textbf{FRANÇAIS} & \textbf{ENGLISH} \\ \hline
\textsc{festival interculturel du lycée} & \textsc{high school intercultural festival} \\
12 -- 14 novembre 2026 & November 12 -- 14, 2026 \\
Cour d'honneur \& Salle polyvalente & School courtyard \& Multipurpose hall \\
\textbf{Communautés invitées :} Bamiléké, Bassa, Peulh & \textbf{Featured communities:} Bamiléké, Bassa, Fulani \\
\textbf{Programme :} & \textbf{Programme:} \\
• Tournois sportifs (foot, hand) & • Sports tournaments (football, handball) \\
• Ateliers cuisine : Nkui, Kpem, Lait caillé & • Cooking workshops: Nkui, Kpem, Curdled milk \\
• Spectacles danse/contes (18h--21h) & • Dance/tale shows (6--9 pm) \\
• Stands artisanat & • Craft stands \\
\textbf{Entrée libre — Inscription stands :} & \textbf{Free entry — Stand registration:} \\
M. Ngono (Salle 12) avant le 05/11 & Mr Ngono (Room 12) before Nov 5 \\
\textit{Organisé par la 1re TMA — Anglais} & \textit{Organised by 1re TMA — English} \\ \hline
\end{tabularx}
\end{center}
\end{exemplebox}

\courssub{Produit 3 — Scénario de médiation (conflit fictif : bruit vs repos)}
\begin{definition}[Structure d'un scénario de médiation scolaire]
\begin{enumerate}
\item \textbf{Contexte} : Lieu, heure, protagonistes (ex. : répétition tambours Bamiléké 20h, dortoir internes à 15 m).
\item \textbf{Positions initiales} : Groupe A (« Nous \emph{must} répéter pour le festival » / « We \emph{have been practising} for weeks ») ; Groupe B (« Nous \emph{need} dormir, examens demain » / « The noise \emph{has prevented} us from sleeping »).
\item \textbf{Intervention du médiateur} : Reformulation, recherche d'intérêts communs, proposition de solutions.
\item \textbf{Accord} : Horaires aménagés (répétitions 16h--18h, silence après 20h), local insonorisé si possible, engagement écrit.
\end{enumerate}
\end{definition}

\begin{exoresolu}[Extrait de scénario — Médiation]
\textbf{Médiateur} : « Good evening. I \textbf{understand} the drum group \textbf{wants to} perfect their performance. \textbf{At the same time}, the boarders \textbf{need} to rest for their exams. \textbf{Let's find} a solution together. \textbf{What if} you rehearsed from 4 pm to 6 pm in the music room, and we \textbf{ensure} quiet after 8 pm? »
\textbf{Chef tambours} : « We \textbf{have been practising} since Monday. We \textbf{could} move to the music room. »
\textbf{Délégué internes} : « That \textbf{would} work. We \textbf{ought to} sign a paper. »
\textbf{Médiateur} : « \textbf{Let's write} the agreement now. \textbf{We have reached} a compromise. »

\textbf{Points grammaticaux intégrés} : \emph{have been practising} (present perfect continuous — durée), \emph{have reached} (present perfect — résultat), \emph{must/wants to/need to} (obligation/besoin), \emph{could/would} (possibilité/hypothétique), \emph{Let's} (suggestion), \emph{At the same time} (concession).
\end{exoresolu}

\courssub{Produit 4 — Discours d'ouverture (2 minutes)}
\begin{methode}[Check-list du discours d'ouverture — 2 min $\approx$ 260 mots]
\begin{enumerate}
\item \textbf{Accueil} (20 s) : \emph{Distinguished Principal, dear teachers, fellow students, honored guests...}
\item \textbf{Présent perfect — bilan} (40 s) : \emph{We \textbf{have prepared} this festival for six weeks. Three communities \textbf{have shared} their traditions. We \textbf{have learned} to listen to each other.}
\item \textbf{Modaux — valeurs} (40 s) : \emph{We \textbf{must} respect every culture. We \textbf{should} see diversity as a strength. We \textbf{ought to} build bridges, not walls.}
\item \textbf{But / Résultat} (30 s) : \emph{This festival \textbf{aims at} promoting unity. \textbf{So that} everyone feels at home, \textbf{let's} celebrate together.}
\item \cle{Let's / Let's not} + appel à l'action (30 s) : \emph{\textbf{Let's} open our hearts. \textbf{Let's not} forget this spirit after the festival. \textbf{Let's} make our school a model of harmony.}
\item \textbf{Conclusion formelle} (20 s) : \emph{Thank you. I now declare the Intercultural Festival open!}
\end{enumerate}
\end{methode}

\begin{exemplebox}[Discours type (≈ 250 mots)]
\textit{Distinguished Principal, dear teachers, fellow students, honored guests, good morning.}

\textit{We \textbf{have worked} hard for six weeks. We \textbf{have contacted} elders, \textbf{have learnt} traditional songs, \textbf{have cooked} together. Today, the Bamiléké, the Bassa and the Fulani communities \textbf{are} here, side by side.}

\textit{This festival \textbf{aims at} strengthening our unity. \textbf{So that} no voice goes unheard, \textbf{let's} listen to every drum, every tale, every flavour. We \textbf{must} remember that culture is not a barrier but a bridge. We \textbf{should} be proud of our roots and curious about others'. We \textbf{ought to} turn difference into dialogue.}

\textit{\textbf{Let's} dance, taste, play, talk. \textbf{Let's not} allow prejudice to spoil this moment. \textbf{Let's} make these three days a memory of peace.}

\textit{On behalf of the organising committee, I \textbf{declare} the Intercultural Festival officially open. Thank you.}
\end{exemplebox}

\courssub{Budget type — Exemple chiffré}
\begin{exemplebox}[Budget prévisionnel festival 3 jours (FCFA)]
\begin{center}
\begin{tabular}{|l|r|r|}\hline
\textbf{Poste} & \textbf{Prévision} & \textbf{Réel (après festival)} \\ \hline
Location sonorisation & 50 000 & 48 000 \\
Décoration (tissus, banderoles) & 30 000 & 32 500 \\
Impression affiches (100 ex.) & 15 000 & 14 000 \\
Prix (3 $\times$ 10 000) & 30 000 & 30 000 \\
Imprévus (5\%) & 25 000 & 18 000 \\ \hline
\textbf{TOTAL DÉPENSES} & \textbf{150 000} & \textbf{142 500} \\ \hline
Cotisations élèves (500 $\times$ 200) & 100 000 & 102 500 \\
Sponsoring local (2 commerçants) & 50 000 & 50 000 \\ \hline
\textbf{TOTAL RECETTES} & \textbf{150 000} & \textbf{152 500} \\ \hline
\textbf{EXCÉDENT} & 0 & \textbf{+10 000} \\ \hline
\end{tabular}
\end{center}
\textit{Note : L'excédent sera reversé au fonds social de l'établissement.}
\end{exemplebox}

\courssub{Grille d'évaluation par compétences (sur 5)}
\begin{center}
\begin{tabularx}{\linewidth}{|l|X|X|X|X|X|}\hline
\textbf{Critère} & \textbf{1 — Insuffisant} & \textbf{2 — Fragile} & \textbf{3 — Satisfaisant} & \textbf{4 — Bon} & \textbf{5 — Excellent} \\ \hline
Interaction orale & Quasi absent, pas de tour de parole & Phrases courtes, hésitations fréquentes & Échanges cohérents, vocabulaire adéquat & Fluide, arguments structurés, registre adapté & Naturel, nuancé, leadership bienveillant \\ \hline
Production écrite & Textes incomplets, nombreux anglicismes & Structure visible, erreurs récurrentes & Lettres/affiche conformes, peu d'erreurs & Registre maîtrisé, connecteurs variés, style soigné & Élégant, persuasif, zéro faute majeure \\ \hline
Grammaire en contexte & Modaux/perfects absents ou faux & Quelques formes justes, confusion fréquente & Modaux/perfects/Let's utilisés à bon escient & Variété de structures, auto-correction & Maîtrise totale, effets de sens recherchés \\ \hline
Culture / Interculturel & Aucune référence culturelle & 1 communauté, stéréotypes & 3 communautés, informations exactes & Approfondissement (symboles, valeurs), respect & Analyse comparative, ouverture critique \\ \hline
\end{tabularx}
\end{center}

\courssub{Remédiation ciblée (ateliers 20 min)}
\begin{aretenir}[Ateliers « clinique » selon besoins détectés]
\begin{itemize}
\item \textbf{Atelier Modaux} : tri de phrases (permission / possibilité / obligation / déduction) $\to$ reformulation contextuelle.
\item \textbf{Atelier Temps parfaits} : ligne du temps visuelle (past simple vs present perfect vs past perfect) $\to$ exercice de transformation.
\item \textbf{Atelier Connecteurs but/résultat} : \emph{to / in order to / so that / so as to / for} $\to$ complétion de phrases à trous, puis production libre.
\item \textbf{Atelier Registre formel} : paires familier/formel $\to$ réécriture de SMS en lettre administrative.
\end{itemize}
\end{aretenir}

\courssub{Présentation finale — Simulation du festival en classe}
\begin{experience}[Simulation « Jour J » — Organisation pratique]
\textbf{Matériel} : Tables pour stands (gastronomie, artisanat), sono portable, projecteur pour affiche, chaises pour public, grille d'évaluation imprimée (jurés).
\textbf{Protocole} :
\begin{enumerate}
\item Installation stands (15 min) : chaque communauté décoré son espace.
\item Accueil du jury (Proviseur adjoint, 2 délégués élèves, 1 parent d'élève, 1 enseignant partenaire).
\item Discours d'ouverture par l'Orateur (2 min) — évalué sur grille.
\item Rotation du public : dégustation, quiz culture, mini-match handball (cour de récréation).
\item Scène ouverte : danses, contes (15 min par communauté).
\item Scénario de médiation joué par le Médiateur + 2 volontaires (5 min).
\item Clôture : remise attestations « Ambassadeur Interculturel », photo de groupe.
\end{enumerate}
\textbf{Observations attendues} : Qualité de l'anglais oral (prononciation, grammaire en situation), cohérence des écrits affichés, esprit d'équipe, respect des horaires.
\textbf{Interprétation} : Le jury complète la grille ; l'enseignant donne un feedback oral immédiat (points forts / axes de progression) et une note finale sur 20 (moyenne des 4 critères $\times$ 5 / 5 = note / 20).
\end{experience}

\courssub{Synthèse de l'unité — Transfert et intégration}
\begin{propriete}[L'intégration comme transfert]
L'intégration n'est pas une simple addition de savoirs (vocabulaire + grammaire + fonctions), mais leur \textbf{mobilisation coordonnée} pour résoudre un problème complexe nouveau : organiser un événement réel, avec contraintes budgétaires, administratives, relationnelles. C'est la preuve que l'élève est capable de \textbf{transférer} ses acquis hors du cadre scolaire strict.
\end{propriete}

\begin{savaistu}[Le MINESEC et les Semaines Culturelles]
Le Ministère des Enseignements Secondaires (MINESEC) institutionnalise les \textbf{Semaines Culturelles} annuelles dans tous les lycées. Les meilleures prestations (danse, théâtre, chant, art oratoire) sont sélectionnées pour les \textbf{Jeux Universitaires et Scolaires (JUSCO)} au niveau régional puis national. Participer activement au festival de classe, c'est donc s'inscrire dans une dynamique nationale de valorisation du patrimoine culturel camerounais.
\end{savaistu}

\courssub{Exercices d'intégration (auto-évaluation)}
\begin{exercice}[Auto-évaluation — Préparation au Probatoire]
\begin{enumerate}
\item \textbf{Reformulation} : Transformez en style formel : « Hey boss, can we use the hall? We wanna do a party. » $\to$ \emph{...}
\item \textbf{Modaux} : Complétez : « Students \_\_\_\_\_\_\_\_ respect the schedule. » (obligation forte) ; « We \_\_\_\_\_\_\_\_ invite parents. » (suggestion) ; « The noise \_\_\_\_\_\_\_\_ disturb the neighbours. » (déduction forte).
\item \textbf{Present perfect / Past perfect} : Mettez le verbe entre parenthèses au temps qui convient : « By the time the Principal arrived, we \_\_\_\_\_\_\_\_ (already / set up) the stands. » ; « We \_\_\_\_\_\_\_\_ (contact) three sponsors so far. »
\item \textbf{Let's / Let's not} : Proposez deux suggestions pour améliorer la communication du festival.
\item \textbf{Connecteurs} : Reliez les deux phrases avec \emph{so that} / \emph{in order to} : « We printed bilingual posters. Everyone understands the programme. »
\end{enumerate}
\end{exercice}

\begin{corrige}[Exercice d'intégration]
\begin{enumerate}
\item \emph{Monsieur le Proviseur, nous sollicitons respectueusement l'autorisation d'utiliser la salle polyvalente pour l'organisation d'un festival interculturel.}
\item \emph{must / ought to / should} (obligation) ; \emph{should / ought to} (conseil) ; \emph{must} (déduction).
\item \emph{had already set up} (past perfect — antériorité par rapport à \emph{arrived}) ; \emph{have contacted} (present perfect — lien avec le présent).
\item \emph{Let's create a WhatsApp group for updates.} / \emph{Let's not forget to post on the school Facebook page.}
\item \emph{We printed bilingual posters so that / in order that everyone understands the programme.} / \emph{We printed bilingual posters in order to / so as to make the programme understandable to everyone.}
\end{enumerate}
\end{corrige}

\begin{aretenir}[Bilan de l'unité « Family and Social Life »]
Vous savez maintenant : (1) interagir poliment avec pairs, adultes, autorités lors d'événements communautaires ; (2) utiliser les modaux pour demander la permission, exprimer la possibilité, l'obligation, la déduction ; (3) comprendre des textes sur le travail communautaire et la médiation ; (4) exprimer but, résultat, préférence ; (5) lire/écrire sur les activités récréatives avec \emph{present perfect} / \emph{past perfect} ; (6) formuler des suggestions avec \emph{Let's} ; (7) \textbf{mobiliser tout cela dans un projet réel}. Ces compétences sont transférables à tout contexte technique : gestion de projet, communication professionnelle, résolution de conflits.
\end{aretenir}

\newpage

\coursec{Activité d'intégration}

\coursec{Family and Social Life -- Intégration : Festival Interculturel}

\courssub{Objectifs de l'activité}
\begin{itemize}
    \item Mobiliser l'ensemble des savoirs de l'unité : vocabulaire des fêtes communautaires, modaux (\cle{can/may}, \cle{could}, \cle{ought to}, \cle{must}, \cle{needn't}, \cle{don't have to}), temps parfaits (\cle{present perfect}, \cle{past perfect}), structures de but/résultat/préférence, \cle{Let's/Let's not}.
    \item Mobiliser les savoir-faire : interaction orale (demander permission, négocier, médiatiser), compréhension orale/écrite, production écrite formelle (lettre, discours, affiche), argumentation.
    \item Travailler en équipe dans une situation authentique : organisation d'un festival interculturel scolaire au Cameroun.
    \item Développer l'autonomie, le sens des responsabilités et l'esprit citoyen (vivre-ensemble, résolution pacifique des conflits).
\end{itemize}

\courssub{Situation}
\textbf{Contexte :} Le Lycée Technique de Dschang (Région de l'Ouest) accueille, les \textbf{15 et 16 mars 2026}, la \textbf{2\textsuperscript{e} édition du Festival Interculturel des Lycées Techniques de l'Ouest}. Trois délégations d'écoles partenaires sont attendues :
\begin{itemize}
    \item Lycée Technique de Bafoussam : \textbf{42 élèves} + 3 encadreurs.
    \item Lycée Technique de Foumban : \textbf{38 élèves} + 3 encadreurs.
    \item Lycée Technique de Mbouda : \textbf{35 élèves} + 2 encadreurs.
\end{itemize}
Le lycée hôte mobilise \textbf{50 élèves volontaires} (commissions : logistique, communication, protocole, médiation, restauration) et \textbf{5 enseignants} référents.

\textbf{Contraintes réelles :}
\begin{itemize}
    \item \textbf{Budget prévisionnel alloué par la Région :} \textbf{750\,000 FCFA} (postes : repas, transport interne, matériel sono, décorations, prix, imprévus 10\%).
    \item \textbf{Règlement intérieur (extrait) :} « Toute manifestation collective nécessite une autorisation écrite du Proviseur 15 jours avant la date. L'utilisation du terrain de sport est interdite après 22 h. Le bruit excessif est prohibé dans les dortoirs après 21 h 30. »
    \item \textbf{Infrastructures :} 1 terrain de sport (foot/hand), 1 salle polyvalente (capacité 200), 2 dortoirs visiteurs (capacité 60 lits chacun), 1 réfectoire, 1 scène extérieure.
    \item \textbf{Lexique Ghomala' de base fourni :} \emph{Ndewo} (Bonjour), \emph{Merci} = \emph{Ayafwa}, \emph{Paix} = \emph{Mbong}, \emph{Ensemble} = \emph{Ngwa}, \emph{Fête} = \emph{Tsï}, \emph{Bienvenue} = \emph{Zoa}.
\end{itemize}

\textbf{Incident déclencheur (Jour 1, 20 h 30) :} La délégation de Bafoussam organise une répétition de danse traditionnelle avec tam-tams sur le terrain de sport. La délégation de Mbouda, logée dans le dortoir adjacent, proteste : le bruit empêche les élèves de se reposer avant les épreuves sportives du lendemain (8 h). Les esprits s'échauffent ; des propos blessants sont échangés (« Vous manquez de respect », « Vous êtes égoïstes »).

\courssub{Consignes}
Les élèves sont répartis en \textbf{5 commissions} (6 à 8 élèves chacune). Chaque commission rend un livrable écrit et/ou oral. Les consignes sont progressives et imbriquées.

\begin{enumerate}
    \item \textbf{Commission « Administration \& Budget » (Lettre officielle + Budget)} \\
    \textbf{1.1} Rédigez la \textbf{lettre formelle} adressée au Proviseur (format administratif : en-tête, objet, formules de politesse). Demandez l'autorisation du festival et le déblocage du budget. Utilisez : \cle{may/can} (permission formelle), \cle{present perfect} (bilan de la 1\textsuperscript{re} édition : \emph{We have successfully organized...}), \cle{ought to/should} (recommandations). \\
    \textbf{1.2} Complétez le \textbf{budget prévisionnel} (tableau vierge fourni) en répartissant les 750\,000 FCFA. Justifiez chaque poste en 1 phrase (ex. : \emph{We need 200\,000 FCFA for meals so that all participants eat properly}).
    \item \textbf{Commission « Communication \& Affiche » (Affiche trilingue + Programme)} \\
    \textbf{2.1} Concevez l'\textbf{affiche} (format A3, portrait) en \textbf{français / anglais / ghomala'}. Champs obligatoires : titre, dates, lieu, logo lycée, partenaires, slogan (\emph{Unity in Diversity / Mbong ngwa}), programme horaire simplifié (8 créneaux). \\
    \textbf{2.2} Écrivez 5 phrases de \textbf{présentation orale} de l'affiche (2 min) utilisant : \cle{Let's} + infinitif, \cle{so that / in order to} (but), \cle{present perfect} (actions préparatoires déjà faites).
    \item \textbf{Commission « Médiation \& Règlement » (Scénario de médiation)} \\
    \textbf{3.1} Rédigez un \textbf{scénario de médiation} (dialogue, 10--12 répliques) entre 2 « \emph{elders} » élèves (médiateurs), 1 représentant de Bafoussam, 1 représentant de Mbouda, 1 enseignant observateur. Intégrez : \cle{must} (obligation/règle), \cle{ought to/ought not to} (conseil moral), \cle{needn't/don't have to} (absence de nécessité), \cle{can/could} (propositions de solution). \\
    \textbf{3.2} Identifiez 3 \textbf{techniques de médiation} utilisées (ex. : reformulation, recherche d'intérêts communs, proposition d'alternative).
    \item \textbf{Commission « Protocole \& Discours » (Discours d'ouverture)} \\
    \textbf{4.1} Rédigez le \textbf{discours d'ouverture} (2 min, env. 260 mots) prononcé par le \emph{Président du Comité d'Organisation} (élève). Intégrez \textbf{obligatoirement} : \\
    \quad -- \emph{« Let's celebrate our diversity »} \\
    \quad -- \emph{« We have prepared this festival so that unity grows »} \\
    \quad -- \emph{« The elders had decided that... »} \\
    \quad -- Au moins 2 modaux (\cle{must}, \cle{should}, \cle{can}) + 1 \cle{past perfect} (antériorité). \\
    \textbf{4.2} Préparez la \textbf{lecture expressive} (ton, pauses, regard).
    \item \textbf{Commission « Évaluation \& Retour » (Grilles + Auto-évaluation)} \\
    \textbf{5.1} Concevez une \textbf{grille d'évaluation par compétences} (4 niveaux : \emph{Not acquired / Acquiring / Acquired / Expert}) pour les 4 livrables ci-dessus (critères : langue, contenu, forme, travail d'équipe). \\
    \textbf{5.2} Chaque élève remplit une \textbf{fiche d'auto-évaluation} (3 forces, 2 axes de progrès, 1 action concrète pour la suite).
\end{enumerate}

\courssub{Ressources à mobiliser}
\begin{definition}[Modaux clés -- Rappel]
\begin{center}
\begin{tabularx}{\linewidth}{|l|X|X|}\hline
\textbf{Modal} & \textbf{Fonction} & \textbf{Exemple} \\ \hline
\cle{can / may} & Permission (may = plus formel) & \emph{May we use the sports field?} \\ \hline
\cle{could} & Possibilité / Politesse & \emph{Could we move the rehearsal to the hall?} \\ \hline
\cle{must} & Obligation forte / Déduction & \emph{You must respect the curfew.} / \emph{They must be tired.} \\ \hline
\cle{ought to / ought not to} & Conseil moral / Devoir & \emph{We ought to listen to each other.} \\ \hline
\cle{should} & Conseil / Recommandation & \emph{You should apologize.} \\ \hline
\cle{needn't} & Absence de nécessité (interne) & \emph{You needn't worry; we have a plan.} \\ \hline
\cle{don't have to} & Absence d'obligation (externe) & \emph{We don't have to pay for the transport.} \\ \hline
\end{tabularx}
\end{center}
\end{definition}

\begin{definition}[Temps parfaits -- Rappel]
\begin{itemize}
    \item \textbf{Present Perfect} (\emph{have/has + V-ed/participe passé}) : lien passé-présent, bilan, expérience. \emph{We have prepared the posters.}
    \item \textbf{Past Perfect} (\emph{had + V-ed/participe passé}) : antériorité dans le passé. \emph{The elders had decided the rules before the conflict started.}
\end{itemize}
\end{definition}

\begin{definition}[Structures de but, résultat, préférence]
\begin{itemize}
    \item \textbf{But} : \cle{so that} + \emph{can/could/will/would} ; \cle{in order to / to} + \emph{infinitif}. \emph{We moved the rehearsal so that the others could sleep.}
    \item \textbf{Résultat} : \cle{so} + phrase ; \cle{such... that}. \emph{The noise was so loud that nobody could sleep.}
    \item \textbf{Préférence} : \cle{prefer + -ing / to + infinitif} ; \cle{would rather + infinitif sans to}. \emph{I prefer dancing to singing.} / \emph{I'd rather rehearse in the hall.}
\end{itemize}
\end{definition}

\begin{definition}[Let's / Let's not]
\emph{Let's} + \textbf{infinitif sans to} = proposition inclusive (1\textsuperscript{re} personne du pluriel). \emph{Let's not} = proposition négative. \emph{Let's celebrate!} / \emph{Let's not argue.}
\end{definition}

\courssub{Corrigés commentés des commissions}

\begin{corrige}[Exercice 1 -- Commission Administration \& Budget]
\textbf{Lettre formelle (extrait corrigé)}

\begin{center}
\begin{tabularx}{\linewidth}{|X|}\hline
\textbf{Lycée Technique de Dschang} \\
B.P. 124, Dschang -- Région de l'Ouest \\
Tél. : 2XX XX XX XX \quad Email : lt.dschang@minsec.cm \\ \hline
\textbf{Objet : Demande d'autorisation et de budget pour le Festival Interculturel des Lycées Techniques de l'Ouest (15--16 mars 2026)} \\ \hline
Dschang, le 28 février 2026 \\ \hline
Monsieur le Proviseur, \\ \hline
Nous, les membres du Comité d'Organisation des Élèves, \textbf{have the honour to} solliciter votre haute bienveillance pour l'organisation de la 2\textsuperscript{e} édition du Festival Interculturel. \textbf{We have successfully organized} la première édition en mars 2025, \textbf{which welcomed} 110 participants et \textbf{strengthened} les liens inter-écoles. \\ \hline
Cette année, nous \textbf{expect} 115 élèves et 8 encadreurs de Bafoussam, Foumban et Mbouda. \textbf{May we} utiliser le terrain de sport, la salle polyvalente et les deux dortoirs visiteurs du 14 mars (installation) au 16 mars (départ) ? \textbf{We ought to} respecter le règlement intérieur, notamment l'interdiction de bruit après 21 h 30. \\ \hline
Le budget prévisionnel s'élève à \textbf{750\,000 FCFA} (détail ci-joint). \textbf{We would be grateful if you could} autoriser le déblocage des fonds avant le 10 mars. \textbf{We needn't} rappeler que la Région a déjà validé l'enveloppe. \\ \hline
En vous remerciant par avance, nous vous prions d'agréer, Monsieur le Proviseur, l'expression de notre respectueuse considération. \\ \hline
\textit{Fait à Dschang, le 28/02/2026} \\
\textbf{Le Président du Comité} \quad \textbf{L'Enseignant Référent} \\
(Nom + signature) \quad (Nom + signature) \\ \hline
\end{tabularx}
\end{center}

\textbf{Commentaires :} \cle{have the honour to} (formule figée formelle) ; \cle{have successfully organized} (present perfect = bilan passé avec effet présent) ; \cle{May we} (permission formelle) ; \cle{ought to} (devoir moral/règlement) ; \cle{would be grateful if you could} (politesse conditionnelle) ; \cle{needn't} (absence de nécessité interne).

\textbf{Budget prévisionnel (répartition justifiée -- Total = 750\,000 FCFA)}

\begin{center}
\begin{tabularx}{\linewidth}{|l|r|X|}\hline
\textbf{Poste} & \textbf{Montant (FCFA)} & \textbf{Justification (but : so that / in order to)} \\ \hline
Repas (123 pers. $\times$ 2 jours $\times$ 1\,500) & 369\,000 & \emph{We need this amount so that all participants eat properly.} \\ \hline
Transport interne (navettes, carburant) & 120\,000 & \emph{We allocate this in order to ensure safe movements.} \\ \hline
Matériel sono / scène & 100\,000 & \emph{This is necessary so that performances are heard by all.} \\ \hline
Décorations / Affiches / Goodies & 50\,000 & \emph{We buy decorations so that the venue looks festive.} \\ \hline
Prix \& Récompenses & 36\,000 & \emph{We provide prizes so that winners feel valued.} \\ \hline
Imprévus (10\%) & 75\,000 & \emph{We keep a reserve in case unexpected costs arise.} \\ \hline
\textbf{Total} & \textbf{750\,000} & \\ \hline
\end{tabularx}
\end{center}
\end{corrige}

\begin{corrige}[Exercice 2 -- Commission Communication \& Affiche]
\textbf{Affiche trilingue (contenu textuel -- Format A3 Portrait)}

\begin{center}
\begin{tabularx}{\linewidth}{|X|}\hline
\rowcolor{popblueL} \textbf{FESTIVAL INTERCULTUREL 2026 / INTERCULTURAL FESTIVAL 2026 / TSÏ NGWA 2026} \\ \hline
\textbf{Unity in Diversity / Mbong ngwa} \\ \hline
\textbf{Lycée Technique de Dschang -- 15 \& 16 Mars 2026} \\ \hline
\textbf{PROGRAMME / SCHEDULE / PROGRAMME (GHOMALA')} \\ \hline
\textbf{15/03} -- 08h00 : Accueil / Welcome / Zoa \\
\textbf{15/03} -- 09h30 : Cérémonie d'ouverture / Opening Ceremony / Tsï za \\
\textbf{15/03} -- 11h00 : Expositions culturelles / Cultural Displays / Tsï nkap \\
\textbf{15/03} -- 14h00 : Tournoi sportif / Sports Tournament / Mba sport \\
\textbf{15/03} -- 19h00 : Soirée danses / Dance Evening / Tsï mbap \\
\textbf{16/03} -- 08h00 : Conférence « Vivre-ensemble » / Talk ``Living Together'' / Mbong ngwa \\
\textbf{16/03} -- 11h00 : Remise des prix / Awards / Ngu'a \\
\textbf{16/03} -- 13h00 : Déjeuner de clôture / Closing Lunch / Nyam za \\ \hline
\textbf{Partenaires :} Région de l'Ouest -- MINSEC -- Collectivités locales \\ \hline
\end{tabularx}
\end{center}

\textbf{Présentation orale (2 min) -- Modèle :}
\emph{``Good morning everyone. \textbf{Let's celebrate our diversity} today! \textbf{We have prepared} this festival \textbf{so that} unity grows among our schools. \textbf{We have designed} this trilingual poster \textbf{in order to} include our local language, Ghomala'. \textbf{Let's not forget} that culture builds peace. Thank you / Ayafwa / Zoa.''}
\end{corrige}

\begin{corrige}[Exercice 3 -- Commission Médiation \& Règlement]
\textbf{Scénario de médiation (12 répliques)}

\begin{center}
\begin{tabularx}{\linewidth}{|l|X|}\hline
\textbf{Personnage} & \textbf{Réplique (avec cibles grammaticales)} \\ \hline
\textbf{Médiateur 1 (Élève)} & \emph{``Good evening. \textbf{We must} respect the school rules: no noise after 21h30. \textbf{Could} each side explain their need calmly?''} \\ \hline
\textbf{Représentant Bafoussam} & \emph{``We \textbf{have to} rehearse for tomorrow's competition. The hall is occupied. \textbf{Can} we use the field just until 21h?''} \\ \hline
\textbf{Médiateur 2 (Élève)} & \emph{``\textbf{You ought to} understand that Mbouda students \textbf{need} to rest. \textbf{They have} sports events at 8h. \textbf{What about} moving to the hall at 20h?''} \\ \hline
\textbf{Représentant Mbouda} & \emph{``We \textbf{can't} sleep with the drums. \textbf{You should} have asked us before. \textbf{We didn't agree} to this noise.''} \\ \hline
\textbf{Enseignant observateur} & \emph{``The rule \textbf{must} be applied. \textbf{However}, \textbf{we could} split the time: field until 20h30, then hall. \textbf{Does that work} for both?''} \\ \hline
\textbf{Médiateur 1} & \emph{``\textbf{Let's} propose a schedule: Bafoussam on field 19h--20h30, then hall 20h30--21h30. Mbouda rests from 21h. \textbf{Do you agree?''} } \\ \hline
\textbf{Représentant Bafoussam} & \emph{``\textbf{We can} accept that. \textbf{We ought not to} disturb others. Sorry for the words earlier.''} \\ \hline
\textbf{Représentant Mbouda} & \emph{``\textbf{We accept} too. \textbf{We needn't} argue. \textbf{Let's} focus on the festival.''} \\ \hline
\textbf{Médiateur 2} & \emph{``\textbf{The elders have decided that} the field closes at 20h30 sharp. \textbf{We will} monitor. \textbf{Thank you} for your maturity.''} \\ \hline
\end{tabularx}
\end{center}

\textbf{Techniques de médiation identifiées :}
\begin{enumerate}
    \item \textbf{Reformulation} (Médiateur 1 résume les besoins : répétition vs repos).
    \item \textbf{Recherche d'intérêts communs} (réussite du festival, respect du règlement, santé des élèves).
    \item \textbf{Proposition d'alternative concrète} (partage horaire terrain/salle + horaire précis).
\end{enumerate}
\end{corrige}

\begin{corrige}[Exercice 4 -- Commission Protocole \& Discours]
\textbf{Discours d'ouverture (version complète, ~260 mots)}

\begin{center}
\begin{tabularx}{\linewidth}{|X|}\hline
\textbf{DISCOURS D'OUVERTURE -- FESTIVAL INTERCULTUREL 2026} \\
\textit{Prononcé par le Président du Comité d'Organisation, Lycée Technique de Dschang, 15 mars 2026, 09h30.} \\ \hline
\textbf{Distinguished Proviseur, Dear Teachers, Dear Delegates from Bafoussam, Foumban and Mbouda, Dear Friends,} \\
\textbf{Let's celebrate our diversity!} Today, our school opens its doors to welcome 115 young talents from three brother schools. \textbf{We have prepared} this festival \textbf{so that} unity grows beyond our differences. Since January, the committee \textbf{has worked} tirelessly: we \textbf{have designed} posters, \textbf{have written} letters, \textbf{have planned} meals and schedules. \textbf{The elders had decided} last year that this festival must become a tradition, and we \textbf{are honoured} to continue it. \\
\textbf{We must} remember that culture is not a competition but a bridge. \textbf{You should} see this festival as a space to share, to listen, to learn. \textbf{You can} dance, sing, play, debate -- but \textbf{you ought to} do it with respect. \textbf{The elders have decided that} the sports field closes at 20h30 to protect everyone's rest; \textbf{we needn't} see this as a limit, but as a care for each other. \\
\textbf{Let's not} allow misunderstandings to spoil our joy. \textbf{Let's} build memories that will last. \textbf{We would rather} create friendship than rivalry. \textbf{Let's} make these two days a model of \emph{Mbong ngwa} -- peace together. \\
\textbf{Thank you. / Ayafwa. / Zoa.} \\ \hline
\end{tabularx}
\end{center}

\textbf{Commentaires grammaticaux :} \cle{Let's celebrate} (proposition inclusive) ; \cle{have prepared / have worked / have designed / have written / have planned} (present perfect = actions passées récentes, résultat visible) ; \cle{had decided} (past perfect = antériorité par rapport à « last year ») ; \cle{must / should / can / ought to / needn't} (modaux variés) ; \cle{so that} (but) ; \cle{would rather} (préférence).
\end{corrige}

\begin{corrige}[Exercice 5 -- Commission Évaluation \& Retour]
\textbf{Grille d'évaluation par compétences (Extrait)}

\begin{center}
\begin{tabularx}{\linewidth}{|l|X|X|X|X|}\hline
\textbf{Critère} & \textbf{Not acquired (1)} & \textbf{Acquiring (2)} & \textbf{Acquired (3)} & \textbf{Expert (4)} \\ \hline
\textbf{Vocabulaire festivals/conflicts} & Mots isolés, répétitifs & Vocabulaire de base, quelques erreurs & Vocabulaire varié, précis, contextuel & Vocabulaire riche, idiomatique, créatif \\ \hline
\textbf{Modaux (permission/obligation/conseil)} & Absents ou mal utilisés & 1--2 modaux corrects & 3--4 modaux corrects, fonctions claires & 5+ modaux, nuances fines (may vs can, needn't vs don't have to) \\ \hline
\textbf{Present / Past Perfect} & Absents ou confusion & 1 forme correcte & Les deux formes correctes, sens justifié & Utilisation naturelle, antériorité claire \\ \hline
\textbf{Structures but/résultat/préférence} & Absentes & 1 structure (so that / to) & 2 structures variées & 3+ structures, enchaînées logiquement \\ \hline
\textbf{Let's / Let's not} & Absent & 1 occurrence correcte & 2 occurrences, sens inclusif & Intégré naturellement dans discours \\ \hline
\textbf{Cohérence / Argumentation} & Décousu, hors sujet & Idées principales présentes, transitions faibles & Structure claire, arguments liés & Argumentation fluide, connecteurs variés, ton adapté \\ \hline
\textbf{Travail d'équipe / Respect délais} & Participation nulle & Participation inégale, retards & Participation active, respect délais & Leadership, entraide, anticipation \\ \hline
\end{tabularx}
\end{center}

\textbf{Auto-évaluation élève (modèle rempli) :}
\begin{itemize}
    \item \textbf{Forces :} (1) J'ai bien utilisé \emph{present perfect} dans la lettre. (2) J'ai proposé une solution créative en médiation (partage horaire). (3) Mon discours avait un ton engageant.
    \item \textbf{Axes de progrès :} (1) Encore hésitant sur \emph{needn't vs don't have to}. (2) Prononciation de \emph{ought to} /ɔːt tuː/.
    \item \textbf{Action concrète :} Réviser les modaux avec l'appli OnBuch+ (exercices ciblés) et m'enregistrer pour la prononciation.
\end{itemize}
\end{corrige}

\courssub{Bilan de l'activité}
Cette activité d'intégration a permis de \textbf{mettre en évidence} :

\begin{enumerate}
    \item \textbf{La transférabilité des savoirs :} Les élèves ont réinvesti le vocabulaire (\emph{celebration, dispute, mediation, rehearsal, curfew, diversity, unity}), les modaux (permission formelle \emph{may}, conseil \emph{ought to}, obligation \emph{must}, absence de nécessité \emph{needn't}), les temps parfaits (bilan \emph{present perfect}, antériorité \emph{past perfect}), les structures de but (\emph{so that, in order to}) et \emph{Let's} dans \textbf{une seule tâche complexe}, et non plus en exercices isolés.
    \item \textbf{L'ancrage camerounais réaliste :} Le contexte (Dschang, Bafoussam, Foumban, Mbouda, budget en FCFA, règlement intérieur, langue ghomala') a suscité l'engagement : les élèves \emph{se reconnaissent} dans la situation.
    \item \textbf{L'articulation oral/écrit :} L'affiche (écrit visuel) soutient la présentation orale ; le scénario de médiation (écrit dialogué) prépare le jeu de rôle ; le discours (écrit formel) devient performance orale.
    \item \textbf{La dimension citoyenne :} La médiation par les \emph{elders} élèves modélise la résolution pacifique des conflits (compétence transversale MINESEC : « vivre-ensemble », « prévention de la violence »). L'auto-évaluation développe la métacognition.
    \item \textbf{Les points de vigilance pour la remédiation :} Confusion persistante \emph{needn't / don't have to} ; difficulté à placer \emph{past perfect} spontanément ; prononciation des modaux (\emph{ought to, mustn't}) ; gestion du temps dans le travail collaboratif (certaines commissions ont rendu incomplet).
\end{enumerate}

\textbf{Pistes de prolongement (remédiation / approfondissement) :}
\begin{itemize}
    \item Atelier ciblé « Modaux : \emph{needn't} vs \emph{don't have to} » (exercices de transformation + mini-dialogues).
    \item Entraînement \emph{past perfect} : réécrire le journal de bord du festival (\emph{By 6 p.m., we had set up the stage...}).
    \item Enregistrement audio des discours + grille d'auto-correction phonologique (accentuation, intonation \emph{Let's}).
    \item Écriture d'un article pour le journal du lycée (genre journalistique, passé composé / \emph{present perfect}).
\end{itemize}

\begin{aretenir}[Synthèse de l'unité -- Family and Social Life]
\begin{itemize}
    \item \textbf{Vocabulaire clé :} \emph{community, festival, celebration, neighbour, authority, dispute, mediation, rehearsal, curfew, diversity, unity, reconciliation, recreational activity}.
    \item \textbf{Grammaire essentielle :}
    \begin{itemize}
        \item Modaux : \cle{can/may} (permission), \cle{could} (possibilité/politesse), \cle{must} (obligation/déduction), \cle{should/ought to} (conseil), \cle{needn't} (pas nécessaire), \cle{don't have to} (pas d'obligation).
        \item Temps : \cle{Present Perfect} (have/has + V3) vs \cle{Past Perfect} (had + V3).
        \item Structures : \cle{so that / in order to} (but), \cle{so / such... that} (résultat), \cle{would rather / prefer} (préférence).
        \item \cle{Let's / Let's not} + BV (proposition/suggestion).
    \end{itemize}
    \item \textbf{Compétences visées :} Interagir en société, résoudre un conflit, organiser un événement, s'exprimer formellement (écrit/oral), s'auto-évaluer.
\end{itemize}
\end{aretenir}

\cle{Bilan final} : L'activité a atteint son objectif d'intégration. Les productions montrent une \textbf{progression nette} par rapport aux séances d'apprentissage isolées. L'évaluation par grille + auto-évaluation donne des données fiables pour le conseil de classe et le livret de compétences.

\newpage

\coursec{Méthodes et exercices résolus}

\coursec{Family and Social Life}

\courssub{Fiche méthode 1 : Choisir le bon modal (permission, possibilité, obligation, déduction)}

\begin{methode}[How to choose the correct modal]
To choose the right modal in an exam sentence, follow these steps:
\begin{enumerate}
\item \textbf{Identify the meaning} the sentence expresses: \cle{permission} (autorisation) -- \emph{can/may}, \cle{possibility} (possibilité) -- \emph{can/could}, \cle{obligation} (obligation) -- \emph{must, should, ought to}, \cle{absence of obligation} (absence d'obligation) -- \emph{needn't, don't have to}, or \cle{deduction} (déduction) -- \emph{must}.
\item \textbf{Check the register} (niveau de langue): \emph{may} is more formal than \emph{can} (use \emph{may} with school authorities, adults; \emph{can} with mates).
\item \textbf{Distinguish} \emph{mustn't} (interdiction -- it is forbidden) from \emph{needn't / don't have to} (pas d'obligation -- it is not necessary).
\item \textbf{For deduction}, use \emph{must + bare infinitive} (infinitif sans \emph{to}) when you are almost sure something is true: \emph{She speaks Ewondo fluently -- she must be from Yaoundé.}
\item \textbf{Check the form}: modal + \cle{bare infinitive} (infinitif sans \emph{to}), except \emph{ought \textbf{to}} and \emph{have \textbf{to}}.
\end{enumerate}
\begin{aretenir}[Key forms / Formes clés]
can/may + BV (permission) ; can/could + BV (possibility) ; should/ought to + BV (advice/conseil) ; needn't / don't have to + BV (no obligation) ; must + BV (obligation or deduction).
\end{aretenir}
\end{methode}

\begin{exoresolu}[Choosing the right modal -- Probatoire type]
Complete each sentence with the correct modal: \emph{can, may, must, mustn't, needn't, should, ought to}.
\begin{enumerate}
\item You \ldots\ wear a uniform to attend the school cultural festival. It is the rule.
\item \ldots\ I use your drum for the traditional dance, Mr Ndi? (polite, to an adult)
\item You \ldots\ bring food; the committee will provide everything. (not necessary)
\item The hall is full of people dancing. It \ldots\ be the Ngondo festival!
\item We \ldots\ respect our elders during community celebrations. (moral duty)
\item You \ldots\ smoke inside the school hall. It is strictly forbidden.
\end{enumerate}
\tcblower
\textbf{Step 1 -- Meaning of each sentence.}
\begin{enumerate}
\item ``It is the rule'' $\Rightarrow$ \textbf{obligation}: \emph{must}. \textbf{You must wear a uniform\ldots}
\item Polite request to an adult (Mr Ndi) $\Rightarrow$ formal permission: \emph{may}. \textbf{May I use your drum\ldots?}
\item ``Not necessary'' $\Rightarrow$ absence of obligation: \emph{needn't}. \textbf{You needn't bring food\ldots}
\item Conclusion from evidence (hall full of dancers) $\Rightarrow$ deduction: \emph{must}. \textbf{It must be the Ngondo festival!}
\item Moral duty, advice $\Rightarrow$ \emph{should} or \emph{ought to}. \textbf{We should respect our elders\ldots} / \textbf{We ought to respect\ldots}
\item ``Strictly forbidden'' $\Rightarrow$ interdiction: \emph{mustn't}. \textbf{You mustn't smoke inside\ldots}
\end{enumerate}
\textbf{Step 2 -- Form check.} Every modal is followed by a bare infinitive (\emph{wear, use, bring, be, respect, smoke}); \emph{ought to} keeps its \emph{to}.
\textbf{Conclusion:} 1. must -- 2. May -- 3. needn't -- 4. must -- 5. should/ought to -- 6. mustn't.
\end{exoresolu}

\courssub{Fiche méthode 2 : Exprimer le but, le résultat, la préférence et l'obligation}

\begin{methode}[Expressing purpose, result, preference and obligation]
\begin{enumerate}
\item \textbf{Purpose} (but) : \cle{to + bare infinitive}, \cle{in order to + bare infinitive}, \cle{so as to + bare infinitive}, or \cle{so that + subject + can/could}. Example: \emph{We organised a meeting \textbf{to settle} the dispute.}
\item \textbf{Result} (résultat) : \cle{so + adjective + that}, \cle{such (a) + noun + that}, \cle{therefore}, \cle{as a result}. Example: \emph{The music was \textbf{so} loud \textbf{that} everybody danced.}
\item \textbf{Preference} (préférence) : \cle{prefer + noun/\emph{-ing} + to}, \cle{would rather + bare infinitive + than}, \cle{would prefer to + bare infinitive}. Example: \emph{I \textbf{would rather play} football \textbf{than} watch TV.}
\item \textbf{Obligation} : \emph{must} (internal obligation, strong rule), \emph{have to} (external obligation), \emph{should/ought to} (advice/conseil).
\item \textbf{Transform the sentence} without changing the meaning: identify the link word, then apply the equivalent structure with the same tense.
\end{enumerate}
\begin{aretenir}[Key structures]
Purpose: \emph{to/in order to/so as to + BV} ; \emph{so that + can/could}. Result: \emph{so/such\ldots that}. Preference: \emph{would rather + BV + than} ; \emph{prefer + -ing/noun + to}.
\end{aretenir}
\end{methode}

\begin{exoresolu}[Rewriting: purpose and result]
Rewrite each sentence without changing its meaning, using the word in brackets.
\begin{enumerate}
\item The youths cleaned the market square because they wanted to prepare the festival. (in order to)
\item The dispute was very serious, so the chief called a palaver. (such\ldots that)
\item I prefer dancing to singing. (would rather)
\end{enumerate}
\tcblower
\textbf{1. Purpose.} ``Because they wanted to'' expresses a purpose. Replace it by \emph{in order to + bare infinitive}:\\
\textbf{The youths cleaned the market square in order to prepare the festival.}\\
\textbf{2. Result.} ``Very serious $\Rightarrow$ consequence'' becomes \emph{such + adjective + noun + that}. The noun is \emph{dispute} (countable, singular) $\Rightarrow$ \emph{such a serious dispute that}:\\
\textbf{It was such a serious dispute that the chief called a palaver.}\\
\textbf{3. Preference.} \emph{prefer X to Y} = \emph{would rather + bare infinitive X + than + bare infinitive Y}. After \emph{would rather}, use the bare infinitive:\\
\textbf{I would rather dance than sing.}\\
\textbf{Conclusion:} the three sentences keep their meaning; only the linking structure changes, and the verb forms (bare infinitive after \emph{in order to} and \emph{would rather}) are respected.
\end{exoresolu}

\courssub{Fiche méthode 3 : Present perfect et past perfect}

\begin{methode}[Using the present perfect and the past perfect]
\begin{enumerate}
\item \textbf{Present perfect} = \cle{have/has + past participle}. Use it for: an action with a \emph{present result} (\emph{I have lost my ticket}), an experience (\emph{Have you ever attended the Ngondo?}), an action that started in the past and continues (\emph{for/since}).
\item \textbf{Past perfect} = \cle{had + past participle}. Use it for the \emph{first of two past actions}: \emph{When we arrived, the match \textbf{had started}.}
\item \textbf{Time markers} (repères temporels): present perfect $\rightarrow$ \emph{just, already, yet, ever, never, for, since} ; past perfect $\rightarrow$ \emph{before, after, when, by the time}.
\item \textbf{Never} use the present perfect with a finished past time (\emph{yesterday, last week, in 2019}) $\rightarrow$ use the simple past.
\item Memorise irregular past participles: go$\rightarrow$gone, take$\rightarrow$taken, see$\rightarrow$seen, write$\rightarrow$written, begin$\rightarrow$begun.
\end{enumerate}
\begin{aretenir}[Time markers]
Present perfect: \emph{just, already, yet, ever, never, for, since}. Past perfect: \emph{before, after, when, by the time}. Simple past: \emph{yesterday, last week, ago, in 2019}.
\end{aretenir}
\end{methode}

\begin{exoresolu}[Present perfect vs past perfect]
Put the verbs in brackets into the present perfect, past perfect or simple past.
\begin{enumerate}
\item Our school (organise) \ldots\ three cultural festivals since 2022.
\item When the guests arrived, the dancers (already / begin) \ldots\ the performance.
\item I (never / see) \ldots\ a juju dance before last year's festival.
\item The community (just / finish) \ldots\ the communal work. Look at the clean square!
\end{enumerate}
\tcblower
\textbf{1.} Marker \emph{since 2022} $\Rightarrow$ action continuing up to now $\Rightarrow$ present perfect: \textbf{Our school has organised three cultural festivals since 2022.}\\
\textbf{2.} Two past actions: the dancers began \emph{before} the guests arrived $\Rightarrow$ earlier action in the past perfect: \textbf{When the guests arrived, the dancers had already begun the performance.}\\
\textbf{3.} \emph{Never\ldots before last year's festival}: experience completed before a finished past moment (\emph{last year}) $\Rightarrow$ past perfect: \textbf{I had never seen a juju dance before last year's festival.}\\
\textbf{4.} Marker \emph{just} + visible present result (the square is clean now) $\Rightarrow$ present perfect: \textbf{The community has just finished the communal work.}\\
\textbf{Conclusion:} 1. has organised -- 2. had already begun -- 3. had never seen -- 4. has just finished.
\end{exoresolu}

\courssub{Fiche méthode 4 : Let's / Let's not pour proposer une activité}

\begin{methode}[Making suggestions with Let's / Let's not]
\begin{enumerate}
\item \textbf{Structure}: \cle{Let's + bare infinitive} (positive suggestion) ; \cle{Let's not + bare infinitive} (negative suggestion).
\item \textbf{Meaning}: \emph{Let's} = ``et si nous\ldots'' -- it includes the speaker and the listener(s).
\item \textbf{Never} write \emph{let's to go} or \emph{let's going}: only the \textbf{bare infinitive}.
\item \textbf{Answers}: accept with \emph{OK, let's! / That's a good idea!} ; refuse politely with \emph{I'd rather not / Let's not\ldots}
\item In a dialogue about recreational activities, combine \emph{Let's} with preference structures: \emph{Let's play football. -- I'd rather play handball.}
\end{enumerate}
\begin{aretenir}[Forme]
\emph{Let's} + BV (base verbale) ; \emph{Let's not} + BV. Pas de \emph{to}, pas de \emph{-ing}.
\end{aretenir}
\end{methode}

\begin{exoresolu}[Dialogue completion -- Let's / Let's not]
Complete the dialogue between two classmates with \emph{Let's}, \emph{Let's not} or \emph{I'd rather}.\par
\textbf{Akono:} The exams are over. (1) \ldots\ do something relaxing this weekend!\par
\textbf{Bella:} Good idea! (2) \ldots\ go to the stadium to watch the football match.\par
\textbf{Akono:} Hmm, (3) \ldots\ go to the stadium; it will be too crowded. (4) \ldots\ organise a dance competition in the school hall instead.\par
\textbf{Bella:} OK, (5) \ldots\ ! And we can invite the whole class.
\tcblower
\textbf{(1)} Akono makes a general suggestion $\Rightarrow$ \textbf{Let's do something relaxing this weekend!}\\
\textbf{(2)} Bella proposes an activity $\Rightarrow$ \textbf{Let's go to the stadium\ldots}\\
\textbf{(3)} Akono rejects the stadium (``too crowded'') $\Rightarrow$ negative suggestion: \textbf{Let's not go to the stadium\ldots}\\
\textbf{(4)} New proposal $\Rightarrow$ \textbf{Let's organise a dance competition\ldots}\\
\textbf{(5)} Bella accepts enthusiastically $\Rightarrow$ \textbf{OK, let's!}\\
\textbf{Form check:} every \emph{Let's / Let's not} is followed by a bare infinitive (\emph{do, go, organise}).\\
\textbf{Conclusion:} 1. Let's -- 2. Let's -- 3. Let's not -- 4. Let's -- 5. let's.
\end{exoresolu}

\courssub{Fiche méthode 5 : Comprendre un texte de lecture (reading)}

\begin{methode}[Reading a text about recreational activities]
\begin{enumerate}
\item \textbf{Read the title and the questions first}: they tell you what to look for.
\item \textbf{Skim} the text (lecture rapide) to get the general idea: who? where? what activity?
\item \textbf{Scan} for details: locate the key words of each question in the text (names, dates, places, numbers).
\item \textbf{Answer in full sentences}, using the tense of the question; copy only the useful information, never the whole paragraph.
\item \textbf{True/False questions}: quote or underline the line that justifies your answer; ``Not stated'' if the text says nothing.
\item \textbf{Vocabulary in context}: guess the meaning of an unknown word from the words around it (contexte).
\end{enumerate}
\end{methode}

\begin{exoresolu}[Reading comprehension]
Read the text and answer the questions in full sentences.\par
\emph{``Every Saturday morning, the youths of Bafoussam meet at the community field. They play football and handball, and the elders organise traditional wrestling. Last month, two quarters were in conflict over a piece of land. Instead of fighting, their football teams played a peace match. After the match, which ended 2--2, the two quarter heads shook hands and settled the dispute. Since then, the peace tournament has taken place every year.''}\par
\textbf{Questions:} (a) When do the youths meet? (b) Which activities do they practise? (c) How was the dispute settled? (d) How often does the peace tournament take place now?
\tcblower
\textbf{(a)} Scan for a time expression: ``Every Saturday morning''. $\Rightarrow$ \textbf{The youths meet every Saturday morning.}\\
\textbf{(b)} Scan for activities: football, handball, traditional wrestling. $\Rightarrow$ \textbf{They play football and handball, and the elders organise traditional wrestling.}\\
\textbf{(c)} Scan for ``dispute'': the teams played a peace match, then the quarter heads shook hands. $\Rightarrow$ \textbf{The dispute was settled after a peace match which ended 2--2: the two quarter heads shook hands and made peace.}\\
\textbf{(d)} Scan for frequency: ``every year''. $\Rightarrow$ \textbf{The peace tournament now takes place every year.}\\
\textbf{Conclusion:} each answer is a full sentence, in the correct tense (simple present for habits, simple past for the past event), and directly supported by the text.
\end{exoresolu}

\courssub{Fiche méthode 6 : Rédiger une composition argumentée (writing)}

\begin{methode}[Writing a composition: recreational activities and peace]
\begin{enumerate}
\item \textbf{Analyse the topic} and make a quick plan: introduction -- 2 or 3 arguments with examples -- conclusion.
\item \textbf{Introduction}: present the situation (recreational activities in the community) and your thesis (they bring peace).
\item \textbf{Body}: one paragraph = one argument + one example (a peace match, communal work, a festival). Use link words: \emph{first, moreover, for instance, as a result, finally}.
\item \textbf{Use the unit's grammar}: modals (\emph{we should\ldots, youths can\ldots}), purpose (\emph{they meet \textbf{in order to} dialogue}), present perfect (\emph{sport \textbf{has brought} peace}).
\item \textbf{Conclusion}: summarise and give your opinion or a recommendation (\emph{Let's promote\ldots}).
\item \textbf{Check}: subject-verb agreement, tenses, spelling, punctuation; 120--150 words for a Probatoire composition.
\end{enumerate}
\begin{aretenir}[Plan type]
Introduction (thesis) -- Argument 1 + example -- Argument 2 + example -- Conclusion (opinion/recommendation). Mots-liens: \emph{first, moreover, for instance, as a result, in conclusion}.
\end{aretenir}
\end{methode}

\begin{exoresolu}[Guided composition]
Write a short composition (about 120 words) on: \emph{``How recreational activities can settle community disputes and bring peace.''}
\tcblower
\textbf{Model answer:}\par
\emph{In our communities, disputes about land, money or politics often divide neighbours. Fortunately, recreational activities can help to settle these conflicts and bring peace.}\par
\emph{First, sport creates friendship. When youths from two rival quarters play football together, they learn to respect one another. For instance, in my village, a peace match was organised in order to reconcile two families, and it succeeded. Moreover, cultural festivals and communal work bring people together: they sing, dance and eat side by side, so that old quarrels are forgotten. As a result, the community becomes united and peaceful.}\par
\emph{In conclusion, recreational activities are not only for pleasure; they are powerful tools for peace. Our school authorities should therefore promote games and festivals. Let's all take part in them!}\par
\textbf{Why this is a good copy:} clear 3-paragraph plan; link words (\emph{first, moreover, for instance, as a result, in conclusion}); unit grammar used correctly (modals \emph{can, should}; purpose \emph{in order to}; suggestion \emph{Let's}); about 120 words; no agreement or tense error.
\end{exoresolu}

\courssub{Fiche méthode 7 : Speaking -- échanger des informations sur les fêtes communautaires}

\begin{methode}[Exchanging information about community celebrations]
\begin{enumerate}
\item \textbf{Prepare question words}: \emph{When} (date), \emph{Where} (place), \emph{Who} (people), \emph{What} (activities), \emph{How long} (duration), \emph{Why} (purpose).
\item \textbf{Ask politely} according to the person: \emph{Can I\ldots?} to a mate ; \emph{May I ask you\ldots? / Could you tell me\ldots?} to an adult or a school authority.
\item \textbf{React and keep the conversation going}: \emph{Really? That's interesting! / I see. / And what happens next?}
\item \textbf{Give information} with the unit's vocabulary: \emph{take part in, organise, attend, traditional dance, communal work, festival, ceremony}.
\item \textbf{Close the exchange}: \emph{Thank you for the information. / You're welcome.}
\end{enumerate}
\begin{aretenir}[Politesse]
Pair/ami: \emph{Can I\ldots?} ; Adulte/autorité: \emph{May I\ldots? / Could you\ldots?}
\end{aretenir}
\end{methode}

\begin{exoresolu}[Dialogue: interviewing a school authority]
You want information about your school's cultural week. Complete the interview with the principal's secretary using appropriate questions.\par
\textbf{You:} Good morning, Madam. (1) \ldots\ ?\par
\textbf{Secretary:} The cultural week will start on Monday 12th May.\par
\textbf{You:} (2) \ldots\ ?\par
\textbf{Secretary:} It will take place in the school hall and on the playground.\par
\textbf{You:} (3) \ldots\ ?\par
\textbf{Secretary:} All students may take part in the dances, songs and drama.\par
\textbf{You:} (4) \ldots\ ?\par
\textbf{Secretary:} You should register with your class teacher before Friday.\par
\textbf{You:} Thank you very much, Madam.
\tcblower
\textbf{(1)} The answer gives a \emph{date} $\Rightarrow$ question with \emph{When}, polite register: \textbf{May I know when the cultural week will start?} / \textbf{Could you tell me when the cultural week starts?}\\
\textbf{(2)} The answer gives a \emph{place} $\Rightarrow$ \emph{Where}: \textbf{Where will it take place?}\\
\textbf{(3)} The answer says \emph{who} participates $\Rightarrow$ \textbf{Who may take part in the activities?}\\
\textbf{(4)} The answer explains \emph{what to do} $\Rightarrow$ \textbf{What should I do to register?} / \textbf{How can I register?}\\
\textbf{Conclusion:} each question matches the information in the answer (question word + correct word order: question word + auxiliary + subject + verb), and the register is polite (\emph{May I, Could you, should}) because the speaker addresses a school authority.
\end{exoresolu}

\begin{attention}[Five common mistakes that cost marks in the Bac / Cinq erreurs qui coûtent des points au Bac]
\begin{enumerate}
\item \textbf{Adding \emph{to} after a modal}: never write \emph{must to go}, \emph{can to play}, \emph{should to respect}. A modal is followed by the \textbf{bare infinitive} (infinitif sans \emph{to}) -- except \emph{ought \textbf{to}} and \emph{have \textbf{to}}.
\item \textbf{Confusing \emph{mustn't} and \emph{needn't}/\emph{don't have to}}: \emph{mustn't} = interdiction (prohibition) ; \emph{needn't} = pas d'obligation (no obligation). Swapping them completely changes the meaning.
\item \textbf{Present perfect with a finished past time}: \emph{I have seen him yesterday} is wrong. With \emph{yesterday, last week, ago}, use the simple past: \emph{I saw him yesterday.}
\item \textbf{\emph{Let's} followed by the wrong form}: never \emph{let's to go} nor \emph{let's going} ; always \emph{Let's + bare infinitive}, and the negative is \emph{Let's \textbf{not} + BV} (not \emph{don't let's}).
\item \textbf{Off-topic composition or no plan}: writing without introduction or conclusion, forgetting connectors (\emph{first, moreover, as a result}), and neglecting subject-verb agreement (\emph{he play} instead of \emph{he plays}) lose both language and coherence marks. Always proofread your work.
\end{enumerate}
\end{attention}

\newpage

\coursec{Exercices}

\courssub{Je vérifie mes connaissances}

\begin{exercice}[QCM : Modaux et interactions sociales]
Choisissez la bonne réponse parmi les propositions A, B, C ou D.

\begin{enumerate}
\item \textbf{Student:} \emph{``\dots\ I leave the classroom to help with the festival preparations, Sir?''} \\
\textbf{Teacher:} \emph{``Yes, you may.''} \\
A) Must \quad B) Could \quad C) May \quad D) Ought to

\item \textbf{Neighbour:} \emph{``The noise from the celebration \dots\ be very loud tonight.''} \\
A) can \quad B) must \quad C) ought not to \quad D) needn't

\item \textbf{Student A:} \emph{``We \dots\ clean the village square before the Ngondo festival.''} \\
\textbf{Student B:} \emph{``Yes, it is our duty.''} \\
A) don't have to \quad B) must \quad C) could \quad D) may

\item \textbf{Principal:} \emph{``You \dots\ bring your parents' consent form for the trip.''} \\
A) can \quad B) may \quad C) have to \quad D) could
\end{enumerate}
\end{exercice}

\begin{exercice}[Vrai ou Faux justifié : Vocabulaire des fêtes communautaires]
Indiquez si les affirmations suivantes sont vraies (V) ou fausses (F). Justifiez votre réponse par une phrase en anglais.

\begin{enumerate}
\item ``A \cle{stakeholder} is a person who holds a stake (part/interest) in an organisation or event.''
\item ``\cle{Communal work} refers to paid labour done by a private company for the community.''
\item ``To \cle{mediate} a dispute means to take sides and punish the guilty party.''
\item ``\cle{Recreational activities} include traditional wrestling, dancing, and storytelling.''
\end{enumerate}
\end{exercice}

\begin{exercice}[Définitions et correspondance : Grammaire -- Temps parfaits]
Associez chaque phrase à la structure grammaticale correcte (Present Perfect ou Past Perfect) et expliquez brièvement la différence de sens (en français).

\begin{center}
\begin{tabularx}{\linewidth}{|l|X|X|}\hline
\textbf{Phrase} & \textbf{Structure} & \textbf{Justification (FR)} \\ \hline
1. ``The elders \textbf{have settled} the dispute.'' & A) Past Perfect & \\ \hline
2. ``The elders \textbf{had settled} the dispute before the festival started.'' & B) Present Perfect & \\ \hline
\end{tabularx}
\end{center}
\end{exercice}

\begin{exercice}[Application directe : Let's / Let's not]
Complétez les phrases suivantes en utilisant \textbf{Let's} ou \textbf{Let's not} suivi de l'infinitif sans \emph{to} (bare infinitive) du verbe entre parenthèses.

\begin{enumerate}
\item \dots\ \emph{(organise)} a cleaning campaign for the school yard tomorrow.
\item \dots\ \emph{(forget)} to invite the quarter head (\emph{chef de quartier}) to the opening ceremony.
\item \dots\ \emph{(prepare)} traditional dishes for the visitors.
\item \dots\ \emph{(make)} too much noise after 10 p.m. during the festival week.
\end{enumerate}
\end{exercice}

\courssub{Je m'entraîne}

\begin{exercice}[Dialogue complété : Demande d'autorisation (Modaux + Politesse)]
Complétez le dialogue entre un élève (E) et le chef de quartier (C) en choisissant le modal ou l'expression appropriée dans la liste : \emph{may, could, must, have to, don't have to, ought to}.

\begin{quote}
\textbf{E:} Good morning, Chief. \dots\ I speak with you for a moment?
\textbf{C:} Of course, young man. What can I do for you?
\textbf{E:} We \dots\ organise the end-of-year school festival next Saturday. We \dots\ ask for your permission to use the community hall.
\textbf{C:} That is a good initiative. You \dots\ submit a written request to my office first.
\textbf{E:} \dots\ we bring it tomorrow morning?
\textbf{C:} Yes, that is fine. But you \dots\ worry about the chairs; the council will provide them.
\end{quote}
\end{exercice}

\begin{exercice}[Transformation de temps : Past Simple $\to$ Present Perfect / Past Perfect]
Transformez les phrases suivantes au \textbf{Present Perfect} ou au \textbf{Past Perfect} selon le marqueur temporel indiqué entre parenthèses.

\begin{enumerate}
\item The youths \emph{cleaned} the stadium. (\emph{just})
\item The traditional ruler \emph{arrived} before the dance \emph{started}. (\emph{already} -- focus on ruler's arrival)
\item We \emph{prepared} the masks for the ceremony. (\emph{since Monday})
\item The dispute \emph{ended} peacefully. (\emph{by the time the police arrived})
\end{enumerate}
\end{exercice}

\begin{exercice}[Réécriture : Exprimer le but, le résultat et la préférence]
Réécrivez les phrases en utilisant les structures indiquées entre parenthèses.

\begin{enumerate}
\item The community organised a wrestling tournament. They wanted to promote unity. (\emph{to + infinitif / in order to})
\item The rain was heavy. The festival was cancelled. (\emph{so / such ... that})
\item Young people prefer dancing modern music. They don't like traditional dances as much. (\emph{prefer ... to ... / would rather ... than ...})
\item You must respect the elders. It is a rule. (\emph{have to / must})
\end{enumerate}
\end{exercice}

\begin{exercice}[Compréhension orale simulée : Travail communautaire]
Lisez le court texte ci-dessous (simulant un script audio) et répondez aux questions en phrases complètes.

\begin{quote}
\textbf{Script:} ``Last Saturday, the quarter head called for \cle{communal labour} (\emph{Appel au travail communautaire}). All residents \textbf{had to} clean the drainage canals. At first, two neighbours argued about the boundary of their plots. The quarter head \textbf{mediated} the dispute calmly. He listened to both sides and proposed a fair solution. Finally, everyone worked together peacefully.''
\end{quote}

\begin{enumerate}
\item What was the communal work about?
\item What problem occurred during the work?
\item How did the quarter head settle the dispute? (Use \emph{mediated} in your answer).
\item Which modal of obligation is used in the text? Quote the sentence.
\end{enumerate}
\end{exercice}

\courssub{Je me prépare au Bac}

\begin{exercice}[Sujet type Probatoire : Dialogue formel -- Fête de fin d'année]
\textbf{Contexte :} Vous êtes le délégué de votre classe de Première Technique. Vous devez obtenir l'autorisation du Chef de Quartier (\emph{Quarter Head}) pour organiser la fête de fin d'année de l'école dans la salle des fêtes du quartier.

\textbf{Consignes :}
\begin{enumerate}
\item Rédigez le dialogue complet (10--12 répliques au total).
\item Utilisez \textbf{au moins 4 modaux différents} parmi : \emph{may, could, must, have to, ought to, don't have to, needn't}.
\item Intégrez \textbf{5 mots/expressions du vocabulaire} de la leçon (ex: \emph{stakeholder, communal work, celebration, permission, recreational activity, mediate}).
\item Respectez les conventions de politesse formelle (salutations, formules de respect).
\end{enumerate}
\end{exercice}

\begin{exercice}[Sujet type Baccalauréat Technique : Lettre formelle -- Parrainage tournoi de lutte]
\textbf{Sujet :} \emph{``En tant que président du Comité d'Organisation du Tournoi de Lutte Inter-quartiers, écrivez une lettre formelle d'environ 150 mots au Maire de votre commune pour solliciter son parrainage.''}

\textbf{Contraintes obligatoires :}
\begin{enumerate}
\item Structure formelle de la lettre (en-têtes, objet, formule d'appel, corps, formule de politesse, signature).
\item Utilisation du \textbf{Present Perfect} pour présenter les antécédents du tournoi (ex: \emph{``We have organised three successful editions...''}).
\item Utilisation d'\textbf{au moins 3 modaux} (obligation, déduction, permission).
\item Utilisation de \textbf{connecteurs logiques} (furthermore, moreover, therefore, however, consequently).
\item Vocabulaire spécifique : \emph{sponsorship, wrestling, community cohesion, youth empowerment, budget}.
\end{enumerate}
\end{exercice}

\begin{exercice}[Sujet type Baccalauréat Technique : Composition guidée -- Jeux traditionnels et paix]
\textbf{Sujet :} \emph{``Traditional games and sports are powerful tools for peacebuilding in Cameroonian communities.'' Do you agree?}

\textbf{Plan imposé :}
\begin{enumerate}
\item \textbf{Introduction} (50 mots) : Accroche, définition des jeux traditionnels, thèse (accord/nuancé), annonce du plan.
\item \textbf{Développement} (150 mots) : Deux arguments principaux illustrés par des exemples camerounais (ex: lutte traditionnelle, \emph{Mbot}, course de pirogues, danse).
\begin{itemize}
    \item Argument 1 : Cohésion sociale / dialogue intergénérationnel (utiliser \emph{Present Perfect} pour expériences passées récentes).
    \item Argument 2 : Règlement pacifique des conflits / canalisation de l'énergie des jeunes (utiliser modaux \emph{should/ought to/must}).
\end{itemize}
\item \textbf{Contre-argument / Limites} (50 mots) : Risques (violence, tribalisme) et solutions (encadrement, règles).
\item \textbf{Conclusion} (30 mots) : Synthèse, ouverture sur le rôle de l'école/autorités.
\end{enumerate}

\textbf{Grille d'évaluation indicative (20 pts) :}
\begin{itemize}
    \item Respect du plan et nombre de mots (250 $\pm$ 10\%) : 4 pts
    \item Grammaire (Modaux, Present/Past Perfect, Let's) : 5 pts
    \item Vocabulaire thématique (fêtes, conflits, paix, communauté) : 4 pts
    \item Cohérence / Connecteurs / Ponctuation : 4 pts
    \item Pertinence des exemples camerounais : 3 pts
\end{itemize}
\end{exercice}

\newpage

\coursec{Corrigés des exercices}

\begin{corrige}[Exercice 1 — QCM : Modaux et interactions sociales]

\begin{enumerate}
\item \textbf{C) May} — The teacher's answer \emph{``Yes, you may''} confirms the question uses \cle{may} to ask for formal \emph{permission} (\emph{demander la permission de façon polie}). \emph{Could} is possible but less formal and does not match the answer.

\item \textbf{A) can} — \emph{Can} expresses here a general \emph{possibility} (\emph{il est possible que le bruit soit fort}). \emph{Must} would express a deduction (too strong without evidence), \emph{ought not to} and \emph{needn't} are negative and do not fit the meaning.

\item \textbf{B) must} — The answer \emph{``it is our duty''} shows a strong \emph{obligation} (\emph{devoir moral}), expressed by \cle{must}. \emph{Don't have to} (absence of obligation) contradicts the context.

\item \textbf{C) have to} — The principal states an \emph{external obligation / rule} (\emph{obligation extérieure, règlement scolaire}), expressed by \cle{have to}. \emph{Can/may/could} express permission or possibility, not a rule.
\end{enumerate}

\textbf{Point de vigilance :} ne pas confondre \emph{must} (obligation ressentie par le locuteur) et \emph{have to} (obligation imposée par une règle extérieure) ; ne pas confondre \emph{don't have to} (pas obligé) et \emph{mustn't} (interdiction).
\end{corrige}

\begin{corrige}[Exercice 2 — Vrai ou Faux justifié : Vocabulaire des fêtes communautaires]

\begin{enumerate}
\item \textbf{V (True).} A \cle{stakeholder} (\emph{partie prenante}) is any person or group with an interest in an event or organisation. \emph{Justification: ``The quarter head, the school and the parents are all stakeholders of the festival.''}

\item \textbf{F (False).} \cle{Communal work} is \emph{unpaid} collective work done by community members themselves, not by a private company. \emph{Justification: ``Communal work is free labour done by the inhabitants, for example cleaning the village square.''}

\item \textbf{F (False).} To \cle{mediate} (\emph{médier}) means to help both sides find an agreement \emph{without taking sides}. \emph{Justification: ``To mediate a dispute means to listen to both parties and propose a fair solution, not to punish anyone.''}

\item \textbf{V (True).} \cle{Recreational activities} (\emph{activités récréatives}) are leisure activities practised for enjoyment. \emph{Justification: ``Traditional wrestling, dancing and storytelling are recreational activities that bring people together.''}
\end{enumerate}
\end{corrige}

\begin{corrige}[Exercice 3 — Définitions et correspondance : Grammaire -- Temps parfaits]

\begin{center}
\begin{tabularx}{\linewidth}{|l|X|X|}\hline
\textbf{Phrase} & \textbf{Structure} & \textbf{Justification (FR)} \\ \hline
1. ``The elders \textbf{have settled} the dispute.'' & \textbf{B) Present Perfect} & Action passée avec un résultat visible \emph{maintenant} (le conflit est réglé, on est en paix). Aucun repère temporel précis n'est donné. \\ \hline
2. ``The elders \textbf{had settled} the dispute before the festival started.'' & \textbf{A) Past Perfect} & Action accomplie \emph{avant} une autre action passée (\emph{the festival started}) : c'est l'antériorité (\emph{antériorité dans le passé}). \\ \hline
\end{tabularx}
\end{center}

\textbf{Différence de sens :} le \cle{Present Perfect} (\emph{have/has + participe passé}) relie le passé au présent ; le \cle{Past Perfect} (\emph{had + participe passé}) situe une action \emph{avant} un autre moment du passé.

\textbf{Point de vigilance :} on n'utilise jamais le Present Perfect avec un repère de temps passé terminé (\emph{yesterday, last year, in 2020}) : dans ce cas, on emploie le Past Simple.
\end{corrige}

\begin{corrige}[Exercice 4 — Application directe : Let's / Let's not]

\begin{enumerate}
\item \textbf{Let's organise} a cleaning campaign for the school yard tomorrow. (Suggestion positive : \emph{Let's + infinitif sans to}.)
\item \textbf{Let's not forget} to invite the quarter head to the opening ceremony. (Suggestion négative : \emph{Let's not + infinitif sans to}.)
\item \textbf{Let's prepare} traditional dishes for the visitors.
\item \textbf{Let's not make} too much noise after 10 p.m. during the festival week.
\end{enumerate}

\textbf{Point de vigilance :} ne jamais écrire \emph{Let's to organise} ni \emph{Let's don't forget} : après \cle{Let's / Let's not}, on utilise toujours l'\emph{infinitif sans to} (bare infinitive).
\end{corrige}

\begin{corrige}[Exercice 5 — Dialogue complété : Demande d'autorisation]

\begin{quote}
\textbf{E:} Good morning, Chief. \textbf{May} I speak with you for a moment? \\
\textbf{C:} Of course, young man. What can I do for you? \\
\textbf{E:} We \textbf{have to} organise the end-of-year school festival next Saturday. We \textbf{must} ask for your permission to use the community hall. \\
\textbf{C:} That is a good initiative. You \textbf{ought to} submit a written request to my office first. \\
\textbf{E:} \textbf{Could} we bring it tomorrow morning? \\
\textbf{C:} Yes, that is fine. But you \textbf{don't have to} worry about the chairs; the council will provide them.
\end{quote}

\textbf{Justifications :}
\begin{itemize}
\item \emph{May I speak...?} : demande de permission très polie (\emph{registre formel}).
\item \emph{have to organise} : obligation extérieure (le calendrier scolaire l'impose).
\item \emph{must ask} : obligation ressentie comme nécessaire par le locuteur.
\item \emph{ought to submit} : conseil / recommandation (\emph{il serait bien de}).
\item \emph{Could we bring...?} : demande de permission polie.
\item \emph{don't have to worry} : absence d'obligation (\emph{ce n'est pas nécessaire}).
\end{itemize}
\end{corrige}

\begin{corrige}[Exercice 6 — Transformation de temps : Past Simple $\to$ Present Perfect / Past Perfect]

\begin{enumerate}
\item The youths \textbf{have just cleaned} the stadium. (\emph{just} $\Rightarrow$ Present Perfect : action récente avec résultat présent.)

\item The traditional ruler \textbf{had already arrived} before the dance \textbf{started}. (\emph{already + before} $\Rightarrow$ Past Perfect pour l'action antérieure ; \emph{started} reste au Past Simple.)

\item We \textbf{have prepared} the masks for the ceremony \textbf{since Monday}. \\
\emph{Variante plus naturelle :} We \textbf{have been preparing} the masks since Monday. (\emph{since} $\Rightarrow$ Present Perfect : action commencée dans le passé, lien avec le présent.)

\item The dispute \textbf{had ended} peacefully \textbf{by the time the police arrived}. (\emph{by the time} $\Rightarrow$ Past Perfect pour l'action antérieure ; \emph{arrived} au Past Simple.)
\end{enumerate}

\textbf{Point de vigilance :} dans les phrases 2 et 4, seule l'action \emph{antérieure} passe au Past Perfect ; l'autre action reste au Past Simple. Les marqueurs typiques du Present Perfect sont \emph{just, already, yet, ever, never, since, for} ; ceux du Past Perfect sont \emph{before, by the time, after} combinés à une autre action passée.
\end{corrige}

\begin{corrige}[Exercice 7 — Réécriture : Exprimer le but, le résultat et la préférence]

\begin{enumerate}
\item \textbf{But (purpose) :} The community organised a wrestling tournament \textbf{to promote unity}. \\
\emph{ou :} The community organised a wrestling tournament \textbf{in order to promote} unity.

\item \textbf{Résultat :} The rain was \textbf{so} heavy \textbf{that} the festival was cancelled. \\
\emph{ou :} It was \textbf{such} heavy rain \textbf{that} the festival was cancelled. \\
(\emph{so + adjectif + that} ; \emph{such + (adjectif +) nom + that}.)

\item \textbf{Préférence :} Young people \textbf{prefer} dancing to modern music \textbf{to} (dancing) traditional dances. \\
\emph{ou :} Young people \textbf{would rather dance to} modern music \textbf{than} (dance) traditional dances. \\
(\emph{prefer + -ing + to + -ing} ; \emph{would rather + infinitif sans to + than + infinitif sans to}.)

\item \textbf{Obligation :} You \textbf{have to} respect the elders; it is a rule. \\
\emph{ou :} You \textbf{must} respect the elders. (\emph{have to} insiste sur la règle extérieure ; \emph{must} sur le devoir.)
\end{enumerate}

\textbf{Point de vigilance :} dans \emph{prefer X to Y}, le \emph{to} est une préposition (suivie d'un nom ou d'un \emph{-ing}), pas la particule de l'infinitif ; dans \emph{would rather X than Y}, les deux verbes sont à l'infinitif sans \emph{to}.
\end{corrige}

\begin{corrige}[Exercice 8 — Compréhension orale simulée : Travail communautaire]

\begin{enumerate}
\item The communal work was about \textbf{cleaning the drainage canals} of the quarter. (\emph{Le travail communautaire consistait à nettoyer les caniveaux.})

\item \textbf{Two neighbours argued about the boundary of their plots.} (\emph{Deux voisins se sont disputés à propos de la limite de leurs parcelles.})

\item The quarter head \textbf{mediated} the dispute: he \textbf{listened to both sides and proposed a fair solution}. (\emph{Il a joué le rôle de médiateur sans prendre parti.})

\item The modal of obligation is \textbf{had to} (past of \emph{have to}): \emph{``All residents \textbf{had to} clean the drainage canals.''} It expresses an \emph{external obligation} imposed by the quarter head's call.
\end{enumerate}
\end{corrige}

\begin{corrige}[Exercice 9 — Sujet type Probatoire : Dialogue formel -- Fête de fin d'année]

\textbf{Modèle de dialogue (corrigé-type) :}

\begin{quote}
\textbf{Delegate (D):} Good morning, Sir. \textbf{May} I come in? \\
\textbf{Quarter Head (Q):} Good morning, young man. Yes, you may. What brings you here? \\
\textbf{D:} I am the class delegate of Première Technique. Our school would like to organise the end-of-year \textbf{celebration}, and we \textbf{must} find a large hall. \textbf{Could} we use the community hall on Saturday 20th June? \\
\textbf{Q:} That is an interesting project. Who are the \textbf{stakeholders} of this event? \\
\textbf{D:} The students, the teachers, the parents' association and the school authorities, Sir. \\
\textbf{Q:} I see. You \textbf{ought to} send me an official letter signed by your principal. \\
\textbf{D:} Of course, Sir. We \textbf{have to} present the programme of the \textbf{recreational activities} as well, don't we? \\
\textbf{Q:} Exactly. But you \textbf{don't have to} pay for the electricity; the council will cover it. \\
\textbf{D:} Thank you very much, Sir. We promise to organise \textbf{communal work} to clean the hall after the festival. \\
\textbf{Q:} Very good. If any dispute arises during the preparations, come and see me and I will help you \textbf{mediate} it. \\
\textbf{D:} We are grateful for your \textbf{permission} and your advice, Sir. \\
\textbf{Q:} You are welcome. Good luck with your festival!
\end{quote}

\textbf{Vérification des consignes :}
\begin{itemize}
\item 12 répliques (10--12 exigées) : respecté.
\item Modaux utilisés : \emph{may, must, could, ought to, have to, don't have to} (6 $\geq$ 4).
\item Vocabulaire : \emph{celebration, stakeholders, recreational activities, communal work, mediate, permission} (6 $\geq$ 5).
\item Politesse formelle : \emph{Good morning Sir, May I come in?, We are grateful...}.
\end{itemize}

\textbf{Barème indicatif (10 pts) :} respect du nombre de répliques et du registre formel (2 pts) ; variété et correction des modaux (3 pts) ; vocabulaire thématique (2 pts) ; cohérence du dialogue (2 pts) ; orthographe et ponctuation (1 pt).

\textbf{Points de vigilance :} saluer avant toute demande ; utiliser \emph{May/Could I...?} pour la permission (pas \emph{I want...}) ; ne pas confondre \emph{don't have to} (pas obligé) et \emph{mustn't} (interdit).
\end{corrige}

\begin{corrige}[Exercice 10 — Sujet type Baccalauréat Technique : Lettre formelle -- Parrainage tournoi de lutte]

\textbf{Modèle de lettre (corrigé-type, environ 150 mots) :}

\begin{quote}
The Organising Committee President \\
Inter-Quarter Wrestling Tournament \\
P.O. Box 123, Douala \\
15th May 2026 \\[4pt]
To His Worship the Mayor of Douala III \\[4pt]
\textbf{Subject: Request for Sponsorship of the Inter-Quarter Wrestling Tournament} \\[4pt]
Dear Sir, \\[4pt]
I am writing on behalf of the Organising Committee to request your \textbf{sponsorship} for our traditional \textbf{wrestling} tournament, scheduled for August 2026. \\[4pt]
We \textbf{have organised} three successful editions since 2023, and each one \textbf{has strengthened} \textbf{community cohesion} and \textbf{youth empowerment} in our municipality. \textbf{Moreover}, the tournament offers young people a peaceful way to spend their energy. \textbf{However}, our \textbf{budget} is limited this year; we \textbf{must} find 500\,000 FCFA for equipment and security. Your support \textbf{could} make a decisive difference. \textbf{Therefore}, we kindly ask for your patronage. Participants \textbf{ought to} respect strict rules, so the event \textbf{must} be a model of fair play. \textbf{Consequently}, your image as a promoter of peace will be reinforced. \\[4pt]
We remain at your disposal for any further information. \\[4pt]
Yours faithfully, \\[4pt]
\emph{Signature} \\
Ngo Bell Jean, President
\end{quote}

\textbf{Vérification des contraintes :}
\begin{itemize}
\item Structure formelle complète : en-têtes, date, objet, \emph{Dear Sir}, corps, \emph{Yours faithfully}, signature.
\item Present Perfect : \emph{have organised}, \emph{has strengthened}.
\item Modaux : \emph{must} (obligation), \emph{could} (possibilité), \emph{ought to} (recommandation), \emph{must} (déduction dans \emph{``must be a model''}).
\item Connecteurs : \emph{Moreover, However, Therefore, Consequently}.
\item Vocabulaire : \emph{sponsorship, wrestling, community cohesion, youth empowerment, budget}.
\end{itemize}

\textbf{Barème indicatif (20 pts) :} structure formelle de la lettre (4 pts) ; grammaire : modaux et Present Perfect (5 pts) ; connecteurs logiques (3 pts) ; vocabulaire spécifique (4 pts) ; pertinence du contenu et longueur (3 pts) ; orthographe et présentation (1 pt).

\textbf{Points de vigilance :} \emph{Yours faithfully} (et non \emph{sincerely}) quand le destinataire n'est pas nommé par son nom ; le Present Perfect exige \emph{have/has + participe passé} (\emph{organised}, pas \emph{organize} dans un contexte britannique cohérent) ; citer au moins un chiffre concret (budget, éditions) renforce la crédibilité.
\end{corrige}

\begin{corrige}[Exercice 11 — Sujet type Baccalauréat Technique : Composition guidée -- Jeux traditionnels et paix]

\textbf{Modèle de composition (corrigé-type, environ 250 mots) :}

\textbf{Introduction (environ 50 mots).}
In many Cameroonian villages, the sound of drums announces wrestling matches and canoe races. Traditional games are sporting and cultural activities inherited from our ancestors, such as wrestling, the \emph{Mbot} game and traditional dances. I strongly agree that they are powerful tools for peacebuilding, because they create social cohesion and help to settle disputes peacefully.

\textbf{Développement (environ 150 mots).}
Firstly, traditional games strengthen social cohesion and intergenerational dialogue. During the Ngondo festival, for example, elders and youths \textbf{have shared} the same canoe races for decades, and these events \textbf{have brought} rival quarters closer together. \textbf{Moreover}, traditional wrestling in the North West and Far North regions gathers whole villages around the same arena: people forget their differences and celebrate together.

Secondly, games channel the energy of young people and prevent conflicts. A young man who trains for a wrestling tournament \textbf{should} learn discipline and respect for rules; he \textbf{ought not to} fight in the streets. Chiefs and elders \textbf{must} use these gatherings to \textbf{mediate} disputes, as they have always done. \textbf{Therefore}, the playground becomes a school of peace.

\textbf{Contre-argument (environ 50 mots).}
\textbf{However}, traditional games can sometimes provoke violence or tribal rivalry when they are badly organised. To avoid this, authorities \textbf{should} supervise the competitions, apply clear rules and punish unfair behaviour.

\textbf{Conclusion (environ 30 mots).}
In conclusion, traditional games truly build peace when they are well supervised. Schools and local authorities \textbf{should} promote them, and \textbf{let's} all take part in them!

\textbf{Vérification du plan et des contraintes :}
\begin{itemize}
\item Plan en 4 parties respecté ; environ 250 mots au total.
\item Present Perfect : \emph{have shared, have brought} ; modaux : \emph{should, ought not to, must} ; \emph{Let's} en conclusion.
\item Exemples camerounais : Ngondo (courses de pirogues du peuple Sawa, à Douala), lutte traditionnelle du Nord-Ouest et de l'Extrême-Nord, \emph{Mbot}.
\item Connecteurs : \emph{Firstly, Secondly, Moreover, Therefore, However, In conclusion}.
\end{itemize}

\textbf{Application de la grille d'évaluation (20 pts) :}
\begin{itemize}
\item Respect du plan et nombre de mots (250 $\pm$ 10\%) : 4 pts
\item Grammaire (modaux, Present/Past Perfect, \emph{Let's}) : 5 pts
\item Vocabulaire thématique (fêtes, conflits, paix, communauté) : 4 pts
\item Cohérence, connecteurs, ponctuation : 4 pts
\item Pertinence des exemples camerounais : 3 pts
\end{itemize}

\textbf{Points de vigilance :} annoncer clairement sa thèse dès l'introduction ; chaque argument doit être illustré par un exemple camerounais précis ; le contre-argument ne doit pas détruire la thèse mais la nuancer ; vérifier l'accord sujet--verbe au Present Perfect (\emph{has} à la 3\up{e} personne du singulier).
\end{corrige}

\newpage

\coursec{Fiche bilan}

\begin{popfigure}
\begin{tikzpicture}[mindmap, grow cyclic, every node/.style={concept, font=\small},
root concept/.append style={concept color=popblue, font=\bfseries},
level 1/.append style={level distance=3.6cm, sibling angle=51.4},
level 1 concept/.append style={font=\footnotesize}]
\node[root concept]{Family and Social Life}
child[concept color=poporange]{node[align=center]{Community festivals\\ \footnotesize greetings, invitations, roles}}
child[concept color=popgreen]{node[align=center]{Modals\\ \footnotesize can, may, must, should, ought to, needn't}}
child[concept color=poppurple]{node[align=center]{Communal work \& disputes\\ \footnotesize listening, peace-making}}
child[concept color=poppink]{node[align=center]{Purpose, result, obligation, preference\\ \footnotesize to, so that, would rather}}
child[concept color=popgold]{node[align=center]{Recreational activities\\ \footnotesize reading, vocabulary, writing}}
child[concept color=popblue!60]{node[align=center]{Perfect tenses\\ \footnotesize present perfect, past perfect}}
child[concept color=popgreen!70!black]{node[align=center]{Let's / Let's not\\ \footnotesize suggestions, integration project}};
\end{tikzpicture}
\legende{Carte mentale de la leçon : les sept blocs de savoirs autour du thème \emph{Family and Social Life} (fêtes communautaires, modaux, travail communautaire, but/résultat/préférence, activités récréatives, temps parfaits, suggestions).}
\end{popfigure}

\begin{aretenir}[L'essentiel de la leçon]
\textbf{1. Key vocabulary (à connaître par c\oe ur)}
\begin{itemize}
\item \cle{Festivals} : celebration, feast, traditional dance, chief, invitation, to take part in, to organise, to welcome guests.
\item \cle{Communal work} : to clean up, to repair, to volunteer, community project, to settle a dispute, mediator, peace, agreement.
\item \cle{Recreation} : leisure, hobby, match, tournament, to relax, to have fun, indoor / outdoor games.
\end{itemize}

\textbf{2. Modals : formes et emplois}
\begin{center}
\begin{tabularx}{\linewidth}{|l|X|X|}
\hline
\rowcolor{popblueL} \textbf{Modal} & \textbf{Meaning} & \textbf{Example} \\
\hline
can / may & permission (may = poli) & May I join the dance group? \\
\hline
can / could & possibility & We could organise a match. \\
\hline
must / have to & obligation & You must respect the elders. \\
\hline
must & deduction (certitude) & He is tired: he must be ill. \\
\hline
should / ought to & conseil & We should clean the square. \\
\hline
needn't / don't have to & absence d'obligation & You needn't bring food. \\
\hline
\end{tabularx}
\end{center}

\textbf{3. Purpose, result, preference}
\begin{itemize}
\item \cle{Purpose} : to + infinitive ; so that + clause. \emph{We meet so that we can plan the festival.}
\item \cle{Result} : so, therefore, as a result. \emph{It rained, so the match was cancelled.}
\item \cle{Preference} : would rather + bare infinitive ; prefer + -ing. \emph{I would rather play football.}
\end{itemize}

\textbf{4. Perfect tenses}
\begin{itemize}
\item \cle{Present perfect} : have/has + past participle. Action passée liée au présent : \emph{I have just finished the work.}
\item \cle{Past perfect} : had + past participle. Action antérieure à une autre action passée : \emph{The guests had left before we arrived.}
\end{itemize}

\textbf{5. Suggestions}
\begin{itemize}
\item \cle{Let's + bare infinitive} : \emph{Let's organise a tournament!}
\item \cle{Let's not + bare infinitive} : \emph{Let's not argue about it.}
\end{itemize}

\textbf{6. Méthode express — Writing (composition)}
\begin{enumerate}
\item Understand the topic and brainstorm ideas (festivals, games, peace).
\item Plan : introduction, two or three body paragraphs, conclusion.
\item Use link words : first, then, moreover, as a result, finally.
\item Check grammar : modals, perfect tenses, spelling, punctuation.
\end{enumerate}
\end{aretenir}

\begin{attention}[Pièges fréquents à éviter]
\begin{itemize}
\item Après un modal, toujours la \cle{base verbale} (infinitif sans \emph{to}) : \emph{You must to respect} est faux ; écris \emph{You must respect}.
\item \emph{Ought to} garde son \emph{to} : \emph{We ought to help}, mais la négation est \emph{ought not to help}.
\item \emph{Mustn't} (interdiction) n'est pas \emph{needn't} (absence d'obligation) : \emph{You mustn't smoke} (interdit) ; \emph{You needn't come} (pas obligé).
\item Au present perfect, n'emploie jamais un marqueur de temps passé précis : \emph{I have seen him yesterday} est faux ; dis \emph{I saw him yesterday} ou \emph{I have just seen him}.
\item \emph{Let's} est suivi de la base verbale sans \emph{to} : \emph{Let's to play} est faux ; écris \emph{Let's play}.
\end{itemize}
\end{attention}

\begin{aretenir}[Checklist avant le Probatoire et le Bac technique]
\begin{itemize}
\item[$\square$] I know the vocabulary of festivals, communal work and recreational activities.
\item[$\square$] I can ask for permission politely with \emph{may} and \emph{can}.
\item[$\square$] I can express obligation with \emph{must} and absence of obligation with \emph{needn't} / \emph{don't have to}.
\item[$\square$] I can give advice with \emph{should} and \emph{ought to}.
\item[$\square$] I can make a deduction with \emph{must} (\emph{She must be tired}).
\item[$\square$] I can express purpose with \emph{to} and \emph{so that}.
\item[$\square$] I can express result with \emph{so} and \emph{as a result}.
\item[$\square$] I can express preference with \emph{would rather} and \emph{prefer}.
\item[$\square$] I can use the present perfect (\emph{have + past participle}) correctly.
\item[$\square$] I can use the past perfect for an action before another past action.
\item[$\square$] I can make suggestions with \emph{Let's} / \emph{Let's not}.
\item[$\square$] I can write a short composition on how recreational activities bring peace to the community.
\end{itemize}
\end{aretenir}

\findecours

\end{document}
